\newcommand{\DoNotLoadEpstopdf}{}
\documentclass[11pt]{article}
\pdfinfoomitdate=1
\pdftrailerid{}
\usepackage[T1]{fontenc}
\usepackage{mathpazo}
\renewcommand{\ttdefault}{pcr}
\usepackage[margin=1in]{geometry}
\usepackage{amsmath,amssymb,amsthm,mathtools}
\usepackage{microtype}
\usepackage{enumitem,booktabs,array}
\usepackage{xcolor,tikz,tikz-cd}
\usetikzlibrary{arrows.meta,positioning,calc}
\usepackage[hidelinks]{hyperref}
\usepackage{xurl}

\newtheorem{theorem}{Theorem}[section]
\newtheorem{lemma}[theorem]{Lemma}
\newtheorem{proposition}[theorem]{Proposition}
\newtheorem{corollary}[theorem]{Corollary}
\theoremstyle{definition}
\newtheorem{definition}[theorem]{Definition}
\theoremstyle{remark}
\newtheorem{remark}[theorem]{Remark}
\newtheorem{notation}[theorem]{Notation}

\newcommand{\C}{\mathbb C}
\newcommand{\Q}{\mathbb Q}
\newcommand{\R}{\mathbb R}
\newcommand{\Z}{\mathbb Z}
\newcommand{\N}{\mathbb N}
\newcommand{\Gm}{\mathbb G_m}
\newcommand{\OO}{\mathcal O}
\newcommand{\PP}{\mathbb P}
\newcommand{\kk}{\mathbf k}
\DeclareMathOperator{\wt}{wt}
\DeclareMathOperator{\Spec}{Spec}
\DeclareMathOperator{\Supp}{Supp}
\DeclareMathOperator{\Div}{Div}
\DeclareMathOperator{\codim}{codim}
\DeclareMathOperator{\Hom}{Hom}
\DeclareMathOperator{\gr}{gr}
\DeclareMathOperator{\DR}{DR}
\DeclareMathOperator{\id}{id}
\DeclareMathOperator{\ord}{ord}
\DeclareMathOperator{\Exc}{Exc}
\DeclareMathOperator{\rk}{rk}
\DeclareMathOperator{\im}{im}
\DeclareMathOperator{\res}{res}
\DeclareMathOperator{\Sym}{Sym}
\DeclareMathOperator{\Frac}{Frac}
\DeclareMathOperator{\vol}{vol}
\DeclareMathOperator{\Pic}{Pic}
\DeclareMathOperator{\Cl}{Cl}
\DeclareMathOperator{\Bir}{Bir}
\DeclareMathOperator{\Aut}{Aut}
\DeclareMathOperator{\tr}{tr}
\DeclareMathOperator{\lct}{lct}
\DeclareMathOperator{\mult}{mult}
\DeclareMathOperator{\Gal}{Gal}
\DeclareMathOperator{\NS}{NS}

\setlist{itemsep=3pt,topsep=5pt}
\numberwithin{equation}{section}
\setcounter{tocdepth}{1}
\allowdisplaybreaks[1]
\raggedbottom
\hypersetup{bookmarksdepth=2,pdftitle={Uniform effective log Iitaka fibrations for fourfolds},pdfauthor={OpenAI}}
\title{Uniform effective log Iitaka fibrations for fourfolds}
\author{OpenAI}
\date{September 26, 2026}

\newcommand{\BANumber}{1.2}
\newcommand{\RDenominatorNumber}{1.1}
\newcommand{\REffectiveNumber}{6.1}
\newcommand{\RCTorsionNumber}{7.1}
\begin{document}
\maketitle
\begin{abstract}
We prove the finite-rational-coefficient case of the effective log Iitaka
conjecture in dimension four. For normal projective log canonical complex
fourfolds with boundary coefficients in a fixed finite rational set and
pseudo-effective rational Cartier adjoint, one uniform degree makes the
complete rounded reflexive system nonempty and its section ratios generate
the full Iitaka field. We also prove a uniform canonical index bound for
projective klt complex fourfolds with rationally trivial canonical divisor.
\end{abstract}

\section{Introduction}\label{sec:introduction}

Let $X$ be a normal projective klt variety with $K_X\sim_\Q0$.
Its canonical index is the least positive integer $r$ for which
$rK_X$ is an integral principal divisor. The index problem asks for
a bound depending only on the dimension. More generally, for a log
Calabi--Yau pair one allows dependence on the boundary coefficients.
In dimension three, Jiang obtains a common canonical multiple for
all klt Calabi--Yau threefolds, using boundedness modulo flops to
handle the noncanonical case \cite[Theorem~1.6 and Corollary~1.7]{Jiang2021}.
Xu's inductive results give the lc threefold index theorem and the
klt fourfold theorem with nonzero boundary
\cite[Theorems~1.13--1.14]{Xu2019}. Jiang and Haidong Liu bound
surface log pluricanonical representations and, through Fujino's
gluing framework, obtain the index theorem for slc pairs through
dimension three and connected non-klt lc fourfold pairs
\cite[Theorem~1.4 and Corollaries~1.6--1.7]{JiangLiu2021}.
Birkar also proves an index bound in every fixed dimension for
rationally connected klt log Calabi--Yau pairs with coefficients in
a fixed rational set satisfying the descending chain condition
(DCC) \cite[Corollary~1.8]{Birkar2025}.
The distinction between zero and nonzero boundary, and between
rationally connected and non-uniruled varieties, is therefore central
to the remaining fourfold argument.

\begin{theorem}[Uniform canonical index]\label{thm:fourfold-index-zero}
There is a positive integer $N_4$ such that $N_4K_X\sim0$ for every
normal projective klt complex fourfold $X$ with $K_X\sim_\Q0$.
\end{theorem}

The conclusion concerns integral Weil divisors: it gives a
trivialization in one degree for all the fourfolds in the statement.
Combined with the preceding non-klt and nonzero-boundary results,
it gives the lc fourfold index theorem for any fixed finite rational
coefficient set. In applying Xu's theorem, a finite set is included
in the hyperstandard set generated by its complementary coefficients.
No assertion for arbitrary real coefficient sets is intended.

Our effective application concerns the map determined by a log
canonical divisor. For a rational Cartier divisor $D$ on a normal
integral projective variety, fix a divisor representative and put
\[
 V_m(D)=H^0\bigl(X,\OO_X(\lfloor mD\rfloor)\bigr),\qquad m\geq1.
\]
The sheaf here is the rank-one reflexive divisorial sheaf, so this
definition does not require $mD$ to be Cartier. The Iitaka field
$K(D)\subset\C(X)$ is generated by ratios of nonzero sections in
the same $V_m(D)$, over all $m$. A complete system gives the Iitaka
fibration when its ratios generate this entire field. Merely obtaining
the correct image dimension would not rule out a finite extension.

\begin{theorem}[Uniform effective log Iitaka fibrations]
\label{thm:main}
Let $\Phi\subset[0,1]\cap\Q$ be finite. There is a positive integer
$m=m(\Phi)$ such that the following holds. Suppose $X$ is a normal
integral projective complex fourfold, $\Delta\geq0$ is a rational
Weil divisor whose nonzero coefficients lie in $\Phi$, $(X,\Delta)$
is lc, and $D=K_X+\Delta$ is rational Cartier and pseudo-effective.
Then the complete system $|\lfloor mD\rfloor|$ is nonempty and its
section ratios generate $K(D)$ inside $\C(X)$. The same integer works
for every choice of canonical divisor on $X$.
\end{theorem}

In Kodaira dimension four the system is birational; in Kodaira
dimension zero it is nonempty with image a point. Nonvanishing follows
from pseudo-effectivity by the good-model theorem and effective
exceptional comparison in OpenAI, \emph{Log abundance in characteristic
zero} \cite[Theorems~1.1 and~11.1, Lemma~2.2]{OpenAIAbundance}.
This qualitative step supplies a section in some positive degree,
not a uniform degree; the uniform $m(\Phi)$ comes from the index and
fibration arguments below.

Earlier effective results separate the big case from lower-dimensional
Iitaka bases. Hacon--McKernan--Xu give a uniform birational degree
for big lc adjoints with coefficients in any fixed DCC set
\cite[Theorem~1.3]{HMX2014}. Viehweg--Zhang treat smooth varieties
of Kodaira dimension two. Their degree depends on invariants of
the general fibre; for fourfolds the fibre is a surface of Kodaira
dimension zero, and boundedness of those invariants makes the degree
universal \cite[Theorem~0.2 and the discussion preceding Corollary~0.4]{ViehwegZhang2009}.
Todorov--Xu treat klt pairs of log Kodaira codimension two in
dimensions two, three and four, with coefficients in a rational DCC set
\cite[Theorem~1.2]{TodorovXu2009}. For fourfolds their hypothesis is
$\kappa(K_X+\Delta)=2$. Chen, Han and Jihao Liu formulate the
effective log Iitaka conjecture for lc pairs of nonnegative Kodaira
dimension with coefficients in an arbitrary DCC subset of $[0,1]$,
and prove it through dimension three
\cite[Conjecture~1.2 and Corollary~1.5]{ChenHanLiu2023}.

The canonical bundle formula relates this question to effective
birationality on the base. Birkar--Zhang prove effective Iitaka
fibrations for smooth projective varieties of nonnegative Kodaira
dimension, with constants depending on the dimension, the least
nonvanishing pluricanonical degree on the general fibre, and the
middle Betti number of a smooth model of its canonical cover.
In fixed dimension and
for a fixed DCC coefficient set, their polarized effective
birationality theorem gives a uniform degree for a big base adjoint
once the nef part has bounded denominator
\cite[Theorems~1.2--1.3]{BirkarZhang2016}.
Theorem~\ref{thm:main} gives the fourfold conclusion for finite
rational coefficient sets, across all nonnegative Kodaira dimensions
and for lc singularities. Its coefficient scope is narrower than
that of the DCC conjecture and the earlier big and klt results.
The uniform degree uses the canonical-index theorem and the
relative denominator theorem below. The passage from
pseudo-effectivity to nonvanishing is the separate qualitative
input described above.

\subsection{The model comparison and the index reductions}

The uniform argument needs a model on which the adjoint defines a morphism. The
comparison in Proposition~\ref{prop:semiample-model} preserves
\emph{every} space $V_m(D)$ under a crepant dlt modification and an
adjoint-negative program. Its semiample output uses the
abundance-after-nonvanishing theorem of OpenAI,
\emph{Lifting sections from the reduced support of an adjoint}
\cite[Theorem~\BANumber]{OpenAILifting}. That theorem
applies to projective lc rational pairs over $\C$ of dimension at
most four: a nef rational Cartier adjoint of nonnegative Kodaira
dimension is semiample. In Kodaira dimension zero it is rationally
linearly trivial. We state the model comparison before the index
reduction because it is needed there as well as in the final
effective application.

The singular Beauville--Bogomolov decomposition gives a finite
quasi-\'etale cover by a product of an abelian variety and nonflat
Calabi--Yau or symplectic factors. Druel proves the needed
dimension-at-most-five result \cite[Theorem~1.2]{Druel2018};
H\"oring--Peternell prove the arbitrary-dimensional version
\cite[Theorem~1.5]{HoringPeternell2019}, using Druel's integrability
results and the holonomy decomposition of Greb--Guenancia--Kebekus
\cite[Theorem~B and Proposition~D]{GGK2019}.
Here quasi-\'etale means \'etale outside a subset of codimension
at least two. Integral cohomology bounds the abelian characters,
and lower-dimensional pluricanonical characters handle products.
The remaining single-factor case has rationally connected proper
rational images. Its noncanonical quotients are handled by a
relative torsion bound; canonical quotients admit a crepant
terminal model $V$ with the same canonical index, irregularity
zero and countable birational self-map group.

An effective divisor $D$ of Iitaka dimension one, two or three on
$V$ would give another bounded-index case. A small rational
boundary $\delta D$ produces a semiample contraction, but the
relative index theorem is applied only after proving that the
underlying canonical divisors remain crepant. The resulting
fibration has boundary zero, so the relative denominator and
rationally connected torsion theorems
\cite[Theorem~\RDenominatorNumber{} and Proposition~\RCTorsionNumber]{OpenAIRelative}
have constants independent of $D$ and $\delta$.
Consequently, in the unbounded-index case every effective divisor
of positive Iitaka dimension is big. This implies the precise
fibration property used in the scalar estimate: every intermediate
equivariant rational fibration of the canonical cover has a
geometric generic fibre of general type on a smooth projective
resolution. It is a conclusion of the reduction.

\subsection{Curves on the canonical cover and its products}

Let $r$ be the canonical index of the remaining terminal fourfold
$V$. A rational function trivializing $rK_V$ defines its canonical
cyclic cover $\pi:Y\to V$, a connected degree-$r$ quasi-\'etale cover
with deck group $\mu_r$ and trivial canonical class upstairs. The
minimality of $r$ is what makes this the primitive cover. Choose an
ample rational divisor $L$ on $V$, scaled so that its volume is
bounded above and below by absolute positive constants. Then
$(\pi^*L)^4=rL^4$.

For an ample rational divisor $P$, two local measurements will be
used. The number $\gamma(P;V)$ is the supremum of the order of
vanishing of a section in degree $k$, divided by $k$, at a very
general smooth point; degrees are sufficiently divisible to make
the divisor Cartier. Its precise definition is given in
Definition~\ref{def:orders}. For an integral curve $C$ through a
smooth marked point $x$, the ratio
\[
 \frac{P\cdot C}{\operatorname{mult}_x C}
\]
compares its degree with its multiplicity at that point. The infimum
over such curves is the Seshadri constant $\varepsilon(P;x)$.
The scalar estimates bound $\gamma$ above and $\varepsilon$ below
by absolute multiples of the fourth root of the volume, under the
fibration hypothesis just obtained. Applied upstairs, the lower
bound makes every such curve ratio grow at least as $r^{1/4}$.

Now take $t$ copies of $Y$ and divide by the simultaneous deck action:
\[
 Z_t=Y^t/\mu_{r,\mathrm{diag}},\qquad
 \theta_t:Z_t\longrightarrow V^t.
\]
The polarization $P_t$ is the pullback by $\theta_t$ of the sum of
the $t$ copies of $L$. On $Z_t$, the product estimate gives curves
through very general points whose degree divided by marked
multiplicity is bounded by a constant times $t^2$, independently of
$r$. This estimate concerns the quotient by the diagonal group;
it does not assert the same bound on $Y^t$.
The maps are displayed in Figure~\ref{fig:diagonal-cover}.

\begin{figure}[ht]
\centering
\begin{tikzcd}[column sep=large,row sep=large]
Y^t \arrow[r,"/\mu_{r,\mathrm{diag}}"] \arrow[d,"\operatorname{pr}_1"']
 & Z_t \arrow[r,"\theta_t"] \arrow[d,"q_1"]
 & V^t \arrow[d,"\operatorname{pr}_1"]\\
Y \arrow[r,"\pi"'] & V \arrow[r,equal] & V
\end{tikzcd}
\caption{The product polarization is pulled back through $\theta_t$.
The geometric generic fibre of $q_1$ is $Y^{t-1}$: choosing one lift of
the marked point in $V$ removes the simultaneous deck action. This
identification allows the growing scalar curve bound upstairs to be
tested inside a product whose curve cost is bounded independently of $r$.}
\label{fig:diagonal-cover}
\end{figure}


The product estimate is proved by reducing a fixed set of data to
positive characteristic and comparing ordinary jets with Frobenius
jets. The local comparison is that of Musta\c{t}\u{a}--Schwede
\cite[Equation (2.8)]{MustataSchwede2014}. The canonical filtration described by Sun
\cite[Theorem 3.7 and Lemma 4.2]{Sun2008}, in the ideal-theoretic
description of Kitadai--Sumihiro
\cite[Definition 3.1 and Remark 3.2]{KitadaiSumihiro2008}, supplies
the truncated symmetric powers and their determinant factors.
We prove the slope estimates and the diagonal rank comparison
needed for these products. For each variety, cover and product size,
the arithmetic model is fixed before the characteristic tends to
infinity. The later contradiction fixes a finite list of product
sizes before allowing $r$ to grow. These are separate limits;
neither requires a uniform arithmetic model for all $r$.

\subsection{From chains of curves to birational maps}

The curve estimates concern different spaces. We compare them by
forming chains from finitely many covering families of marked
curves on $Z_t$, chosen compatibly with coordinate projections and
the group actions. Each movement follows a chosen curve from its
marked point to a generic point. The closures of the generically reachable sets
are fibres of a rational quotient, called \emph{chain leaves}.
The construction uses Campana's stabilization strategy
\cite[Theorem~1.1 and its proof sketch]{Campana2004};
parameter differentiation in the manner of Ein--K\"uchle--Lazarsfeld
\cite[Proposition~2.3]{EinKuchleLazarsfeld1995} retains a degree
bound. In our setting a generic point of an $e$-dimensional leaf
can be reached in at most $e$ movements. If the marked curve
ratios are at most $B$, the leaf degree is at most $(eB)^e$.
We prove the projection and finite-quotient compatibilities needed
to use these bounds together.

Write $H_t\subset Z_t$ for a generic leaf and
$W_t=\theta_t(H_t)\subset V^t$. A positive-dimensional fibre of
$H_t\to V$ would lie, after choosing one lift of the marked base
point, in $Y^{t-1}$. The scalar lower bound there grows with $r$,
contradicting the bounded leaf degree. Thus, for sufficiently large
$r$, the maps to each coordinate are generically finite onto their
images and the leaf dimensions lie between one and four. Among a
sufficiently long finite list of product sizes, stabilization of
dimensions and a comparison of forgetting-map degrees force one
map $W_{t+1}\to W_t$ to have degree one.

This birational forgetting map produces more than a dimension
contradiction. The leaf tangent spaces identify positive-dimensional
subspaces $A_x\subset T_xV$ at different points. We show that the
resulting linear transports depend only on their two endpoints
and compose. A vector transported from one fixed point then defines
a rational vector field. Degree-one forgetting supplies rational
evaluation of its local flow, and inverse flow gives birational
self-maps. The resulting uncountable family contradicts the
countability of $\Bir(V)$, completing the index argument.

Finally, the model comparison turns the effective theorem into a
question on the base of a semiample contraction. For every
positive-dimensional base the relative big-base theorem
\cite[Proposition~\REffectiveNumber]{OpenAIRelative} gives the
whole base field in a uniform complete floor system, including when
the contraction is birational. In
Kodaira dimension zero, the three index cases give a nonzero section
in one uniform degree. The comparison of all
integer degrees brings these systems back to the original pair.
The relative results used in this argument are proved independently
of Theorem~\ref{thm:fourfold-index-zero}. The abundance-after-nonvanishing
input from \cite{OpenAILifting} and the uniform and index arguments above
do not require logarithmic or orbifold subadditivity. The full
pseudo-effective statement of Theorem~\ref{thm:main}, however, first uses
\cite{OpenAIAbundance} to obtain nonvanishing and inherits that companion's
logarithmic Iitaka subadditivity input.

\subsection{Organization and conventions}

Section~\ref{sec:reductions} constructs the semiample model and compares
all rounded systems. Section~\ref{sec:relative-inputs} records the exact
relative inputs. Section~\ref{sec:index-reductions} removes product and
rationally connected cases and proves the fibration condition needed by
the scalar estimate. Sections~\ref{sec:chains}--\ref{sec:transports}
prove the chain, scalar, product and transport statements in that order.
Section~\ref{sec:index-zero} applies them to the remaining canonical
covers. Section~\ref{sec:main-proof} derives the effective theorem.

Varieties are over $\C$ unless a geometric generic fibre or an
arithmetic reduction is specified. Canonical divisors in a
birational or finite comparison are chosen compatibly by rational
top forms. We distinguish linear equivalence $\sim$, rational
linear equivalence $\sim_{\Q}$, and numerical equivalence
$\equiv$. A pair has an effective rational boundary and a
\(\Q\)-Cartier actual adjoint. Subpairs may have negative boundary
coefficients. Log discrepancies are ordinary discrepancies plus
one; canonical and terminal singularities for zero boundary are
defined by nonnegative and positive ordinary discrepancies on
exceptional divisors. We use the usual dlt and slc conventions,
including conductor gluing in slc adjunction; see \cite{Kollar2013}.

A contraction is a surjective projective morphism $f:X\to Z$ of
normal varieties with $f_*\OO_X=\OO_Z$. A rational fibration is
understood through a relatively algebraically closed base field
inside the total function field. A log Calabi--Yau adjoint is
\(\Q\)-linearly trivial. A uniform index means a common positive
integer giving an integral linearly trivial multiple, not merely
a numerical relation. Quasi-\'{e}tale means finite surjective and
\'{e}tale outside codimension at least two. Rational connectedness
on singular spaces refers to a smooth projective model when used
in a birational argument.

Divisorial sheaves test rational sections at prime divisors. In
particular, $v\ne0$ belongs to
$H^0(X,\OO_X(\lfloor mD\rfloor))$ precisely when
$\Div(v)+mD\geq0$. Numerical positivity of a rational divisor
always concerns its rational Cartier class. Section-order and
volume limits may be taken in sufficiently divisible degrees:
powering a section preserves its normalized order. DCC and ACC
stand for the descending and ascending chain conditions.

\tableofcontents

\section{Semiample models and rounded section spaces}
\label{sec:reductions}

We pass from the original pair to a semiample model while preserving
every rounded section space. The following comparison keeps every
integer degree, even when the corresponding adjoint multiple is not
Cartier.

We use the following precise theorem of the companion
\emph{Lifting sections from the reduced support of an adjoint}
\cite[Theorem~\BANumber]{OpenAILifting}.

\begin{theorem}[Abundance after nonvanishing]\label{thm:fourfold}
Let $(X,\Delta)$ be a projective lc rational pair over $\C$ of
 dimension at most four, with effective boundary and nef
 $\Q$-Cartier adjoint $D=K_X+\Delta$. If $\kappa(X,D)\ge0$,
 then $D$ is semiample. If $\kappa(X,D)=0$, then $D\sim_{\Q}0$.
\end{theorem}

This is the companion's independent after-nonvanishing theorem. Its proof
uses the supported-boundary theorem and lower-dimensional abundance;
it does not use the separate fourfold nonvanishing companion. The boundary
coefficients may be arbitrary rational numbers in $[0,1]$. The conclusion
is qualitative: it supplies neither a uniform index nor a uniform
section degree. If a generated Cartier multiple has Iitaka dimension
zero, its morphism has image a point and a generating section is nowhere
zero, which explains the actual rational linear triviality.

For the original pair we first run an adjoint-negative program and
compare all section spaces. For the moving-divisor application in
Lemma~\ref{lem:index-moving-divisor}, the same procedure applies to a
small rational perturbation $(V,\delta D)$. There it supplies only a
contraction. Uniform statements enter after the boundary is restored
to zero and canonical crepancy is proved separately.

\begin{proposition}\label{prop:semiample-model}
Let $(X,\Delta)$ be a projective lc rational pair of dimension four
over $\C$, with $D=K_X+\Delta$ \(\Q\)-Cartier and
$\kappa(X,D)\geq0$. There is a projective \(\Q\)-factorial dlt pair
$(X',\Delta')$ with semiample adjoint $D'=K_{X'}+\Delta'$ such that,
under the identification of rational function fields and compatible
canonical divisors,
\begin{equation}\label{eq:all-rounded-sections}
 H^0\bigl(X,\OO_X(\lfloor mD\rfloor)\bigr)
 =H^0\bigl(X',\OO_{X'}(\lfloor mD'\rfloor)\bigr)
 \qquad(m\geq1).
\end{equation}
The coefficients of $\Delta'$ belong to those of $\Delta$ together
with $1$. The map is obtained by a crepant dlt modification followed
by adjoint-negative birational steps. For a \(\Q\)-factorial klt
input the latter steps can be taken directly, with klt outputs.
The semiample contraction $f:X'\to Z$ has
$\dim Z=\kappa(X,D)$ and
\[
 D'\sim_{\Q}f^*A
\]
for an ample rational Cartier divisor $A$ on the normal projective
variety $Z$.
\end{proposition}

\begin{proof}
We first establish the section comparison used throughout the
construction. For a nonzero rational function $v$ and a rational
divisor $L$ on a normal variety,
\begin{equation}\label{eq:floor-pole-test}
 v\in H^0\bigl(\OO(\lfloor mL\rfloor)\bigr)
 \quad\Longleftrightarrow\quad
 \Div(v)+mL\geq0.
\end{equation}
The equivalence holds because the coefficients of $\Div(v)$ are
integers. Effective \(\Q\)-Cartier pullback and pushforward give
both inclusions for any crepant birational morphism. More generally,
suppose that two models have a common resolution with maps $p,q$
to their input and output, and that their adjoints satisfy
\[
 p^*D_{\mathrm{in}}=q^*D_{\mathrm{out}}+E,
 \qquad E\geq0,\qquad E\text{ is }q\text{-exceptional}.
\]
A section on the output pulls back, remains effective after adding
$mE$, and pushes down to a section on the input. A section on the
input pulls back and pushes by $q$ to a section on the output, since
$q_*E=0$. Applying \eqref{eq:floor-pole-test} proves equality for
every $m$, without assuming that $mD$ is Cartier. Thus the equality
survives any sequence of these comparisons.

We now construct the models to which this comparison applies.
Take a crepant projective \(\Q\)-factorial dlt modification. This is
the usual dlt modification theorem; see, for example,
\cite[Theorem~2.8]{FujinoGongyo2014}. Its boundary consists of the
strict transform of $\Delta$ and exceptional components of coefficient
one.

Run an adjoint-negative minimal model program. We recall the existence
and termination inputs needed here. On a fixed negative extremal ray,
slightly decreasing the coefficient-one part gives a klt boundary
for which that ray is still negative. Klt contraction and flip
existence apply \cite{BCHM2010}. To identify its flip with that of the
original adjoint, write the perturbed klt adjoint as $D_\epsilon$ and
the contraction as $g$. Relative Picard number one gives a positive
rational number $a$ such that $L=D-aD_\epsilon$ is numerically trivial
over $g$. For sufficiently divisible $k$, both $kL$ and $-kL$ are
Cartier, and their differences from $D_\epsilon$ are relatively ample.
The relative base point free theorem makes both relatively semiample.
Their associated maps are constant on each connected fibre of $g$,
so a further multiple of $L$ descends to a Cartier divisor on the
normal contraction base. Thus $L$ is a rational pullback. On the
small flipped model, the transform of $D$ is consequently a positive
rational multiple of the relatively ample transform of $D_\epsilon$
plus that pullback. It is relatively ample as required. The original
dlt condition is preserved by these negative steps. A fibre-type end
is excluded by \(\Q\)-linear effectivity of the adjoint, which follows
from $\kappa\geq0$. Indeed, a curve in a general fibre not contained in the
support of an effective representative cannot have negative
intersection with that representative. Divisorial contractions
decrease the Picard number, and the remaining flips terminate by
\cite[Theorem~1.1]{ChenTsakanikas2023}, applied with zero nef
data to the pseudo-effective rational lc adjoint. The output is
therefore nef and \(\Q\)-factorial dlt. Starting with a klt pair
keeps all these outputs klt.

The initial modification is crepant, and discrepancy comparison for
each adjoint-negative step has exactly the effective exceptional form
used above. Iteration proves \eqref{eq:all-rounded-sections}; in
particular the nef output still has nonnegative Kodaira dimension.
Theorem~\ref{thm:fourfold} makes it semiample. A generated Cartier
multiple, followed by Stein factorization, gives the stated
contraction and ample rational divisor.

We also identify its field with the whole embedded Iitaka field.
For sufficiently divisible $a$, the divisor $aD'$ differs from $f^*(aA)$
by an integral principal divisor. Complete systems of sufficiently high
multiples of the ample divisor $aA$ generate $\C(Z)$, so
$\C(Z)\subset K(D')$. Conversely, take two nonzero sections in any
$V_m(D')$. Raising both to a sufficiently divisible common power puts
them in such a pullback degree. Their ratio to that power lies in
$\C(Z)$ by $f_*\OO_{X'}=\OO_Z$. The original ratio is algebraic over
$\C(Z)$, which is relatively algebraically closed in $\C(X')$ because
$f$ is a contraction. Thus every rounded-degree ratio already lies in
$\C(Z)$. Together with \eqref{eq:all-rounded-sections}, this proves
$K(D)=K(D')=\C(Z)$ inside the common function field.
\end{proof}

\begin{lemma}\label{lem:multiples-full-field}
Let $D$ be a rational Cartier divisor on a normal projective variety
such that $D\sim_{\Q}f^*A$ for a contraction $f:X\to Z$ and an
ample rational Cartier divisor $A$. If one nonempty complete system
$|\lfloor mD\rfloor|$ generates $\C(Z)$ inside $\C(X)$, then every
positive integer multiple of $m$ does so as well.
\end{lemma}

\begin{proof}
The field identification just proved uses only $D\sim_\Q f^*A$
and the contraction property, so every rounded-degree ratio lies in
$\C(Z)$. Fix a nonzero section $s$ in degree $m$. For each $k\ge1$,
multiplication by $s^{k-1}$ embeds $V_m(D)$ in $V_{km}(D)$: the
pole inequalities add, without a Cartier assumption. Ratios of the
resulting sections are unchanged. Hence the degree-$km$ ratios contain
the whole field generated in degree $m$, and cannot generate anything
beyond $\C(Z)$.
\end{proof}

\section{Relative adjunction inputs}\label{sec:relative-inputs}

The index reductions use uniform torsion on rationally connected bases;
the effective application uses full field generation over a big base.
We record the three exact results proved in the companion
\emph{Relative denominators and effective systems for log Calabi--Yau
fibrations} \cite{OpenAIRelative}. Their proofs are independent of the
canonical-index theorem proved here and of abundance after nonvanishing.

A generalized pair $(Z,B_Z+M_Z)$ includes a rational b-divisor
$\mathbf M$ that is the pullback of a nef rational Cartier divisor
on a higher projective model. The divisor $M_Z$ is its trace on $Z$.
The expression ``$p\mathbf M$ b-Cartier'' means that this nef divisor
has a Cartier multiple $pM_W$ on one determination $W$, whose
pullbacks give the data on all higher models. Generalized discrepancies
are computed after subtracting the nef data on such a model. The
following result concerns an exact presentation by rational divisors.

\noindent\textit{Source: \cite[Theorem~\RDenominatorNumber]{OpenAIRelative}.}
\begin{theorem}[Uniform relative denominators]
\label{thm:relative-denominators}
Fix a finite set $\Phi\subset[0,1]\cap\Q$ and an integer $1\leq d\leq4$.
There are positive integers $p_0\mid p$, with $p_0$ clearing $\Phi$, and a
DCC set $\mathcal B\subset[0,1]\cap\Q$, depending only on $d$ and $\Phi$,
with the following property.

Let $(X,B)$ be a projective lc pair of dimension $d$, with effective rational
boundary having coefficients in $\Phi$.  Suppose that $f:X\to Z$ is a
contraction onto a projective base $Z$, with $\dim Z>0$, and
\[
 K_X+B\sim_\Q f^*D
\]
for a $\Q$-Cartier divisor $D$ on $Z$.  There exist
$\psi\in\C(X)^*$ and a $\Q$-Cartier divisor $D_Z\sim_\Q D$ such that
\begin{equation}\label{eq:relative-exact-presentation}
 K_X+B+\frac1{p_0}\Div(\psi)=f^*D_Z,
 \qquad D_Z=K_Z+B_Z+M_Z.
\end{equation}
Here $B_Z$ is effective with coefficients in $\mathcal B$,
$(Z,B_Z+M_Z)$ is generalized lc, and $\mathbf M$ is b-nef with
$p\mathbf M$ b-Cartier.  If $(X,B)$ is klt, the generalized pair on $Z$ is
generalized klt.
\end{theorem}

The constant $p_0$ trivializes the geometric generic adjoint, while
$p$ controls the moduli divisor on a determination. These are actual
divisor statements, not numerical equivalences. A varying coefficient
introduced only to construct a contraction will not enter the fixed
coefficient set when this theorem is applied with boundary zero.

\noindent\textit{Source: \cite[Proposition~\REffectiveNumber]{OpenAIRelative}.}
\begin{proposition}\label{prop:effective-base}
Under the hypotheses of Theorem~\ref{thm:relative-denominators}, suppose in
addition that $D$ is big.  There is a positive integer $m=m(d,\Phi)$ such
that, for every positive multiple $l$ of $m$, the complete divisorial
system
\[
 \left|\left\lfloor l(K_X+B)\right\rfloor\right|
\]
is nonempty and generates precisely the subfield $\C(Z)\subset\C(X)$.
Its rational map is therefore birationally equivalent to $f$ and to the
Iitaka fibration of $K_X+B$.
\end{proposition}

Here the complete floor system generates the entire relatively
algebraically closed base field, including its algebraic part. The
conclusion holds for every positive multiple of the stated degree.

\noindent\textit{Source: \cite[Proposition~\RCTorsionNumber]{OpenAIRelative}.}
\begin{proposition}\label{prop:rc-torsion}
Fix an integer $1\leq d\leq 4$, a DCC set
$I\subset [0,1]\cap\Q$, and a positive integer $p$.
There exists a positive integer $\ell=\ell(d,I,p)$ with the following
property.  Let $Z$ be a normal projective complex variety of dimension $d$
with a rationally connected smooth projective resolution.  Suppose that
$(Z,B+M_Z)$ is generalized klt, that $B\geq 0$ has coefficients in $I$,
and that the b-nef rational b-divisor $\mathbf M$ has $p\mathbf M$
b-Cartier.  If
\[
 D=K_Z+B+M_Z\sim_{\Q}0,
\]
then $\ell D$ is an integral principal divisor.  Consequently, for every
integer $a\geq1$, the divisor $a\ell D$ is principal and
\[
 h^0\bigl(Z,\OO_Z(a\ell D)\bigr)=1.
\]
The same integer can be chosen for all dimensions in any fixed subset of
$\{1,2,3,4\}$.
\end{proposition}

The last statement will be used directly with zero boundary and zero
nef data on a klt variety with rationally connected resolution. It also
applies to the generalized base of the first statement. After taking a
common multiple with $p_0$, its principal divisor pulls back to an integral
principal canonical multiple on the total space. This is the only
uniform-torsion step in Lemma~\ref{lem:index-moving-divisor}.

\section{Structural reductions for the canonical index}\label{sec:index-reductions}

We reduce the canonical-index theorem to a particular terminal case.
First we bound the volume-form characters when the decomposition has
several factors or is purely abelian. A remaining single nonflat factor
has only rationally connected proper rational images. This bounds the
index for noncanonical quotients and for terminal models carrying an
effective divisor of intermediate Iitaka dimension. The cases left over
have irregularity zero and the fibration property needed for the scalar
bounds of Section~\ref{sec:scalar}.

All covers and all varieties in this
section are connected unless otherwise specified.  For a normal variety,
\(\Omega_X^{[j]}\) denotes the reflexive extension of \(\Omega^j_{X_{\rm reg}}\).
The order of \(K_X\sim_{\Q}0\) is the least positive integer \(r\) such that
\(rK_X\sim0\), with linear equivalence understood for integral Weil divisors.

\subsection{Finite covers and canonical characters}

We first record the relation between the canonical order and the action on a
volume form.  Suppose that \(q:R\to X\) is a finite Galois quasi-\'etale cover
of normal projective varieties, with group \(G\), and that \(K_R\sim0\).
Choose a generating reflexive top form \(\omega\) on \(R\).  Since a regular
function on \(R\) is constant, the action on this one-dimensional space has
the form
\[
                   g^*\omega=\chi(g)\omega,
             \qquad \chi:G\longrightarrow\C^*.
\]
Then
\begin{equation}\label{eq:index-character}
          mK_X\sim0\quad\Longleftrightarrow\quad\chi^m=1.
\end{equation}
Indeed, an invariant rational pluricanonical form on \(R\) descends at the
function-field level to \(X\).  Since \(q\) is \'etale at the generic points
of divisors, the descended form has neither zeros nor poles in codimension
one when the original form does not.  It therefore trivializes the
divisorial sheaf \(\OO_X(mK_X)\).  Conversely, a trivializing form downstairs
pulls back to a scalar multiple of \(\omega^{\otimes m}\), so that this
power is invariant.  In particular, bounding the order of \(\chi\) bounds
the exact linear-equivalence order, including the Cartier index.

For any klt \(X\) with \(K_X\sim_{\Q}0\), its cyclic index cover is obtained
by normalizing the cover of \(X_{\rm reg}\) defined by an \(r\)-th root of
a primitive trivialization of \(rK_X\).  Minimality of \(r\), or equivalently
the Kummer description of this cyclic extension, makes the cover connected
of degree \(r\).  It is quasi-\'etale and has a generating reflexive volume
form.  The finite discrepancy formula shows that it is klt.  Its canonical
divisor is Cartier and linearly trivial, so its ordinary discrepancies are
integers greater than \(-1\); it is consequently canonical.  We use the
finite discrepancy formula in this form throughout; see
\cite[Proposition~5.20]{KollarMori1998}.

The singular Beauville--Bogomolov decomposition theorem applied to this
cover gives a finite quasi-\'etale cover
\begin{equation}\label{eq:index-product-cover}
        P=A\times P_1\times\cdots\times P_s\longrightarrow X,
\end{equation}
where \(A\) is an abelian variety, possibly a point, and the positive
dimensional \(P_i\) are irreducible Calabi--Yau or irreducible holomorphic
symplectic varieties in the singular sense
\cite[Definition~1.4 and Theorem~1.5]{HoringPeternell2019}.
In dimension four the decomposition is already supplied by Druel
\cite[Theorem~1.2]{Druel2018}. The arbitrary-dimensional formulation
above combines the algebraic integrability and
holonomy methods discussed by H\"oring--Peternell in the introduction
to \cite{HoringPeternell2019}; in particular they identify the
contributions of Druel and of Greb--Guenancia--Kebekus
\cite[Theorem~B and Proposition~D]{GGK2019}.
We use the following precise properties.  Each \(P_i\) is normal,
projective and canonical, with \(K_{P_i}\sim0\) and dimension \(d_i\ge2\).
For every connected normal finite quasi-\'etale cover of \(P_i\), its algebra
of reflexive forms is generated by a volume form in the Calabi--Yau case,
and by a generically nondegenerate symplectic form in the symplectic case.
In particular, there are no reflexive one-forms or reflexive
\((d_i-1)\)-forms on these covers.  A one-dimensional factor is included in
\(A\).

\begin{proposition}\label{prop:index-product-case}
There is a uniform positive integer \(N_{\rm prod}\) with the following
property.  If a projective klt fourfold \(X\), with \(K_X\sim_{\Q}0\),
admits a cover as in \eqref{eq:index-product-cover} having at least two
positive dimensional factors, or having only an abelian factor, then
\(N_{\rm prod}K_X\sim0\).
\end{proposition}

\begin{proof}
\textbf{Preserving the product directions.}
First suppose that the product has at least two positive dimensional
factors.  Denote these factors by \(B_1,\ldots,B_t\), counting the abelian
factor, when present, as a single factor.  Thus \(t\le4\),
\(1\le\dim B_i\le3\), and at most one \(B_i\) is abelian.  Take the normal
Galois closure \(R\to X\) of the cover, with group \(G\).  This is still
quasi-\'etale: over the smooth locus of \(X\) the original cover is finite
\'etale by purity, and its Galois closure is finite \'etale there as well.
The map \(R\to P\) is also quasi-\'etale.  In particular, \(R\) is canonical
and \(K_R\sim0\).  Here and below purity is used in its usual form for a
finite generically \'etale map to a regular variety
\cite[Tag~0BMB]{StacksProject}.

The product directions pull back, and extend reflexively, to a splitting
\begin{equation}\label{eq:index-directions}
                     \mathcal T_R=\bigoplus_{i=1}^t\mathcal E_i.
\end{equation}
There are no homomorphisms between two distinct summands.  To check this,
choose a nonabelian factor among the corresponding two factors, say
\(B_i\), and fix general smooth points on all the other factors.  After
shrinking the complementary parameter space, the morphism
\(R\to\prod_{j\ne i}B_j\) has normal, pure-dimensional fibres.  This follows
from normality of the geometric generic fibre in characteristic zero,
generic flatness, and openness of geometric normality.  Each of its
connected components is therefore normal and finite surjective over
\(B_i\).  The codimension-two locus excluded from the quasi-\'etale map
\(R\to P\) meets a general slice in codimension at least two, so these
component covers are quasi-\'etale.

On the smooth big open of such a component, \(\mathcal E_i\) restricts to
its tangent sheaf and every complementary summand restricts to a trivial
bundle.  The component has no reflexive one-forms by the defining property
of \(B_i\).  It has no vector fields either: contraction with its generating
volume form identifies its reflexive tangent sheaf with its sheaf of
reflexive \((\dim B_i-1)\)-forms.  A homomorphism in either direction between
the two summands thus restricts to zero: the resulting forms or vector
fields on the big open extend reflexively across its codimension-two
complement.  These slices sweep a dense open
of \(R\), proving the assertion.

Consequently the projections in \eqref{eq:index-directions} are central
idempotents in the finite-dimensional algebra
\(\operatorname{End}_R(\mathcal T_R)\).  Its center is a finite-dimensional
commutative algebra.  Its primitive idempotents determine nonzero direct
summands of positive generic rank, whose ranks sum to four; there are at
most four of them.  The group \(G\) permutes these primitive idempotents.
The kernel \(G_0\) of that permutation action has index at most \(4!\) and
fixes every block projection, since each block projection is a sum of
primitive central idempotents.  It therefore preserves every
\(\mathcal E_i\).

\smallskip\noindent\textbf{Recovering the factor fields.}
Let \(F_i\) be the relative algebraic closure of \(\C(B_i)\) in \(\C(R)\),
and let \(R_i\to B_i\) be the normalization in this finite field extension.
At the generic point, the complementary directions in
\eqref{eq:index-directions} span
\(\operatorname{Der}_{\C(B_i)}\C(R)\): this follows first on the product
and then under the finite separable extension \(\C(P)\subset\C(R)\).
In characteristic zero the elements annihilated by all these derivations
are exactly those algebraic over \(\C(B_i)\).  Thus \(G_0\) preserves
\(F_i\), and acts by birational maps on \(R_i\).

The finite map \(R_i\to B_i\) is quasi-\'etale.  For completeness, the
regular extension obtained by adjoining the geometrically integral
product of the other factors is
linearly disjoint from the finite extension \(F_i/\C(B_i)\).  Hence the
function field of
\(R_i\times\prod_{j\ne i}B_j\) is an intermediate field between
\(\C(P)\) and \(\C(R)\).  This product is normal, so the normalization
property gives finite maps
\[
 R\longrightarrow R_i\times\prod_{j\ne i}B_j
       \longrightarrow P.
\]
Ramification over a prime divisor of \(B_i\) would persist after taking this
product.  Its ramification index would then divide the ramification index
at any prolongation to \(\C(R)\), contradicting quasi-\'etaleness of
\(R\to P\).  Thus \(R_i\) is projective and canonical with
\(K_{R_i}\sim0\).  Choose its generating volume \(\omega_i\) to be the
pullback of a generating volume on \(B_i\).

\smallskip\noindent\textbf{Bounding the characters in lower dimension.}
The varieties $R_i$ have dimension at most three. We now use their
finite birational actions to bound the characters on their volume lines.
Let \(J_i\) be the image of \(G_0\) in \(\Bir(R_i)\).  Resolve the graph
of its finite birational action to obtain a smooth projective
\(J_i\)-variety \(W_i\) with a birational morphism to \(R_i\).  One can
first take the closure of the graph in the product indexed by the group,
where the action permutes the factors, and then use functorial resolution
\cite[Theorem~1.1]{BierstoneMilman2008}.  Canonicality implies that the pullback
of \(\omega_i\) is a regular top form on \(W_i\).  For every positive
integer \(m\),
\begin{equation}\label{eq:index-block-sections}
          H^0(W_i,mK_{W_i})=\C\,\omega_i^{\otimes m}.
\end{equation}
Indeed, any section on \(W_i\) gives a rational pluricanonical form on
\(R_i\) with no pole at any prime divisor there.  It extends reflexively,
and \(\OO_{R_i}(mK_{R_i})\simeq\OO_{R_i}\) has only constant global
sections.  This also proves that the character in degree \(m\) is the
\(m\)-th power of the character \(\chi_i\) on the volume line.

Form the finite quotient \(q_i:W_i\to Q_i=W_i/J_i\).  Its branch boundary is
\[
             B_{Q_i}=\sum_D(1-1/e_D)D,
\]
where \(e_D\) is the ramification order along \(D\); unramified primes
have coefficient zero.  The quotient ramification formula gives
\[
                 K_{W_i}=q_i^*(K_{Q_i}+B_{Q_i}).
\]
The total divisor on the right is \(\Q\)-Cartier and the pair
\((Q_i,B_{Q_i})\) is klt.  These are the finite quotient and discrepancy
properties of a smooth variety with zero boundary
\cite[Proposition~5.20]{KollarMori1998}; \(\Q\)-Cartierness can also be
seen by descending a sufficiently divisible local canonical power, or
by taking a norm after trivializing near the finite fibre.  Since \(J_i\)
is finite, some power of \(\omega_i\) is invariant.  Galois descent of
rational forms, followed by the displayed ramification formula, gives
\[
                     \kappa(Q_i,K_{Q_i}+B_{Q_i})\ge0.
\]

The coefficients of these boundaries lie in the fixed DCC set
\(\{1-1/e:e\in\Z_{>0}\}\), and \(1\le\dim Q_i\le3\).  Effective log Iitaka
fibrations in these dimensions
\cite[Corollary~1.5]{ChenHanLiu2023} therefore give a uniform integer
\(m_*\), made common to the three possible dimensions, for which
\[
     H^0\bigl(Q_i,\OO_{Q_i}(\lfloor m_*(K_{Q_i}+B_{Q_i})\rfloor)\bigr)
                    \ne0.
\]
A section here pulls back as a rational \(m_*\)-canonical form on
\(W_i\).  The ramification formula and coefficientwise effectiveness
show that it has no divisorial poles, so it is a nonzero regular form.
It is invariant because it was pulled back from \(Q_i\).  Equation
\eqref{eq:index-block-sections} now implies \(\chi_i^{m_*}=1\).

\smallskip\noindent\textbf{Combining the factor characters.}
The wedge of the rational pullbacks of the \(\omega_i\) to \(R\) is
the pullback of the product volume on \(P\).  It therefore generates
\(K_R\), even though \(R\) need not itself split as a product.  Its
character on \(G_0\) is the product of the block characters, and is
killed by \(m_*\).  A further factor depending only on
\([G:G_0]\le4!\) kills the character on all of \(G\).  For example,
one may multiply by the least common multiple of the integers from
one to \(4!\).  Equation \eqref{eq:index-character} proves the desired
uniform canonical multiple in this case.

\smallskip\noindent\textbf{The purely abelian case.}
If \(P=A\) is an abelian fourfold, take the same normal Galois closure
\(R\to X\).  The quasi-\'etale map \(R\to A\) is \'etale by purity.
A connected finite \'etale cover of a complex abelian variety is again
abelian: analytically it is obtained from a finite-index sublattice of
the lattice in its universal cover, and it is projective here.  Each
element of \(G\) acts integrally on
\(H^4(R,\Z)\), a free group of rank \(\binom84=70\).  Its eigenvalue on
the one-dimensional space \(H^{4,0}(R)\) is the canonical character.
If that eigenvalue has order \(d\), its cyclotomic polynomial divides an
integral characteristic polynomial of degree 70, and hence
\(\varphi(d)\le70\).  There are only finitely many such \(d\).  Their
least common multiple kills every value of the character, and
\eqref{eq:index-character} again gives a uniform trivial multiple of
\(K_X\).  Taking a common multiple for the two cases proves the
proposition.
\end{proof}

\subsection{Rational images of a single nonflat factor}

\begin{lemma}\label{lem:index-single-factor-images}
Suppose that a projective klt fourfold \(X\), with \(K_X\sim_{\Q}0\),
has a finite quasi-\'etale cover \(P\to X\) in which \(P\) is one
four-dimensional irreducible Calabi--Yau or irreducible holomorphic
symplectic factor as above.  Then the following statements hold.
\begin{enumerate}[label=\textup{(\roman*)}]
\item There is no dominant rational map from \(X\) to a positive
dimensional non-uniruled smooth projective variety of dimension less
than four.
\item Every positive dimensional rational image of dimension less
than four has rationally connected smooth projective resolution.
\item Every smooth projective model of \(X\) has irregularity zero.
\end{enumerate}
\end{lemma}

\begin{proof}
Suppose that \(X\dashrightarrow T\) contradicts (i), and set
\(j=\dim T\in\{1,2,3\}\).  Non-uniruledness implies that \(K_T\) is
pseudo-effective by \cite{BDPP2013}.  The minimal-model and abundance
theorems in dimension at most three imply \(\kappa(T,K_T)\ge0\).
In dimension three, one takes a projective \(\Q\)-factorial terminal
minimal model and applies abundance for minimal threefolds
\cite{FlipsAbundance1992,Kawamata1992}; canonical sections are unchanged
on a resolution.  The lower dimensional cases are the corresponding
curve and surface statements, or follow by first taking a product with
an abelian variety of dimension \(3-j\) and applying the same
three-dimensional results.

Choose a nonzero section \(\eta\in H^0(T,mK_T)\).  Resolve the composite
rational map from \(P\) to \(T\), obtaining a smooth projective
\(u:W\to P\) and a morphism \(f:W\to T\).  Pullback by differentials
gives a nonzero tensor
\begin{equation}\label{eq:index-pulled-tensor}
         s=f^*\eta\in H^0\bigl(W,(\Omega_W^j)^{\otimes m}\bigr).
\end{equation}
It is nonzero because the map is dominant in characteristic zero.
At the function-field level it has the form
\(s=a\beta^{\otimes m}\), where \(a\in\C(P)^*\) and \(\beta\) is a
nonzero decomposable rational \(j\)-form, obtained by pulling back a
rational top form on \(T\).

We claim that the divisorial zero locus of \(s\) is \(u\)-exceptional.
Choose an ample Cartier divisor \(H_P\) on \(P\), and put
\(H=u^*H_P\).  Since \(P\) is canonical with trivial canonical class,
\(K_W\) is \(\Q\)-linearly equivalent to an effective exceptional
divisor.  Thus it is pseudo-effective and \(K_W\cdot H^3=0\).
Generic semipositivity gives nonnegative degrees for torsion-free
quotients of \(\Omega_W^1\) with respect to any ample polarization
\cite[Theorem~2.1]{CampanaPaun2015}.  It gives the same conclusion for
torsion-free quotients of
\(\mathcal V=(\Omega_W^j)^{\otimes m}\).

Here is the tensor step explicitly.  Use an ample rational perturbation
\(H+\epsilon H_0\), where \(H_0\) is ample and \(\epsilon>0\).
The Harder--Narasimhan factors of \(\Omega_W^1\) have nonnegative slopes.
On a sufficiently high general complete-intersection curve the
filtration restricts to a filtration by semistable bundles
\cite[Theorem~6.1]{MehtaRamanathan1982}; the curve can avoid all the
codimension-two loci where the sheaves or quotients fail to be locally
free.  Tensor products of semistable bundles in characteristic zero
are semistable, with the sum of the slopes
\cite[Theorem~3.18]{RamananRamanathan1984}.  Tensor powers consequently
have filtrations whose slopes are nonnegative.  Exterior powers are
quotients of tensor powers, so the same lower slope bound applies to
\(\mathcal V\) and to its torsion-free quotients.  Letting
\(\epsilon\) decrease to zero proves the required nonnegative degree
against \(H^3\).

Saturate the rank-one subsheaf generated by \(s\) in \(\mathcal V\).
It is a line bundle \(\OO_W(Z)\), with \(Z\ge0\) the divisorial zero
divisor of \(s\), and the quotient is torsion-free.  The first Chern
class of \(\mathcal V\) is a positive integral multiple of \(K_W\).
It has degree zero against \(H^3\); the preceding quotient inequality
therefore gives
\[
                              Z\cdot H^3\le0.
\]
Every nonexceptional prime divisor has strictly positive intersection
with \(H^3\), by the projection formula and ampleness on \(P\).
Since \(Z\) is effective, it follows that \(Z\) is exceptional, as
claimed.

At the generic point of a prime divisor on \(P\), write the rational
form \(\beta\) in a regular frame of \(\Omega_P^j\), and let \(b\in\Z\)
be the minimum order of its coefficients.  The absence of divisorial
zeros of \(s\) at that point says
\[
                            \ord(a)+mb=0.
\]
Hence every coefficient of \(\Div_P(a)\) is divisible by \(m\).
Normalize \(P\) in a field component of the extension obtained by
adjoining \(z\) with \(z^m=a\).  This cover is finite and quasi-\'etale:
at any divisorial discrete valuation, removing the uniformizer to its
\(m\)-divisible order reduces the equation to a root of a unit, which
is unramified in characteristic zero.  The rational form \(z\beta\)
has no divisorial poles.  It is therefore a nonzero reflexive
\(j\)-form on the cover, and remains decomposable at its generic point.

This contradicts the defining form algebra of \(P\).  For a
four-dimensional irreducible Calabi--Yau factor there are no such
forms in degrees one, two or three on any quasi-\'etale cover.  For a
four-dimensional irreducible symplectic factor the only possible
degree is two, and every such form is a scalar multiple of the
generically nondegenerate symplectic form.  Its square is nonzero,
whereas the square of a decomposable two-form is zero.  This proves (i).

For (ii), resolve the rational image and take its maximal rationally
connected fibration.  Its smooth projective base is non-uniruled
\cite[Corollary~1.4]{GHS2003}.  A positive dimensional such base would
be a rational image of \(X\) of dimension less than four, contrary to
(i).  The base must be a point, which is exactly rational
connectedness of the resolution.

Finally let \(W_X\) be a smooth projective model of \(X\).  A regular
one-form on \(W_X\) pulls back under the dominant generically finite
rational map \(P\dashrightarrow W_X\) to a reflexive one-form on
\(P\).  Indeed, the rational map to the proper target is defined at
all codimension-one points, and the pulled-back form is regular there;
reflexive extension completes it.  Pullback is injective at the
function-field level in characteristic zero.  The absence of
one-forms on \(P\) forces the original form to vanish.  Hodge
decomposition and symmetry on the smooth projective \(W_X\) then give
\(h^1(W_X,\OO_{W_X})=h^0(W_X,\Omega^1_{W_X})=0\), proving (iii).
\end{proof}

We also need the following numerical fact about exceptional divisors.
Its constants depend on the fixed birational morphism only; no uniform
birational boundedness is involved.

\begin{lemma}\label{lem:index-exceptional-psef}
Let \(u:W\to X\) be a projective birational morphism from a smooth
projective fourfold to a normal projective variety.  A pseudo-effective
\(u\)-exceptional real divisor on \(W\) is effective.
\end{lemma}

\begin{proof}
Fix a very ample Cartier divisor \(L\) on \(X\), put \(A=u^*L\), and
fix an ample Cartier divisor \(H\) on \(W\).  Because \(A\) is big,
there are an integer \(b>0\) and an effective Cartier divisor \(F\)
such that \(bA\sim H+F\).  If \(S\) is a nonexceptional prime divisor
not contained in \(\Supp F\), then \(F\cdot S\) is an effective
codimension-two cycle.  Both \(bA\) and \(H\) are nef, and factoring
their cubes gives
\[
 \begin{split}
 (bA)^3\cdot S-H^3\cdot S
   &=F\cdot S\cdot\bigl((bA)^2+bA\,H+H^2\bigr)\ge0.
 \end{split}
\]
There are only finitely many nonexceptional primes in \(\Supp F\), and
each has positive \(A^3\)-degree.  Increasing a constant to handle these
primes gives
\begin{equation}\label{eq:index-strict-degree}
              H^3\cdot S\le C A^3\cdot S
                    =C L^3\cdot u_*S
\end{equation}
for every nonexceptional prime \(S\).

Let \(E=\sum_i a_iE_i\) be exceptional and pseudo-effective.  Choose
effective real divisors \(D_n\) whose numerical classes converge to
\([E]\).  Write \(D_n=S_n+T_n\), where \(S_n\) has no exceptional
components and \(T_n\) is supported on the finitely many exceptional
prime divisors of \(u\).  Projection gives
\[
      A^3\cdot S_n=A^3\cdot D_n\longrightarrow A^3\cdot E=0.
\]
Equation \eqref{eq:index-strict-degree} implies \(H^3\cdot S_n\to0\).
Choose divisor classes \(Q_1,\ldots,Q_\rho\) forming a basis of
\(N^1(W)_{\R}\), and choose \(c_j>0\) so that both
\(c_jH+Q_j\) and \(c_jH-Q_j\) are ample.  Effectivity gives
\[
       |Q_j\cdot H^2\cdot S_n|\le c_j H^3\cdot S_n\longrightarrow0.
\]
The Hodge index theorem makes the form
\((D,Q)\mapsto D\cdot Q\cdot H^2\) nondegenerate on
\(N^1(W)_{\R}\).  Thus \([S_n]\to0\), and consequently
\([T_n]\to[E]\).

The numerical classes of the exceptional prime divisors are linearly
independent.  In fact, an exceptional numerically trivial linear
combination and its negative are both relatively nef.  The negativity
lemma forces both to be nonpositive, hence forces the combination to
vanish as a divisor \cite[Lemma~3.39]{KollarMori1998}.  Their positive
cone is therefore closed and has unique coordinates in these classes.
Since all \(T_n\) belong to that cone, so does \([E]\); independence
then gives \(a_i\ge0\) for every \(i\).
\end{proof}

\begin{corollary}\label{cor:index-noncanonical}
The canonical orders of the noncanonical varieties in
Lemma~\ref{lem:index-single-factor-images} are uniformly bounded.
\end{corollary}

\begin{proof}
Choose a smooth projective resolution \(u:W\to X\) exhibiting a
negative ordinary discrepancy.  Its canonical divisor satisfies
\[
                  K_W=u^*K_X+E\sim_{\Q}E,
\]
where \(E\) is exceptional with a negative coefficient.  By
Lemma~\ref{lem:index-exceptional-psef}, \(K_W\) is not pseudo-effective.
The uniruledness criterion of \cite{BDPP2013} implies that \(W\) is
uniruled.  Its maximal rationally connected quotient consequently has
dimension less than four, and its base is non-uniruled
\cite[Corollary~1.4]{GHS2003}.  If that base had positive dimension, the
induced rational map from \(X\) would contradict
Lemma~\ref{lem:index-single-factor-images}(i).  Thus \(W\) is rationally
connected.  Apply Proposition~\ref{prop:rc-torsion} directly to
\((X,0)\), with zero nef data.  Its hypotheses are precisely klt
singularities, \(K_X\sim_{\Q}0\), and a rationally connected smooth
resolution, all of which have been checked.  Its uniform integer
bounds the canonical order.
\end{proof}

\subsection{Terminal models and intermediate fibrations}

\begin{lemma}\label{lem:index-terminalization}
A normal projective canonical variety \(X\) with \(K_X\sim_{\Q}0\)
has a projective crepant \(\Q\)-factorial terminalization \(V\to X\).
The canonical orders of \(V\) and \(X\) are equal.
\end{lemma}

\begin{proof}
On a fixed smooth resolution, the discrepancy divisor of \(X\) is
effective.  Any divisor exceptional over that resolution has positive
ordinary discrepancy over the smooth resolution, and adding the
pullback of an effective discrepancy divisor preserves positivity.
Consequently only the finitely many exceptional divisors already on
the fixed resolution can have ordinary discrepancy zero over \(X\).

The projective extraction theorem
\cite[Corollary~1.4.3]{BCHM2010} extracts this finite list and gives a
\(\Q\)-factorial variety.  One may arrange its klt application as
follows.  Fix an integer \(r\) with \(rK_X\) Cartier, and choose an
effective ample Cartier divisor \(H\) containing all the centers of
the selected valuations.  Since \(X\) is klt, choose \(c>0\) such that
\((X,cH)\) is lc.  Then, for every divisorial valuation \(E\),
\[
                      \ord_E(H)\le A_X(E)/c.
\]
Every exceptional log discrepancy greater than one is at least
\(1+1/r\).  Choosing a positive rational \(\delta\) with
\(\delta/c<1/(r+1)\) keeps all those discrepancies greater than one,
while the selected log discrepancies equal to one decrease strictly
below one.  Take \(\delta\) still smaller if necessary so that
\((X,\delta H)\) is klt, and apply the extraction theorem to this
pair and exactly the selected list.  The list can be empty, in which
case the theorem gives a small \(\Q\)-factorialization.

Restore boundary zero.  The extracted divisors have ordinary
discrepancy zero for \(X\), so the resulting morphism is crepant.
All divisors still exceptional over the extracted variety have
positive ordinary discrepancy; hence the variety is terminal.
Finally, compatible canonical divisors satisfy \(K_V=u^*K_X\).
A principal multiple downstairs pulls back to the same principal
multiple upstairs, and a principal multiple upstairs pushes forward
to the same one downstairs.  This proves equality of their exact
orders.
\end{proof}

\begin{lemma}\label{lem:index-moving-divisor}
There is a uniform positive integer \(N_{\rm mov}\) with the following
property.  Let \(V\) be a projective \(\Q\)-factorial terminal fourfold
with \(K_V\sim_{\Q}0\), and suppose that every positive dimensional
rational image of dimension less than four has rationally connected
smooth projective resolution.  If \(V\) has an effective Weil divisor
\(D\) with \(1\le\kappa(V,D)\le3\), then
\(N_{\rm mov}K_V\sim0\).
\end{lemma}

\begin{proof}
By \(\Q\)-factoriality, \(D\) is \(\Q\)-Cartier.  Choose a small
positive rational number \(\delta\) for which \((V,\delta D)\) is klt.
Its adjoint has Iitaka dimension \(\kappa(V,D)\), since
\(K_V\sim_{\Q}0\).  Apply Proposition~\ref{prop:semiample-model},
starting its birational program at \(V\), to obtain a \(\Q\)-factorial
klt model \(V'\) and its semiample contraction
\[
                       f:V'\longrightarrow Z,
          \qquad 1\le\dim Z=\kappa(V,D)\le3.
\]

We check the canonical divisor separately from the perturbed adjoint.
Choose compatible canonical representatives and a rational function
\(h\) with \(rK_V=\Div(h)\).  Each birational contraction in the program
pushes this same relation to the target, and each flip carries it to
the new model by its codimension-one identification.  All these
models are \(\Q\)-factorial.  On a common resolution their canonical
pullbacks are therefore the same divisor \(r^{-1}\Div(h)\).
In particular, \(K_{V'}\sim_{\Q}0\), the zero-boundary varieties are
crepant, and their canonical orders are equal.  They remain klt;
alternatively, this follows from crepant equality starting at the
terminal \(V\).

Now use Theorem~\ref{thm:relative-denominators} on the contraction
\(f:(V',0)\to Z\), with zero boundary and base class zero.  It gives
uniform integers \(p_0\mid p\), a rational function \(\psi\), and
generalized klt data satisfying
\begin{equation}\label{eq:index-zero-relative}
 K_{V'}+\frac1{p_0}\Div(\psi)
       =f^*D_Z,\qquad
 D_Z=K_Z+B_Z+M_Z\sim_{\Q}0,
\end{equation}
where \(B_Z\ge0\) has coefficients in a fixed DCC set and the nef
b-divisor has Cartier denominator dividing \(p\).  The base \(Z\) is a
rational image of \(V\), so its smooth projective resolution is
rationally connected.  Proposition~\ref{prop:rc-torsion} gives a
uniform \(l\), which we enlarge to a multiple of \(p_0\), such that
\(lD_Z\sim0\).  Multiplying \eqref{eq:index-zero-relative} by \(l\)
then gives \(lK_{V'}\sim0\), hence \(lK_V\sim0\).

The coefficient \(\delta\) and the divisor \(D\) were used only to
produce a semiample model.  The two uniform statements were applied
afterward with boundary exactly zero, in one of the three fixed base
dimensions.  Their integers thus depend on neither \(D\) nor
\(\delta\).  A common multiple for these dimensions is
\(N_{\rm mov}\).
\end{proof}

\begin{lemma}\label{lem:index-general-type-fibres}
Let \(V\) be a projective \(\Q\)-factorial terminal fourfold with
\(K_V\sim_{\Q}0\).  Suppose that every effective Weil divisor on
\(V\) of positive Iitaka dimension is big.  Let
\(\pi:Y\to V\) be its cyclic index cover, with deck group \(G\).
Then every \(G\)-equivariant rational fibration of \(Y\) with positive
dimensional base and fibre has a geometric generic fibre of general
type on a smooth resolution.  The same conclusion holds for every
such rational fibration of \(V\).
\end{lemma}

\begin{proof}
The variety \(Y\) is terminal.  In fact, an exceptional divisorial
valuation over \(Y\) restricts to an exceptional valuation over \(V\),
and the finite discrepancy formula multiplies its log discrepancy
by its positive ramification index.  The log discrepancy remains
strictly greater than one.  Also \(K_Y\sim0\).

Consider a \(G\)-equivariant rational fibration \(Y\dashrightarrow T\).
Its base field \(\C(T)\) is relatively algebraically closed in
\(\C(Y)\).  The invariant field \(\C(T)^G\) lies in
\(\C(Y)^G=\C(V)\), has the same positive transcendence degree as
\(\C(T)\), and is relatively algebraically closed in \(\C(V)\).
For the last assertion, an element of \(\C(V)\) algebraic over
\(\C(T)^G\) is algebraic over \(\C(T)\), hence belongs to
\(\C(T)\), and then is invariant.  A projective model of the invariant
field therefore gives a descended dominant rational map
\(V\dashrightarrow Z\), where \(0<\dim Z<4\).

Fix smooth projective resolutions of the maps and the bases.  Thus
we may use a smooth projective birational morphism \(u:U\to Y\),
a morphism \(f:U\to T'\) to a smooth model of the unquotiented base,
and a morphism \(g:U\to Z'\) to a projective model of the quotient
base.  The map \(g\) is constant on the generic fibre of \(f\).
Take a general hyperplane divisor \(H\) for a very ample embedding
of \(Z'\).  On the domain of \(V\dashrightarrow Z'\), pull it back
and take its closure to obtain an effective Weil divisor \(D\) on
\(V\).  The domain contains every codimension-one point, because
\(V\) is normal and \(Z'\) is proper.  The hyperplane pulls form a
moving system of linearly equivalent divisors: the rational-function
ratios on the domain extend as the same principal-divisor relations
on \(V\).  Some ratio is nonconstant, since \(\dim Z'>0\).  Thus
\(\kappa(V,D)\ge1\), and \(D\) is big by hypothesis.

Write \(q=\pi\circ u\).  We have
\begin{equation}\label{eq:index-exceptional-fibre}
                           q^*D=g^*H+E,
\end{equation}
where \(E\) is an effective \(u\)-exceptional rational divisor.
The pullback exists because \(V\) is \(\Q\)-factorial.  At every
nonexceptional generic divisorial point of \(U\), the two terms
\(q^*D\) and \(g^*H\) agree: the finite cover takes that point to a
codimension-one point of \(V\), where the original rational map is
defined.  Thus their difference is exceptional.  After the resolution
has been fixed, choose \(H\) general enough that it contains none of
the generic images under \(g\) of its finitely many exceptional
prime divisors.  The coefficient of \(g^*H\) at each such prime is
then zero, while \(q^*D\) is effective.  This proves the claimed
effectivity of \(E\).

Let \(F\) be the smooth geometric generic fibre of \(f\).  The divisor
\(q^*D\) is big on \(U\), and its restriction to \(F\) is big as
well.  Indeed, write a big rational class as the sum of an ample
rational class and an effective rational class.  The integral
geometric generic fibre is not contained in the support of the
effective term, and the ample term restricts amply.  The class
\(g^*H\) restricts trivially to \(F\), so
\eqref{eq:index-exceptional-fibre} implies that \(E|_F\) is big.

Terminality and \(K_Y\sim0\) give
\[
                        K_U\sim_{\Q}\sum_i a_iE_i,
                       \qquad a_i>0,
\]
where the sum runs over all \(u\)-exceptional prime divisors.
Since \(E\) is effective and supported on this finite set, there is
a positive rational number \(c\) such that
\(\sum_i a_iE_i-cE\ge0\).  Restriction to the geometric generic
fibre is effective.  Smooth-fibre adjunction identifies
\(K_U|_F\) with \(K_F\), and therefore \(K_F\) is the sum, up to
\(\Q\)-linear equivalence, of the big divisor \(cE|_F\) and an
effective divisor.  It is big, which proves that \(F\) is of
general type.

For a rational fibration of \(V\), use the identical argument with
\(Y=V\), \(\pi=\id\), and the trivial group.  Terminality of \(V\)
provides the same positive discrepancy comparison.
\end{proof}


We have now removed the product, noncanonical and intermediate-divisor
cases by uniform index bounds. In a hypothetical sequence of unbounded
indices, its remaining terminal models $V$ therefore have irregularity
zero and no effective nonbig divisor of positive Iitaka dimension.
Lemma~\ref{lem:index-general-type-fibres} gives the requisite general-type
fibres for both $V$ and its canonical cover $Y$. The next three sections
establish chain degree bounds and the scalar and product curve estimates
for these varieties;
Section~\ref{sec:transports} will turn their discrepancy into birational
self-maps, and Section~\ref{sec:index-zero} will finish the contradiction.

\section{Generic chains and orders of sections}
\label{sec:chains}

The canonical-index argument will produce curves of uniformly bounded
degree divided by their marked multiplicity on finite product covers.
We need to replace those curves by the fibres of a rational map and
retain quantitative degree bounds under projection and finite quotient.
This section carries out that step. Its order propagation uses the
parameter-differentiation method of Ein--K\"uchle--Lazarsfeld
\cite[Proposition~2.3]{EinKuchleLazarsfeld1995}, applied here along
generic chains. For the quotient we use the stabilization strategy in
Campana's chain-equivalence construction
\cite[Theorem~1.1 and its proof sketch]{Campana2004}, with the
generic-point formulation and quantitative estimates required here.
We prove the projection and finite-quotient compatibilities as well.

We work in characteristic zero and first specify the generic-point
conventions used in the arguments.  All finite collections of varieties,
line bundles, morphisms, and sections over $\C$ descend to countable
subfields.  After enlarging such a field and taking its algebraic closure,
we may take successive geometric generic points inside a sufficiently
large algebraically closed extension.  A point chosen generically at a
particular stage is generic over the field of all the data already chosen.
We always specify a smaller field when taking the closure of that point.

Conditions on jet-kernel dimensions involve only countably many pairs of
integers: a degree and an order.  For each pair, generic base change and
semicontinuity on smooth opens give the corresponding generic condition.
Thus their simultaneous validity at a geometric generic point is
consistent with the usual very-general-point convention over $\C$.
Countably many exceptional closed sets can be included in the initial
field of definition.  Conversely, finite data at a geometric generic
point, including a section or a marked curve, can be realized over $\C$
while preserving its genericity over that field.

We also use the following spreading convention.  If new finite-type data
are chosen on a geometric generic member of a family, they are defined
over a field finitely generated over its function field.  Spreading that
field gives a \emph{dominant} extension of the parameter space, and the
new geometric generic fibre realizes the chosen data.  Resolutions and
graphs can be spread in this fashion.  Section spaces and their jet kernels
commute with algebraically closed field extension.  No assertion about
constancy of all-degree section invariants on special fibres is intended.

\subsection{The rational quotient generated by marked families}

Fix an integral projective variety $V$. A \emph{marked covering family}
on $V$ consists of an integral parameter
space $S$, a dominant marked-point map $u:S\to V$, and a family of
subvarieties $C_b\subset V$ containing $u(b)$.  After restricting to a
dense open of $S$, we require its geometric generic member to be integral
of positive dimension.  We use the reduced incidence in $S\times V$;
its endpoint evaluation dominates $V$ because it contains the graph of
$u$.  All generic choices below are made on these opens.

Starting from a generic point $a$ of $V$, a \emph{generic movement}
chooses one of the marked families, a generic point $b$ of a geometric
component of $u^{-1}(a)$ over the field of the existing path, and then a
generic endpoint on $C_b$ over that field together with $b$.

\begin{lemma}[Generic chains]\label{lem:generic-chain}
Let a nonempty finite list of marked covering families be given on an
integral projective variety $V$ in characteristic zero.  There is a
dominant rational map
\[
 \xi:V\dashrightarrow M
\]
with geometrically integral generic fibre, whose fibre closures
$H_p\ni p$ have a constant positive dimension $e$ for general $p$.
Starting at a generic $p$, a path of at most $e$ generic movements reaches
a point generic in $H_p$ over the algebraic closure of the field of $p$.
Every generic movement preserves the value of $\xi$.  In particular, a
generic family member $C_b$ is contained in $H_{u(b)}$.

The assignment is equivariant under any finite automorphism group
preserving the family list.
\end{lemma}

\begin{proof}
Choose an algebraically closed field of definition $k_*$ and a point $p$ generic in $V$ over $k_*$.  For a tuple
$a$ of points write $K_a=\overline{k_*(a)}$ inside one sufficiently large
algebraically closed extension.  Fields for successive choices include
the entire preceding path.

Every point encountered in a generic path is individually generic in
its ambient space over $k_*$.  Indeed, over a generic mark of $V$, all
components of the geometric generic fibre of a dominant map $S\to V$
have dimension $\dim S-\dim V$.  A parameter generic in such a component
is therefore generic in $S$.  The endpoint is then generic in the
incidence with that parameter, and hence in $V$.

We will compare paths after enlarging $K_p$.  The closure of a finite
path tuple over the algebraically closed field $K_p$ is geometrically
integral.  A generic point of its base change realizes the same path
after an algebraically closed field extension, with the same projected
loci.  For a fixed sequence of families, a prefix of length $j$ has total
transcendence degree
\[
 \sum_{i=1}^j
 \bigl(\dim S_i-\dim V+\dim C_{b_i}\bigr)
\]
over its initial-point field: these are precisely the increments from
the parameter and endpoint choices.  Any tuple satisfying those
incidences has at most this transcendence degree.  A path generic after
an algebraically closed field extension has the displayed degree there,
and cannot have a smaller degree before the extension.  Thus all prefix
degrees agree with their upper bounds, proving the converse comparison
as well.  This observation also applies when a path is begun at a later
endpoint and earlier path data are subsequently adjoined to its ground
field.

Let $d_j$ be the largest dimension of the closure \emph{over $K_p$} of
an endpoint after $j$ generic movements; set $d_0=0$.  These loci can be
obtained from components of the finite-step path spaces and their
projections.  Appending a step makes the new endpoint locus contain the
old one.  To see this, the new locus contains, after extension to the
field of the preceding path and the new parameter, the whole integral
member of that last step.  It therefore contains the previous endpoint,
and hence its closure over $K_p$.  Consequently $d_j$ is nondecreasing,
$d_1>0$, and the first equality has the form
\[
 d_h=d_{h+1}=e,\qquad 1\le h\le e\le\dim V.
\]
Choose an endpoint locus $H_p$ of dimension $e$ at length $h$ and a
corresponding generic endpoint $y$.

Every possible generic next movement, made over the preceding path,
has its member and endpoint in $H_p$: its endpoint locus contains
$H_p$, and the equality of dimensions forces equality.  By the
field-extension observation this is also true for all generic choices
over $K_{p,y}$, without retaining the rest of the path.  The new endpoint
is again generic in $H_p$ over $K_p$.  For each family and each geometric component of its parameter fibre
over the generic point of $H_p$, choose the new parameter independently
of the previous path. The stabilized dimension calculation says that
the endpoint image of this generic incidence is contained in $H_p$.
This is an identity on the incidence over $H_p$, so it persists after
each subsequent independent generic base change. Induction therefore
keeps all later paths in $H_p$, with endpoints generic there over $K_p$.
We use fresh generic realizations in this argument; a previously chosen
endpoint need not remain generic after arbitrary old data are adjoined.

Now form endpoint loci of paths intrinsically over $K_y$, and realize
their generic choices also over $K_{p,y}$.  They are contained in $H_p$.
Because $y$ is generic in $V$ over $k_*$, their maximum dimension at
length $h$ is the same number $e$.  Every maximizing locus therefore
becomes $H_p$ after extension.  This proves uniqueness of the maximizing
locus at $y$.  Transporting the statement between generic points of $V$
proves uniqueness at a generic starting point and shows
\begin{equation}\label{eq:chain-leaf-equality}
 H_y=H_p
 \quad\text{when }y\text{ is generic in }H_p\text{ over }K_p.
\end{equation}

Uniqueness gives descent to $k_*(p)$.  More explicitly, choose a
projective embedding over $k_*$ and take the equation subspaces of the
saturated homogeneous ideal of $H_p$ through a degree generating that
ideal.  Every algebraic conjugate is again the unique maximizing locus.
Its Grassmann coordinates are therefore invariant, and belong to
$k_*(p)$ in characteristic zero.  They define a rational map $\xi$ to
the closure of their image, naming $H_p$.  We may replace the image by
its normalization.  Equality of names means equality of subvarieties.

For a next step from the point $y$ just considered, both endpoints are
generic in $H_p$ over $K_p$, so \eqref{eq:chain-leaf-equality} gives
equality of their names.  Realizing each choice generically after field
extension shows that this equality holds at a generic parameter and
endpoint of \emph{each} family.  It is therefore a rational identity on
its incidence.  It follows that every generic movement from a generic
point preserves $\xi$, including movements that are not part of the
original maximizing path.

For completeness, fix a generic name $m=\xi(p)$ and work over $K_m$.
A generic point of any component of the geometric fibre of $\xi$ is
generic in $V$ over $k_*$.  It belongs to the subvariety bearing its own
name, so each such component is contained in the integral subvariety
named $m$.  Conversely, a generic point of that subvariety, chosen also
generically over $K_p$, has name $m$ by
\eqref{eq:chain-leaf-equality}.  Thus the closures of these fibres equal
that integral subvariety.  Characteristic zero gives geometric
reducedness as well.  The graph closure supplies a rational family with
geometrically integral generic fibre; geometric integrality and constant
dimension spread to an open of the base, by generic flatness and the
constructibility of geometric fibre properties
\cite[\S\S9,12]{EGAIV3}.  This proves the fibre assertions.  Finally, a
finite automorphism preserving the family list takes a maximizing locus
to a maximizing locus.  Uniqueness proves equivariance.  The induced rational action on the
name space can, if desired, be made regular by taking the closure of its
graph in the product of the finitely many translates.  This leaves the
generic leaves unchanged.
\end{proof}

\subsection{Projection and finite-quotient compatibility}

\begin{corollary}\label{cor:chain-functoriality}
Suppose a dominant rational map between two ambient varieties sends
generic movements of an upper family list to generic movements of a
lower family list, or to stationary steps.  Then the image of a generic
upper leaf is contained in the corresponding lower leaf.

If a finite group preserves a family list on $V$ and
$\theta:V\to V'$ is its quotient, the subvarieties
$\theta(H_p)$ are the closures of the geometrically integral generic
fibres of a rational fibration on $V'$.
\end{corollary}

\begin{proof}
The lower rational name is unchanged along every projected step.
Following a path that fills the upper leaf proves the first assertion.
It is enough here that equality of the lower names holds at the generic
point of each projected incidence: that rational identity persists when
the step is realized over the additional upper path data.

For the second assertion, let $p'$ be generic in $V'$ and lift it to
$p\in V$.  All lifts are conjugate, and equivariance gives the same
image
\[
 W_{p'}=\theta(H_p)
\]
for every lift.  Because $p$ is algebraic over $p'$, their algebraic closures in the fixed
overfield agree: $K_p=K_{p'}$. The image equation subspaces descend
to $k_*(p')$ by conjugation invariance, just as in the proof of
Lemma~\ref{lem:generic-chain}.  They give a rational name for $W_{p'}$.
A point generic in $W_{p'}$ over $K_{p'}$ lifts to a point generic in
$H_p$ over $K_p$, and hence has the same name as $p'$.  Conversely every
ambient generic point belongs to the subvariety with its own name.
The preceding fibre-of-the-name argument proves geometric integrality
and identifies its fibre closure with $W_{p'}$.
\end{proof}

\subsection{From marked-member bounds to leaf degree}

The quotient and its compatibilities depend only on the marked families.
We now put a polarization on the ambient variety. The estimate below
bounds the degree of a leaf by propagating section orders along a path
that fills it. We first define the section-order invariant used in this
argument and later in the scalar estimates.

\begin{definition}\label{def:orders}
Let $V$ be an integral projective variety of positive dimension, and let
$P$ be a big, nef, semiample rational Cartier divisor on $V$.  For a smooth
projective birational resolution $\pi:\widetilde V\to V$, set
\[
 \gamma(P;V)=
 \sup_{\substack{k>0\text{ sufficiently divisible}\\
                    0\ne s\in H^0(\widetilde V,k\pi^*P)}}
            \frac{\ord_x(s)}{k},
 \qquad
 \rho(P;V)=(P^{\dim V})^{1/\dim V},
\]
where $x$ is very general in the smooth isomorphism locus, identified
with its lift to $\widetilde V$.  If a specified
smooth point $x$ is used instead, write $\gamma(P;V,x)$, with the
resolution chosen to be an isomorphism near $x$.  For a subvariety $A$,
the notation $\gamma(P;A)$ means the same invariant for the restriction
of $P$, pulled back to a resolution of $A$.
\end{definition}

These definitions do not depend on the chosen smooth resolution: proper
birational pushforward of the structure sheaf between normal varieties
and the projection formula preserve the pulled-back section spaces.
Common resolutions then compare any two choices.  They also allow any
additional fixed divisibility requirement on $k$, since taking a power of
a section preserves its normalized order.  The invariant is finite:
on a smooth model, restrict a given nonzero section to a general
complete-intersection curve through the mark, chosen not to lie in its
zero divisor.  Its order is at most the degree of the restricted line
bundle, which is a fixed constant times $k$.  The invariant is positively
homogeneous under rational scaling of $P$ and is unchanged by rational
linear equivalence at a very general mark.

We will repeatedly use
\begin{equation}\label{eq:orders-volume}
 \gamma(P;V)\ge \rho(P;V).
\end{equation}
Indeed, put $n=\dim V$ and choose $q>0$ such that $q\pi^*P$ is Cartier
and basepoint-free on $\widetilde V$. Its morphism $f$ to its projective
image is generically finite by bigness. The Hilbert polynomial of
$f_*\OO_{\widetilde V}$, twisted by the hyperplane bundle of the image,
gives
\[
 h^0(\widetilde V,qt\pi^*P)
   =\frac{P^n}{n!}(qt)^n+O(t^{n-1}).
\]
The leading coefficient is the generic rank of
$f_*\OO_{\widetilde V}$ times the degree of the image, divided by $n!$;
it equals $(q\pi^*P)^n/n!$ by the projection formula.
Vanishing to order at least $l$ at a smooth point imposes at most
$\binom{n+l-1}{n}$ conditions.  Taking $l/k$ below, and then arbitrarily
close to, $(P^n)^{1/n}$ proves \eqref{eq:orders-volume}.

\begin{lemma}[Orders and degrees of chain leaves]\label{lem:chain-order}
Let the marked families and their rational quotient be as in
Lemma~\ref{lem:generic-chain}, and let $e$ be the dimension of a generic
leaf. Suppose that $P$ is a big, nef, semiample rational Cartier
divisor on $V$.  Assume that, for every family, its geometric generic
marked member satisfies one of the following bounds, with a common
number $B\ge0$:
\begin{enumerate}[label=\textup{(\roman*)}]
 \item the mark is smooth on $C_b$ and
       $\gamma(P;C_b,u(b))\le B$;
 \item $C_b$ is a curve and
       $P\cdot C_b/\mult_{u(b)}C_b\le B$.
\end{enumerate}
Then on the geometric generic leaf,
\begin{equation}\label{eq:chain-order-degree}
 \gamma(P;H_p)\le eB,
 \qquad
 P^e\cdot H_p\le(eB)^e.
\end{equation}
\end{lemma}

\begin{proof}
Choose an algebraically closed field of definition $k_*$ for the data.
Let $m$ be a generic point of the quotient base and put
$K_m=\overline{k_*(m)}$ inside the fixed algebraically closed overfield.
The restriction of $P$ to the geometric
generic leaf over $K_m$ is big: the leaf sweeps $V$, so it meets the generically
finite locus of the morphism defined by a basepoint-free multiple of
$P$.  Take a smooth projective resolution $J$ of this leaf over $K_m$.
The generic marked families and paths lift to $J$ wherever the resolution
is an isomorphism.  A point generic in the leaf over $K_m$ is generic in
$V$ over $k_*$.  For a generic mark $p$ on this leaf, put $K_p=\overline{k_*(p)}$.
Its name $m$ belongs to $k_*(p)$. Lemma~\ref{lem:generic-chain} supplies
a path of length $h\le e$ reaching a point generic in the whole leaf over
$K_p$.  Every
intermediate mark and endpoint is generic in $J$ over $K_m$: it is
generic in $V$ and has name $m$, and consequently has transcendence
degree $e$ over $K_m$.  In particular all their evaluation maps dominate
$J$.

Suppose that, in a divisible degree $k$, a nonzero section of $kP$ has
order at least an integer $m_0>hkB$ at the variable generic mark $p$ of
$J$.  In a fixed $K_m$-basis of $H^0(J,kP)$, the kernel of the
generic jet-evaluation matrix lets us choose such a section
$\sigma(p)$ with coefficients rational in $p$.  Follow a path of length
$h$ that fills the leaf generically over $K_p$.

Assume inductively that the same section has multiplicity at least
$m_j$ at the current mark.  Let $T$ be the closure \emph{over $K_m$}
of the path tuple through the parameter $b$ chosen for the next step;
the original point $p$ varies in this closure.  Let
$I\subset T\times J$ be the incidence obtained by adjoining the new
endpoint.  Both the graph of the current mark and $I$ dominate $J$.
Over the generic point of $T$, the generic fibre of $I$ is the chosen
member, lifted to $J$.

Differentiate in parameter directions, keeping the variable on $J$
fixed.  Concretely, use derivations of $K_m(T)$ over $K_m$ on the
coefficients of $\sigma$ in the fixed section basis.  A line-bundle
frame from $J$ is regular and nonvanishing at the generic point of $J$,
so it is unaffected by these derivations.  At the graph of the moving
mark, the multiplicity assumption is membership in the $m_j$th power
of the graph ideal, after localizing over the generic parameter.
Parameter differentiation drops its powers by at most one.  This can
also be tested after field extension to make the mark rational, using
faithful flatness.

Put
\[
 r_j=m_j-\lfloor kB\rfloor.
\]
Every iterated derivative of length less than $r_j$ still has order
strictly greater than $kB$ at the moving mark.  Each such derivative
vanishes identically on the chosen member.  In case (i), a nonzero
restriction, pulled back to its smooth resolution, would violate the
assumed order bound at the smooth mark.  The resolutions used in this
comparison agree near the mark and can otherwise be compared on a
common resolution.

In case (ii), a section with order at least $l$ at the ambient mark
restricts to an element of the $l$th power of the maximal ideal of the
local curve ring $A$.  On its normalized branches its order is at least
$l\,v_i(\mathfrak m_A)$.  The identity
\[
 \mult A=\sum_i v_i(\mathfrak m_A)
\]
then shows that a nonzero restriction has degree at least
$l\mult A$.  One way to verify this identity is to use the finite
semilocal normalization $\bar A$.  Its quotient by
$\mathfrak m_A^t\bar A$ has length
$t\sum_i v_i(\mathfrak m_A)$ over the algebraically closed residue
field.  The conductor gives a fixed $c$ such that
\[
 \mathfrak m_A^{t+c}\bar A
 \subset \mathfrak m_A^t A
 \subset \mathfrak m_A^t\bar A.
\]
Comparing leading terms in the Hilbert--Samuel lengths gives the
identity.  Here $l>kB$, whereas the degree of a nonzero restricted
section is $kP\cdot C_b\le kB\mult A$, a contradiction.

Thus all parameter derivatives of length less than $r_j$ vanish at the
generic point of $I$.  We explain why this forces order at least $r_j$
along that incidence.  Let
\[
 R=\OO_{T\times J,\eta_I},\qquad F=K_m(J).
\]
Since $I$ dominates both factors, this ring localizes over their generic
points.  It is regular, using a smooth open of $T$.  The conormal map
\[
 \mathfrak m_R/\mathfrak m_R^2
 \longrightarrow \Omega_{R/F}\otimes\kappa(I)
\]
is injective.  Indeed the residue extension is separable in
characteristic zero; in the cotangent sequence the dimensions of the
ambient differentials and residue differentials differ by $\dim R$,
which is the conormal dimension.  Equivalently, this is the generic
smoothness computation for the incidence over $J$.
Parameter derivations span the dual of this conormal after reduction,
since $R$ is a localization of
$K_m(T)\otimes_{K_m}F$.  Regularity identifies the associated graded
ring with the symmetric algebra of the conormal.  A nonzero initial
form of degree $a<r_j$ is detected by a length-$a$ iterated parameter
derivative, by the symmetric-power pairing and the invertibility of
$a!$.  The derivative vanishings therefore imply that our section, in
the chosen frame, lies in $\mathfrak m_R^{r_j}$.

Pass to the generic parameter field and then extend it so that the
independently generic endpoint is rational. The map from the ambient
product to this fibre induces a local homomorphism from $R$ to the local
ring of $J$ at that endpoint. A denominator outside the incidence prime
has nonzero value at its generic point and stays a unit after this
extension. The incidence ideal maps into the endpoint maximal ideal;
its $r_j$th power therefore maps into that maximal ideal's $r_j$th power.
The same section consequently has multiplicity at least
\[
 m_{j+1}=r_j=m_j-\lfloor kB\rfloor
\]
at the next endpoint.  After $h$ steps it still has positive
multiplicity, since $m_0>hkB$.  But the endpoint is generic in $J$ over
$K_p$, while $\sigma(p)$ is a nonzero section defined over $K_p$.
This is impossible.  Hence $\gamma(P;H_p)\le hB\le eB$.
Applying \eqref{eq:orders-volume} on the leaf gives the second inequality
of \eqref{eq:chain-order-degree}.
\end{proof}


\section{Scalar jet bounds}
\label{sec:scalar}

The next result compares section orders and curve degrees with the
volume of an ample rational divisor. Its hypothesis on equivariant
fibrations is essential to both estimates. In the index argument,
Lemma~\ref{lem:index-general-type-fibres} establishes this
hypothesis on the varieties to which we apply the result.

We use the normalized order invariant of Definition~\ref{def:orders}.
The following estimates apply to rational polarizations of arbitrary
denominator.  All orders and section spaces are taken in sufficiently
divisible degrees; taking powers shows that imposing any further fixed
divisibility does not change the normalized supremum.

\begin{theorem}[Scalar bounds]\label{thm:scalar-bounds}
There are constants $C,c>0$ with the following property.  Let $V$ be a
normal projective canonical fourfold with $K_V\sim_{\Q}0$, and let a finite
group $G$ act on $V$.  Suppose that every $G$-equivariant rational fibration
with geometrically integral generic fibre and with both base and fibre of
positive dimension has geometric generic fibre of general type on a
smooth projective resolution.  It is enough to impose this condition on
the rational quotients supplied by Lemma~\ref{lem:generic-chain}.
If $P$ is an ample rational Cartier divisor satisfying
$g^*P\sim_{\Q}P$ for every $g\in G$, then
\begin{equation}\label{eq:scalar-bounds}
 \gamma(P;V)\le C(P^4)^{1/4},
 \qquad
 \varepsilon(P;x)\ge c(P^4)^{1/4}
\end{equation}
at a very general smooth point $x\in V$, where
\[
 \varepsilon(P;x)=
 \inf_{\substack{x\in D\subset V\\D\text{ an integral curve}}}
 \frac{P\cdot D}{\mult_xD}.
\]
The constants are independent of $V$, $G$, and $P$.
\end{theorem}

The upper estimate is the main geometric step. If a section can have
large normalized order on a polarization of very small volume, its
moving base locus produces smaller-dimensional covering members.
Numerical subadjunction bounds their canonical degrees. Repeatedly
passing to their Iitaka fibres will produce members of Kodaira dimension
zero with very small section orders. If their chains fill $V$, the
chain-order bound contradicts the initial large section order. Otherwise
they give a proper general-type leaf covered by Kodaira-dimension-zero
members, contradicting Lemma~\ref{lem:scalar-covering}. Once the upper
bound is proved, a thickening of a small-degree chain leaf gives the
lower curve bound.

We first establish the local numerical input.  For rational Cartier
classes $A,B$, we write $A\preceq B$ when $B-A$ is rationally linearly
equivalent to an effective rational divisor.  A restriction of such a
comparison will always be made to a member not contained in that effective
divisor.

\subsection{Numerical subadjunction and moving base components}

\begin{lemma}[Numerical subadjunction at a generic lc centre]
\label{lem:scalar-subadjunction}
Let $Z$ be a projective $\Q$-factorial klt variety, let $\Theta\ge0$ be a
rational divisor, and let $S\subset Z$ be an integral subvariety of
dimension $d>0$.  Suppose that $(Z,\Theta)$ is lc over a neighbourhood of
the generic point of $S$, and that $S$ is an lc centre there.  For a nef
rational Cartier divisor $A$ on $Z$ and a smooth projective resolution
$S^*\to S$, one has
\begin{equation}\label{eq:scalar-subadjunction}
 K_{S^*}\cdot(A|_{S^*})^{d-1}
 \le (K_Z+\Theta)|_S\cdot(A|_S)^{d-1}.
\end{equation}
Here the right-hand side means the intersection of the restricted
rational line bundle with the cycle $S$.
\end{lemma}

\begin{proof}
The arbitrary-boundary dlt blow-up theorem
\cite[Theorem~2.10]{FujinoHashizume} gives a projective birational morphism
$p:Q\to Z$, with $Q$ $\Q$-factorial, such that
\begin{equation}\label{eq:scalar-dlt-excess}
 K_Q+B_Q+E=p^*(K_Z+\Theta),
 \qquad (Q,B_Q)\text{ dlt},\qquad E\ge0,
\end{equation}
and $\Supp E$ lies above the non-lc locus of $(Z,\Theta)$.
The theorem allows coefficients of $\Theta$ greater than one: $B_Q$ is
the truncated strict transform together with the reduced exceptional
divisor, and the remaining crepant boundary is the effective excess $E$.

Choose a dlt stratum $R$ minimal by inclusion among the strata mapping
onto $S$.  Such a stratum exists.  Indeed, over the lc neighbourhood of
the generic point of $S$, the morphism in
\eqref{eq:scalar-dlt-excess} is crepant, and the centre on $Q$ of an lc
place with centre $S$ is a dlt stratum mapping onto $S$.
Dlt adjunction, in the form of the lc-centre and different theorem
\cite[4.16]{Kollar2013}, gives a normal $R$ and an effective
different.  Adding the restriction of $E$ gives an effective rational
divisor $\Delta_R$ with
\[
 K_R+\Delta_R\sim_{\Q}(K_Z+\Theta)|_R.
\]
The restriction of $E$ is defined and effective: $E$ is rational Cartier
and misses the generic fibre of $R\to S$.  Moreover $(R,\Delta_R)$ is
klt over the generic point of $S$.  Any lc centre of the dlt different
dominating $S$ would give a smaller dlt stratum of $Q$ dominating $S$,
contrary to the choice of $R$; the added excess is zero over that generic
point.

Let $S^\nu$ be the normalization of $S$ and factor the map through its
normal Stein base:
\[
 R\xrightarrow{f}S'\xrightarrow{h}S^\nu\longrightarrow S,
 \qquad f_*\OO_R=\OO_{S'},\qquad h\text{ finite}.
\]
Write $D$ for the pullback of $(K_Z+\Theta)|_S$ to $S^\nu$.
The morphism $f:(R,\Delta_R)\to S'$ satisfies the definition of a
klt-trivial fibration in \cite[Definition~3.1]{FujinoGongyoModuli}.
It is projective, generically subklt, and its adjoint is rationally
linearly equivalent to $f^*h^*D$.  The rank-one condition also holds:
on a resolution over the generic base point, the rounded discrepancy
divisor is effective and exceptional over $R$.  Allowing poles along
such a divisor does not enlarge regular functions on the normal variety
$R$, and connectedness of the fibres then gives rank one.

Ambro's canonical bundle formula, in the form recalled in
\cite[Theorem~3.5]{FujinoGongyoModuli}, therefore gives
\begin{equation}\label{eq:scalar-cbf}
 h^*D\sim_{\Q}K_{S'}+B_{S'}+M_{S'},
\end{equation}
where $M_{S'}$ is the trace of b-nef rational data.  The discriminant
$B_{S'}$ is effective.  To see this even where the total pair is not lc,
let $P$ be a prime of $S'$.  Over its generic point, $f^*P$ has a divisor
of multiplicity at least one and $\Delta_R$ has nonnegative coefficient
on that divisor.  The corresponding threshold is therefore at most
one.  It can be negative, but its discriminant coefficient, one minus
the threshold, remains nonnegative.

Let $A'$ be the pullback to $S'$ of $A|_{S^\nu}$.  Intersect
\eqref{eq:scalar-cbf} with $(A')^{d-1}$.  Both the effective discriminant
and the b-nef moduli trace contribute nonnegatively.  For the latter
assertion, compute on a model where the moduli divisor is nef and apply
the projection formula.  If $e=\deg h$, codimension-one ramification
for the finite map of normal varieties gives
\[
 K_{S'}\cdot(A')^{d-1}
 \ge e\,K_{S^\nu}\cdot(A|_{S^\nu})^{d-1}.
\]
These are intersections of Weil divisors with Cartier powers, so the
calculation can be made on the smooth codimension-one loci.  Finally,
the pushforward of $K_{S^*}$ is $K_{S^\nu}$, and its exceptional
divisors pair to zero with the pulled-back power of $A$.  Dividing by
$e$ proves \eqref{eq:scalar-subadjunction}.
\end{proof}

The moving-base-component strategy below adapts the jets construction in
\cite[Sections~7--8]{OpenAIAbundance}.  This method credit is distinct
from the qualitative nonvanishing input in the introduction; the public
jet-base-locus method and the present estimates are specified in the proof.

\begin{lemma}[Tracking a moving base component]\label{lem:scalar-tracking}
Let $Z$ be a projective $\Q$-factorial klt variety of dimension $n\ge2$,
and let $N$ be a big nef semiample rational Cartier divisor.  Suppose
there are $n+1$ positive rational numbers, with consecutive spacing
$\delta>0$, all strictly between $(N^n)^{1/n}$ and $\gamma(N;Z)$.
There is a covering family of integral subvarieties $S\subset Z$ of
some dimension $0<d<n$ such that, on a smooth projective resolution
$S^*$ of its geometric generic member,
\begin{align}
 N^d\cdot S&\le (2/\delta)^{n-d}N^n,
       \label{eq:scalar-track-degree}\\
 K_{S^*}\cdot(N|_{S^*})^{d-1}
 &\le \bigl(K_Z+(2n/\delta)N\bigr)|_S
             \cdot(N|_S)^{d-1}.
       \label{eq:scalar-track-canonical}
\end{align}
\end{lemma}

\begin{proof}
The moving-parameter argument follows the jet-base-locus method of
\cite[Proposition~2.3]{EinKuchleLazarsfeld1995}. We keep the
subadjunction and degree estimates explicit for the rational
polarizations used here.
Choose a sufficiently divisible positive integer $l$, clearing the
denominators of $N$ and of the levels.  It can be chosen arbitrarily
large so that every jet space under consideration is nonzero: first
choose a section witnessing an order above the highest level, and
then take powers.  At a variable very general smooth point $x$, write
\[
 V_\tau(x)=
 H^0\bigl(Z,\OO_Z(lN)\otimes\mathfrak m_x^{l\tau}\bigr).
\]
These spaces are the fibres of generic jet kernels, and thus have
rationally varying bases after restricting the parameter space.

At each level the base locus through $x$ is proper and has no isolated
component at $x$.  For the second assertion, if its local base ideal
were primary to $\mathfrak m_x$, then $n$ general members would have
local intersection multiplicity at least $(l\tau)^n$.  Successive
proper cuts, discarding other base components before each cut, would
give
\[
 l^nN^n\ge(l\tau)^n,
\]
contrary to $\tau>(N^n)^{1/n}$.  This degree argument is also detailed
below for the positive-dimensional component that we use.

The base loci are nested as the level increases.  Choose successively
an irreducible component through $x$, each contained in the next.
Their dimensions are in $\{1,\ldots,n-1\}$, so two consecutive
components agree.  Denote their common reduced component by $S$,
and spread it over a parameter space containing the varying mark $x$.
A finite extension of the parameter field may be used to name a
geometric component.  The resulting incidence dominates $Z$ because
it contains the graph of the mark.

Let $\tau$ and $\tau+\delta$ be the two levels.  We claim that the
base ideal of $V_{\tau+\delta}(x)$ has order at least $l\delta/2$ at
the generic point of $S$.  Differentiate a rationally varying basis
of this jet space in parameter directions, keeping the spatial
variable and a fixed basis of $H^0(Z,lN)$ constant.  A derivative of
length $r<l\delta$ has order at least $l(\tau+\delta)-r>l\tau$
at the varying mark, so it belongs to the lower jet space and
vanishes along the incidence of $S$.

Here these derivative vanishings detect the order of the original
section along the incidence.  At its generic point, the ambient
product is regular and both projections are dominant.  Since the
incidence dominates the spatial factor in characteristic zero,
parameter derivations span the dual of its conormal space.  The
associated graded ring of the regular local ring is the symmetric
algebra of that conormal space.  A nonzero initial form of degree
$r$ is detected by some iterated parameter derivative of length $r$,
by the characteristic-zero symmetric-power pairing.  Thus all the
vanishings just obtained force order at least $l\delta$, and in
particular the asserted weaker bound.  This is the same local Taylor
argument used in Lemma~\ref{lem:chain-order}.  Passing to the
geometric generic parameter fibre preserves it.

Put $a=n-d$.  The local ring of $Z$ at the generic point of $S$ is
regular of dimension $a$: a sweeping member meets the smooth open of
$Z$.  Since $S$ is a component of the higher base locus, its base
ideal there is primary to the maximal ideal $\mathfrak m$, and it is
contained in $\mathfrak m^{\lceil l\delta/2\rceil}$.  General choices
of $a$ members form a parameter ideal.  Their intersection length,
equal to its multiplicity in this regular local ring, is at least
$(l\delta/2)^a$.  Globally, take the cuts successively and discard
components already contained in the base locus before the next cut.
No discarded component contains the generic point of $S$, since $S$
is itself an irreducible component of that base locus.  The cuts are
therefore proper at the generic point being tracked.  Nefness of $N$
makes all discarded degree contributions nonnegative, and gives
\[
 (l\delta/2)^a(N^d\cdot S)\le l^aN^n.
\]
This proves \eqref{eq:scalar-track-degree}, as well as the isolated
base-point calculation used above.

Let $t$ be the local lc threshold of this base ideal over the generic
point of $S$.  A log resolution shows that $t$ is positive and rational.
Blowing up the smooth transverse centre, whose log discrepancy is
$a$, gives
\[
 t\le\frac{2a}{l\delta}.
\]
For $a=1$ the same inequality follows by testing $S$ itself.  Average
sufficiently many general members of the higher jet space and
multiply this average by $t$.  On a log resolution of the ideal the
free parts are general and their coefficients can be made arbitrarily
small.  This constructs an effective rational divisor
\[
 \Theta\sim_{\Q}bN,
 \qquad b=lt\le\frac{2a}{\delta},
\]
such that $(Z,\Theta)$ is lc near the generic point of $S$ and $S$ is
an lc centre there.  Indeed a place computing the threshold has
positive base-ideal order; the primary condition makes its centre
equal to $S$ after localization.  Lemma~\ref{lem:scalar-subadjunction}
with $A=N$ now gives \eqref{eq:scalar-track-canonical}, since $a\le n$.
\end{proof}

We will also need the following elementary consequence of a covering
family.

\begin{lemma}[Canonical dimension of a covering member]
\label{lem:scalar-covering}
Let $X$ be a projective variety birational to a smooth projective
variety with a nonzero pluricanonical section.  A smooth projective
resolution of the geometric generic member of a covering family of
integral subvarieties of $X$ has nonnegative Kodaira dimension.  If
$X$ is of general type, that resolution is of general type.
\end{lemma}

\begin{proof}
Pass to a smooth projective model of $X$ and to strict transforms of
the moving members.  Shrink the parameter space so that the incidence
has fibres of pure constant dimension.  If the incidence has dimension
larger than $\dim X$, intersect the parameter space in a projective
closure with general hyperplanes, in number equal to that excess.
Over the generic point of $X$ the incidence fibre embeds into the
parameter space, so a general such section meets its top-dimensional
part and makes the evaluation generically finite.  A component of
the cut incidence dominating $X$ dominates a component of the parameter
slice, by dimension.

Choose the cutting hyperplanes with independent generic coefficients
over the original field of definition.  A generic point of the selected
parameter slice remains generic in the original parameter space over
that field: every original proper closed subset has smaller dimension
after the same general cuts.  Thus the argument tests the original
geometric generic member, rather than a special subfamily.

Resolve and compactify the cut incidence and its evaluation to $X$.
The resulting generically finite morphism between smooth projective
varieties satisfies
\[
 K_I=q^*K_X+R,\qquad R\ge0.
\]
Restrict to a smooth geometric generic fibre of the parameter map.
It is birational to the corresponding member, and adjunction identifies
the restriction of $K_I$ with its canonical divisor.  A nonzero
pluricanonical section of $X$ pulls back to a section not identically
zero on this generic fibre, by dominance.  This proves nonnegative
Kodaira dimension.  If $K_X$ is big, restrict a rational decomposition
of a multiple into an ample class and an effective divisor.  The
ample pullback is big on the generic parameter fibre, since evaluation
is generically finite on its image.  Its canonical class is then big.
\end{proof}

\subsection{Small adjoint volumes on Iitaka fibres}

The tracking lemma produces a moving member of dimension at most three
whose polarization has small volume and whose canonical degree is bounded
by that volume.  The next lemma isolates the lower-dimensional argument:
its Iitaka fibres have Kodaira dimension zero, and adding the restricted
polarization to their canonical divisor still gives small volume.

\begin{lemma}\label{lem:scalar-iitaka-volume}
Fix $1\le d\le3$ and a positive rational number $C'$.  Let
$(W_i,L_i)$ be a sequence of smooth projective $d$-folds with
$\kappa(W_i)\ge0$ and big, nef, semiample rational Cartier divisors $L_i$.
Suppose that
\begin{equation}\label{eq:scalar-fibre-hypotheses}
 L_i^d\longrightarrow0,
 \qquad K_{W_i}\cdot L_i^{d-1}\le C'L_i^d.
\end{equation}
After passing to a subsequence and resolving the Iitaka fibrations,
their smooth geometric generic fibres $U_i$ have a fixed positive
dimension, satisfy $\kappa(U_i)=0$, and obey
\[
 \vol(K_{U_i}+L_i|_{U_i})\longrightarrow0.
\]
Here $L_i$ denotes its pullback to the resolved model; its restriction
to $U_i$ is big, nef, and semiample.  When $\kappa(W_i)=0$, the fibre
$U_i$ means $W_i$ itself.
\end{lemma}

\begin{proof}
Suppress the sequence index.  We first bound the volume of every fixed
positive adjoint combination on $W$, then compare it with the volume on
an Iitaka fibre.
For every fixed positive rational number $a$, we claim that
\begin{equation}\label{eq:scalar-small-adjoint-volume}
 \vol(aK_W+L)\longrightarrow0.
\end{equation}
Choose a general divided member of a sufficiently divisible free
multiple of $L/a$.  It gives an effective klt boundary
$\Gamma\sim_{\Q}L/a$.  This boundary is big, and so is
$K_W+\Gamma$, since $\kappa(W)\ge0$.
The klt log-general-type model and finite-generation theorem
\cite[Theorem~1.2]{BCHM2010} gives a $\Q$-factorial klt log terminal model.
Its nef big adjoint is semiample by the klt basepoint-free theorem
\cite[Theorem~3.3]{KollarMori1998}: apply it to a Cartier multiple of
the adjoint, whose sufficiently large multiple minus the adjoint is
again nef and big.

On a common smooth resolution, let $P_a$ be $a$ times the pullback of
this semiample adjoint.  The negative-map comparison, using the
negativity lemma \cite[Lemma~3.39]{KollarMori1998}, gives
\[
 aK_W+L\sim_{\Q}P_a+F_a,
 \qquad F_a\ge0
\]
after pullback, with $F_a$ exceptional over the model.  Sections in
divisible degrees are unchanged by this exceptional part, so
$P_a^d=\vol(aK_W+L)$.  All occurrences of $L$ on the resolution mean
its nef pullback.  The Khovanskii--Teissier inequality gives
\[
 (aK_W+L)\cdot L^{d-1}
 \ge P_a\cdot L^{d-1}
 \ge (P_a^d)^{1/d}(L^d)^{(d-1)/d}.
\]
For completeness, the nef mixed-intersection inequalities used here
follow by ample perturbation and passage to the limit from the
surface Hodge index theorem, applied on general complete-intersection
surfaces.  Combining with \eqref{eq:scalar-fibre-hypotheses}, we
obtain the explicit bound
\begin{equation}\label{eq:scalar-volume-bound}
 \vol(aK_W+L)\le (aC'+1)^dL^d,
\end{equation}
which proves the claim.

Pass to a subsequence on which $j=\kappa(W)$ is fixed.
Since $d\le3$, effective Iitaka fibrations in the zero-boundary case
give an integer $m_0=m_0(d)>0$ such that $|m_0K_W|$ gives the Iitaka
fibration; see \cite[Corollary~1.5]{ChenHanLiu2023}.
The case $j=d$ is impossible for large sequence index.  Resolving
that linear system would produce an integral free moving part with
generically finite map and top self-intersection at least one.  Thus
$\vol(K_W)\ge m_0^{-d}$, whereas
$\vol(K_W+L)\to0$ by \eqref{eq:scalar-small-adjoint-volume} and
$L$ has effective multiples.

If $0<j<d$, resolve the Iitaka system, its Stein base, and the induced
map to obtain a morphism of smooth projective varieties
\[
 f:W\longrightarrow T
\]
with connected fibres.  There is an integral big basepoint-free divisor
$H$ on $T$, pulled back from the hyperplane class of the system's image,
such that
\begin{equation}\label{eq:scalar-iitaka-hyperplane}
 f^*H\preceq m_0K_W,
 \qquad H^j\ge1.
\end{equation}
We continue to write $W$ for the resulting smooth model and $L$ for
its pullback.  The previous intersection estimates persist: canonical
divisors under further smooth resolutions acquire effective exceptional
terms, and these pair to zero against $L^{d-1}$.

Let $U$ be a smooth integral geometric generic fibre of $f$ and put
$g=d-j>0$.  Then $K_W|_U\sim K_U$, and $L|_U$ is big because these
fibres move.  We recall explicitly that $\kappa(U)=0$.  A nonzero
canonical multiple on $W$ restricts nontrivially to the generic fibre,
so $\kappa(U)\ge0$.  If a system $|kK_U|$ had positive-dimensional
image, its sections would extend over the generic base after twisting
$kK_W$ by $vf^*H$ for a sufficiently large $v$.  Indeed the coherent
direct image of $kK_W$ has the required generic fibre of sections,
and a sufficiently large multiple of the big divisor $H$ makes it
generically generated, using a decomposition into ample and effective
parts.  By \eqref{eq:scalar-iitaka-hyperplane}, the twist is absorbed
into a further canonical multiple.  The resulting section ratio is
nonconstant on the geometric generic fibre and hence transcendental
over $\C(T)$.  Products with sections of $|m_0K_W|$ also preserve the
$j$ independent base ratios.  This would force $\kappa(W)>j$, a
contradiction.  When $j=0$, simply put $U=W$ and $g=d$.

We next show that
\begin{equation}\label{eq:scalar-small-fibre-volume}
 \vol(K_U+L|_U)\longrightarrow0.
\end{equation}
For $j=0$ this is \eqref{eq:scalar-small-adjoint-volume}.  If $j>0$
and $g=1$, then $\kappa(U)=0$ implies $\deg K_U=0$, and the volume
in question equals $\deg(L|_U)$.  Set $A=f^*H$.  Nef mixed-intersection
log-concavity and \eqref{eq:scalar-iitaka-hyperplane} give
\[
 A^jL^{d-j}
 \le\left(\frac{AL^{d-1}}{L^d}\right)^jL^d
 \le (m_0C')^jL^d\longrightarrow0.
\]
One may add a positive multiple of $L$ to $A$ and pass to the limit
if a mixed intersection vanishes.  The projection formula and
$H^j\ge1$ show that this bounds $\deg(L|_U)$ from above.

The only remaining case is $d=3$, $j=1$, and $g=2$.  Now $T$ is a
smooth curve and $H$ is ample.  A basepoint-free pencil in $|H|$ and
its ramification divisor give $K_T+2H\succeq0$.  Consequently
\begin{equation}\label{eq:scalar-relative-domination}
 K_{W/T}+L+2f^*H\preceq(1+4m_0)K_W+L.
\end{equation}
Choose an effective klt rational boundary linearly equivalent to $L$
from its semiample system.  The twisted weak positivity theorem
\cite[Theorem~1.1]{FujinoWeakPositivity} applies to $f$ with this
boundary.  Thus for sufficiently divisible $l$ the locally free sheaf
\[
 E_l=f_*\OO_W\bigl(l(K_{W/T}+L)\bigr)
\]
is weakly positive.  Here smoothness, projectivity, connected fibres,
and effectivity of the klt boundary are all satisfied; on the smooth
curve a torsion-free direct image is locally free.

Choose $l\ge1$, separately for each member of the sequence, so that
the $l$-th Veronese of the section algebra of $K_U+L|_U$ is generated
in degree one.  Such an $l$ exists by the same big klt adjoint model
and finite-generation theorem used above.  Enlarge it to clear the
chosen rational linear equivalence.  Weak positivity gives $b>0$
for which
\[
 \Sym^bE_l\otimes\OO_T(bH)
\]
is generated by global sections at the generic point.  For every
positive multiple $q$ of $b$, symmetric multiplication and the
degree-one generation on the generic fibre therefore give global
sections of
\[
 ql(K_{W/T}+L)+qf^*H
\]
spanning the degree-$ql$ fibre sections.  Select a basis over
$\C(T)$ among their restrictions, of cardinality
$h^0(U,ql(K_U+L|_U))$.  Multiplying by a basis of $H^0(T,qlH)$
gives linearly independent products over $\C$: any relation can
first be tested using independence over $\C(T)$ and then independence
of the base sections.  Multiplication by a nonzero section of
$(ql-q)H$ places these products in degree $ql$ of
$K_{W/T}+L+2f^*H$.

Curve Riemann--Roch and the asymptotic fibre section formula, with
$l$ fixed and $q\to\infty$ through these multiples, yield
\begin{align}
 \vol\bigl((1+4m_0)K_W+L\bigr)
 &\ge\vol(K_{W/T}+L+2f^*H)\notag\\
 &\ge\frac{d!}{g!}\deg(H)\vol(K_U+L|_U).
 \label{eq:scalar-weak-positive-volume}
\end{align}
The first expression tends to zero by
\eqref{eq:scalar-small-adjoint-volume}, and $\deg H\ge1$.
This proves \eqref{eq:scalar-small-fibre-volume} in every case.

\end{proof}

\subsection{The upper order estimate}

\begin{proof}[Proof of the upper bound in Theorem~\ref{thm:scalar-bounds}]
If no uniform $C$ exists, rational rescaling gives a sequence of the
data in the theorem, with polarizations denoted $M$, such that
\begin{equation}\label{eq:scalar-small-volume}
 M^4\longrightarrow0,
 \qquad \gamma(M;V)>1.
\end{equation}
The proof decreases the dimension while retaining small volume and
an effective canonical upper bound. At each stage we first extract a
moving base component, then pass to a Kodaira-dimension-zero fibre of
its Iitaka fibration, and finally replace a big adjoint by its semiample
model. The process stops when its section-order invariant tends to zero.
Tracing the effective divisor comparisons back will construct covering
families of positive-dimensional proper subvarieties $A\subset V$ whose
geometric generic members satisfy
\begin{equation}\label{eq:scalar-small-zero-members}
 \kappa(A^*)=0,
 \qquad \gamma(M;A)\longrightarrow0.
\end{equation}

All constructions can be made successively on geometric generic
members.  To be precise, choose countable fields of definition for
the sequence.  A geometric generic datum over such a field can be
realized over $\C$ to apply the model and positivity theorems used here.
The resulting finite collection of varieties, morphisms, and divisors
descends again to a field finitely generated over the preceding
generic parameter field.  Spreading it out therefore gives a dominant
extension of the preceding parameter space: the old generic parameter
is never specialized.  Formation of section spaces and of their jet
kernels commutes with field extension, so the generic all-degree order
comparisons are preserved.  There are only finitely many stages for
each variety, and countably many varieties in the sequence.

Begin with a small projective $\Q$-factorialization $Z\to V$, which
exists for the klt variety $V$ by \cite[Corollary~1.4.3]{BCHM2010},
and let $N$ be the pullback of $M$ to $Z$.  We describe a stage of the construction.  After
passing to a subsequence, its dimension $n$ is fixed and the following
properties hold:
\begin{enumerate}[label=(\roman*)]
 \item $Z$ is projective $\Q$-factorial klt and $N$ is big, nef, and
       semiample, with $N^n\to0$;
 \item a smooth projective model of $Z$ has nonnegative Kodaira dimension;
 \item $K_Z\preceq a_0N$ for a fixed positive rational number $a_0$;
 \item $\gamma(N;Z)$ is bounded below by a positive constant.
\end{enumerate}
Initially these assertions follow from \eqref{eq:scalar-small-volume},
canonicality, and crepantness of the small morphism.  If $n=1$, they
cannot all hold, since the order of a section on a smooth curve is
bounded by its degree.  We may therefore assume $n\ge2$.

Choose $n+1$ fixed rational jet levels, equally spaced by $\delta>0$,
below the lower bound in (iv).  For all sufficiently large sequence indices they are above
$(N^n)^{1/n}$.  Apply Lemma~\ref{lem:scalar-tracking}.  Write
$W=S^*$ for a resolution of its geometric generic moving member and
$L=N|_W$, and pass to a subsequence with fixed $d=\dim W<n$.
The divisor $L$ is nef and semiample, and is big because a sweeping
member meets the open set on which the semiample map of $N$ is
generically finite.  Lemma~\ref{lem:scalar-covering} gives
$\kappa(W)\ge0$.  Restricting the effective comparison in (iii) to
the moving member and using the tracking bounds yields
\begin{equation}\label{eq:scalar-small-intersections}
 L^d\longrightarrow0,
 \qquad K_W\cdot L^{d-1}\le C' L^d
\end{equation}
with the fixed constant $C'=a_0+2n/\delta$ at this stage.

Apply Lemma~\ref{lem:scalar-iitaka-volume} to these moving members.
After passing to a subsequence and resolving their Iitaka systems,
let $U$ be the smooth geometric generic Iitaka fibre, with $U=W$ when
$\kappa(W)=0$, and put $g=\dim U$.  Then
\[
 0<g\le d<n,\qquad \kappa(U)=0,\qquad
 \vol(K_U+L|_U)\longrightarrow0.
\]
The restricted divisor $L|_U$ is big, nef, and semiample.  The generic
field and spreading convention above keeps these fibres in covering
families over the preceding parameter space.

To continue the dimension descent, we need a semiample polarization
that still bounds the original restricted sections from above.
Take an effective klt boundary $\Gamma\sim_{\Q}L|_U$ from the free
system, and take a big klt good model of $K_U+\Gamma$, as in the proof
of Lemma~\ref{lem:scalar-iitaka-volume}.
Write $Z'$ for its $\Q$-factorial klt underlying variety and $N'$
for its big semiample adjoint.  Then
\[
 (N')^g=\vol(K_U+L|_U)\longrightarrow0.
\]
On a smooth common resolution with maps $a$ to $U$ and $b$ to $Z'$,
the negative-map comparison reads
\begin{equation}\label{eq:scalar-good-model-comparison}
 a^*(K_U+\Gamma)=b^*N'+F,
 \qquad F\ge0\text{ exceptional over }Z'.
\end{equation}
For any effective $\Delta_U\sim_{\Q}K_U$, we have
\begin{equation}\label{eq:scalar-fixed-part-bound}
 F\le a^*\Delta_U.
\end{equation}
Indeed, for sufficiently divisible $m$, every effective divisor in
$|mb^*N'+mF|$ contains $mF$.  Its pushforward is an effective divisor
in $|mN'|$ on the normal variety $Z'$, and pulling that Cartier divisor
back leaves exactly the exceptional summand $mF$.  But
$ma^*\Delta_U+Y$ belongs to this system when $Y$ is a general member
of the free system $|ma^*(L|_U)|$.  Choose $Y$ avoiding every component
of $F$.  Comparing coefficients proves
\eqref{eq:scalar-fixed-part-bound}.  In particular,
\begin{equation}\label{eq:scalar-polarization-domination}
 a^*(L|_U)\preceq b^*N'.
\end{equation}

If $\gamma(N';Z')\to0$ on a subsequence, we stop.  Otherwise, on a
subsequence it is bounded below by a positive constant, and we repeat
with $(Z',N')$.  The four stage conditions hold: $K_{Z'}\preceq N'$
because the pushed-forward boundary is effective, and a smooth model
of $Z'$ is birational to $U$, which has Kodaira dimension zero.
The dimension has decreased strictly, from $n$ to $g\le d<n$.
The construction must stop, since on a curve maximum normalized
section order is at most the degree.

We verify that the stopping condition proves
\eqref{eq:scalar-small-zero-members} for the original polarization.
Every moving member and fibre used above spreads to a covering
family over the preceding generic parameter space.  Its map to its
image in that space, and ultimately in $V$, is birational: a moving
member meets the open set of isomorphism of each relevant birational
model map.  For a subsequent member in $Z'$, take its strict transform
on the common resolution, so it has an actual morphism back to $U$.
The effective comparisons \eqref{eq:scalar-polarization-domination}
restrict to these members, since they are not contained in the extra
divisors.  Their successive restrictions compare the pullback of the
original $M$ with the final $N'$ on a common resolution.  Multiplication
by the effective differences injects section spaces in divisible
degrees, and at a generic marked point those differences do not
vanish.  Birational pullback between normal projective models preserves
the section spaces.  Thus the final image $A\subset V$ satisfies
\[
 \gamma(M;A)\le\gamma(N';Z')\longrightarrow0,
 \qquad \kappa(A^*)=\kappa(U)=0.
\]
Composing the dominant incidence evaluations proves that these images
form a covering family.  They have positive dimension and are proper
because the first tracking step already decreased dimension.

Mark a generic smooth point on these members and add every $G$-translate
of the marked family.  All data can be defined over one countable
field; invariance of $M$ preserves the order bound.
Lemma~\ref{lem:generic-chain} gives a $G$-equivariant rational quotient,
and Lemma~\ref{lem:chain-order} bounds the section order on its
geometric generic leaf by its dimension times $B$, where $B\to0$.
If the leaf has dimension four, this gives $\gamma(M;V)\le4B$, contradicting
\eqref{eq:scalar-small-volume}.  Otherwise it gives a $G$-equivariant
rational fibration with positive-dimensional base and fibre.  The
geometric generic leaf is of general type by hypothesis.

That leaf is nevertheless covered by members of Kodaira dimension
zero.  To see that the assertion is about generic members, restrict
the marked incidence to its generic leaf name.  A component still
dominates the leaf by marked-point evaluation.  Its generic parameter
is generic in the original parameter space over the original field,
because it is obtained by geometric extension over the generic name.
The quotient is generically constant on each such member, and its
strict transforms give a covering family on a resolution of the leaf.
This contradicts Lemma~\ref{lem:scalar-covering}.  The upper bound
follows.
\end{proof}

\subsection{The lower curve estimate}

\begin{proof}[Proof of the lower bound in Theorem~\ref{thm:scalar-bounds}]
Rescale rationally so that $1\le P^4\le2$.  If no uniform positive
lower bound exists, there is a sequence of marked covering families
of integral curves with
\[
 \frac{P\cdot D}{\mult_xD}\le B,
 \qquad B\longrightarrow0.
\]
Here the marks can be chosen generic over countable fields of
definition.  Indeed, failure of a very general bound still gives a
bad mark after excluding the countably many proper closed subsets
over such a field.  Choose a curve witnessing a strict upper bound
there and spread the curve and its mark over their field of definition.
The mark evaluation is dominant, and degree and multiplicity agree
with those of the geometric generic data.  Include all $G$-translates
if necessary; they obey the same ratio bound.

Take the rational quotient from Lemma~\ref{lem:generic-chain}.
After passing to a subsequence with fixed leaf dimension $h>0$,
Lemma~\ref{lem:chain-order} gives, on its geometric generic leaf $H$,
\begin{equation}\label{eq:scalar-small-leaf}
 P^h\cdot H\le(hB)^h.
\end{equation}
For large sequence index it cannot have dimension four, since
$P^4\ge1$.  Put $s=4-h>0$.  Resolve the leaf fibration to a morphism
$f:X\to T$ between smooth projective varieties with geometrically
integral generic fibre, and let $\pi:X\to V$ be the resulting
birational morphism.  Continue to write $P$ for $\pi^*P$.

Choose a smooth integral very general fibre $F$ of $f$.  It can be
chosen not contained in any of the countably many proper closed sets
excluded by the all-degree upper order bound, nor in the exceptional
locus of $\pi$.  For clarity, over a smooth open of the base, each
proper closed set of $X$ contains the whole fibre only over a proper
closed subset of $T$, by properness and upper semicontinuity of fibre
dimension.  Exclude these countably many subsets first, then choose
a point $x\in F$ outside their intersections with $F$.  Thus $\pi$
is an isomorphism near $x$ and
\begin{equation}\label{eq:scalar-point-upper}
 \gamma(P;V)\le C2^{1/4}
\end{equation}
holds at $\pi(x)$ in every divisible degree.  The degree
$P^h\cdot F$ equals the geometric generic degree in
\eqref{eq:scalar-small-leaf}.

Let $I=I_F\subset\OO_X$.  Since $F$ is a smooth fibre over a smooth
point of the $s$-dimensional base, it is regularly embedded and has
trivial conormal bundle of rank $s$.  Consequently
\[
 I^k/I^{k+1}\simeq
 \Sym^k(\OO_F^{\oplus s})
 \simeq\OO_F^{\oplus\binom{k+s-1}{s-1}}.
\]
The filtration by powers of $I$ therefore bounds the rank of
restriction to the $u$-th thickening, for every $u\ge1$ and every
divisible $l$, by
\begin{equation}\label{eq:scalar-thickening}
 h^0\bigl(X,\OO_X(lP)\otimes(\OO_X/I^u)\bigr)
 \le\binom{u+s-1}{s}\,h^0(F,lP|_F).
\end{equation}
No vanishing theorem uniform in $u$ is required for this inequality.

Fix $D>C2^{1/4}$ and set $u=\lceil Dl\rceil$.  The pullback of $P$
is big and semiample on $X$, and its restriction to $F$ is big and
semiample because $F$ meets the birational isomorphism locus.
The semiample Hilbert-polynomial asymptotics show that the leading
coefficient, after division by $l^4$, on the right-hand side of
\eqref{eq:scalar-thickening} is
\[
 \frac{D^s(P^h\cdot F)}{s!\,h!}
 \le\frac{D^s(hB)^h}{s!\,h!}.
\]
Choose the sequence index first so that this is strictly less than
$1/4!$, hence less than $P^4/4!$.  Then choose a sufficiently large
divisible $l$ for this fixed example.  Since
\[
 h^0(X,lP)=\frac{P^4}{4!}l^4+O(l^3),
\]
restriction to the thickening has a nonzero kernel.  A kernel section
lies in $I^u$, and $I_x^u\subset\mathfrak m_x^u$.  It therefore has
order at least $u\ge Dl$ at $x$.  Normality gives
$\pi_*\OO_X=\OO_V$, so it is a section downstairs as well, with
the same order at the point where $\pi$ is an isomorphism.  This
contradicts \eqref{eq:scalar-point-upper}.  Rescaling back proves the
lower bound, and completes the theorem.
\end{proof}

\section{Quadratic jet cost for diagonal cyclic products}
\label{sec:quadratic-cost}

We prove a bound for diagonal products of the canonical cyclic cover.  The
constant in the bound depends only on an upper bound for normalized section
orders on the original fourfold.  In particular, it is independent of the
order of the cover.

The diagonal rank-deficit and determinant-order strategy adapts the
Frobenius construction of \cite[Section~9]{OpenAIAbundance}.
The cyclic many-factor weights, estimates, and constants needed here
are supplied below; the abundance conclusion of that reference is not
an input to this argument.

\begin{proposition}\label{prop:quadratic-cost}
Let $X$ be a normal projective canonical fourfold over $\C$ with
$K_X\sim_{\Q}0$, and let $r$ be the least positive integer such that
$rK_X\sim0$.  Let $\pi:Y\to X$ be its canonical cyclic cover, with group
$\mu_r$.  Fix $C_1\in\Q_{>0}$, and choose $b_0\in\Q_{>0}$ satisfying
\begin{equation}\label{eq:frob-choice-b0}
 3(2b_0)^4<\frac14,
 \qquad b_0(C_1+1)<\frac34.
\end{equation}
Suppose that $L$ is an ample rational Cartier divisor on $X$ such that
\[
 1\le L^4\le2,
 \qquad \gamma(L;X)\le C_1,
\]
where $\gamma$ is as in Definition~\ref{def:orders}.  For each integer $t\ge2$,
set
\[
 Z_t=Y^t/\mu_{r,\mathrm{diag}},
 \qquad
 \theta_t:Z_t\longrightarrow X^t,
 \qquad
 P_t=\theta_t^*\left(\sum_{i=1}^t\operatorname{pr}_i^*L\right).
\]
Then $P_t$ is ample and, at a very general point $z$ of $Z_t$,
\begin{equation}\label{eq:frob-cost}
 \varepsilon(P_t;z)\le\frac{5t^2}{b_0}.
\end{equation}
\end{proposition}

Suppose the asserted curve bound fails. The resulting surplus of
ordinary jets will produce, after reduction to characteristic $p$, a
map of vector bundles that has full rank at a general product point.
On the diagonal its rank is at most one quarter of the target rank.
A nonzero maximal minor must consequently vanish to high order there.
The section-order bound downstairs will give a strictly smaller upper
bound for that same minor at one diagonal point.

The next three lemmas supply the estimates in this comparison. The
truncated-power slope bound controls the graded pieces on the
Frobenius diagonal. The fixed flag converts those slopes into a bound
for their spaces of sections, hence for the diagonal rank. The
external-product order lemma bounds the minor by the sum of its slot
orders and rules out a larger order caused by cancellation.
We first establish these estimates on vector bundles.  For a vector bundle $E$ on a smooth
projective curve, write $\mu_{\min}(E)$ and $\mu_{\max}(E)$ for the smallest
and largest slopes in its Harder--Narasimhan filtration.

\begin{lemma}\label{lem:frob-truncated-slopes}
Let $C$ be a smooth projective integral curve of genus $g$ over an
algebraically closed field $k$ of characteristic $p>0$.  Let $V$ be a vector
bundle of rank $n$ with $\mu_{\min}(V)\ge0$, and put
\[
 G=\max\{0,2g-2\},\qquad u=n(p-1),\qquad c_2=n(n-1)G.
\]
Let $A_\bullet(V)$ denote the symmetric algebra of $V$ modulo the ideal
locally generated by the $p$-th powers of its linear generators.  Thus, in
a local frame, $A_\bullet(V)$ is the graded algebra
$\OO_C[v_1,\ldots,v_n]/(v_1^p,\ldots,v_n^p)$.  Its graded pieces are locally
free, it is zero in degrees above $u$, and
\begin{equation}\label{eq:frob-truncated-upper-slope}
 \mu_{\max}\bigl(A_j(V)\bigr)
 \le (p-1)\deg(V)+c_2
 \qquad (0\le j\le u).
\end{equation}
\end{lemma}

\begin{proof}
We use the canonical-connection approach to Frobenius instability;
see Langer \cite[\S2, Corollaries~2.4 and~6.2]{Langer2004}.
We give the estimate and tensor argument explicitly to retain the
constants needed for the many-factor diagonal.
Let $F_C$ denote absolute Frobenius.  For every vector bundle $U$ of rank
$n$, we first claim that
\begin{equation}\label{eq:frob-one-step-slope}
 \mu_{\min}(F_C^*U)
 \ge p\mu_{\min}(U)-(n-1)G.
\end{equation}
The pullback $F_C^*U$ carries the connection given by componentwise
differentiation in pulled-back frames: the transition matrices have
$p$-th-power entries.  Write the strictly decreasing slopes of its
Harder--Narasimhan factors as $\nu_1>\cdots>\nu_a$.  If some successive gap
is greater than $G$, take the last such gap, and let $S$ be the saturated
filtration step immediately before that gap.  Then
\[
 \mu_{\min}(S)>
 \mu_{\max}\bigl((F_C^*U)/S\bigr)+\deg(\Omega_C^1).
\]
Consequently the second fundamental map
\[
 S\longrightarrow ((F_C^*U)/S)\otimes\Omega_C^1
\]
vanishes, so $S$ is preserved by the connection.

Here this invariance gives Frobenius descent directly.  At the function
field $K=k(C)$, choose a matrix whose columns form
a basis of the subspace $S_K$ and whose rows in a suitable index set form
the identity matrix.  Differentiating its columns gives vectors in
$S_K\otimes\Omega^1_{K/k}$.  Their identity rows are zero, so these
differentiated columns are all zero.  Every matrix entry therefore lies
in $K^p$.  Indeed $[K:K^p]=p$ for a one-variable function field over a
perfect field, and the nonzero derivation $d:K\to\Omega^1_{K/k}$ has
kernel $K^p$.  Extracting roots of the matrix entries produces a subspace
of $U_K$.  Let $S_0\subset U$ be its saturation.  Since $C$ is smooth,
$U/S_0$ is locally free.  Frobenius is flat, so $F_C^*S_0$ is saturated
in $F_C^*U$.  It has the same generic fibre as $S$, hence equals $S$.

Put $Q=(F_C^*U)/S$.  By descent it is the Frobenius pullback of a locally
free quotient of $U$, and therefore
$\mu(Q)\ge p\mu_{\min}(U)$.  By the choice of the last large gap, all
successive Harder--Narasimhan slope gaps of $Q$ are at most $G$.  Its
average slope exceeds its smallest slope by at most $(n-1)G$.
Moreover $\mu_{\min}(Q)=\mu_{\min}(F_C^*U)$.  These observations prove
\eqref{eq:frob-one-step-slope}.  If there is no gap greater than $G$, the
same argument applies with $S=0$ and $Q=F_C^*U$.

Iteration, using $\mu_{\min}(V)\ge0$, gives for every $e\ge1$
\begin{equation}\label{eq:frob-iterated-slope}
 \mu_{\min}\bigl((F_C^e)^*V\bigr)
 \ge -(n-1)G\frac{p^e-1}{p-1}.
\end{equation}
We next deduce a tensor estimate without assuming tensor semistability in
characteristic $p$.  For each $e$, choose a line bundle $T_e$ of integral
degree $d_e$ such that
\[
 2g-\mu_{\min}\bigl((F_C^e)^*V\bigr)
 \le d_e<
 2g+1-\mu_{\min}\bigl((F_C^e)^*V\bigr).
\]
The bundle $V_e=(F_C^e)^*V\otimes T_e$ is globally generated.  In fact,
after subtracting any point its minimum slope is at least $2g-1$, so
Serre duality and the slope criterion for a nonzero homomorphism give
$H^1(C,V_e(-x))=0$; evaluation at $x$ is therefore surjective.

Let $Q_i$ be any positive-rank torsion-free quotient of $V^{\otimes i}$.
Its pullback, twisted by $T_e^{\otimes i}$, is a quotient of
$V_e^{\otimes i}$, hence is globally generated and has nonnegative
degree.  Thus
\[
 p^e\mu(Q_i)+i d_e\ge0.
\]
Divide by $p^e$ and let $e$ tend to infinity.  The upper bound for $d_e$
and \eqref{eq:frob-iterated-slope} imply
\begin{equation}\label{eq:frob-tensor-minimum}
 \mu(Q_i)\ge-\frac{i(n-1)G}{p-1}.
\end{equation}
Since $A_i(V)$ is a quotient of $V^{\otimes i}$, the same lower bound
holds for its minimum slope.  For $0\le i\le u$, it is at least $-c_2$.

Multiplication in the truncated algebra is a perfect pairing
\[
 A_j(V)\otimes A_{u-j}(V)\longrightarrow A_u(V).
\]
In a frame, each monomial pairs with its unique complementary exponent
monomial, with coefficient one.  The top line has transition character
$(\det V)^{p-1}$: this follows for diagonal changes of frame and for
elementary changes of frame, which generate the general linear group.
It follows that
\[
 A_j(V)^*\otimes(\det V)^{p-1}\simeq A_{u-j}(V).
\]
Its minimum slope is at least $-c_2$.  Slope duality now gives
\eqref{eq:frob-truncated-upper-slope}.
\end{proof}

\begin{lemma}\label{lem:frob-flag-sections}
Let $W$ be a smooth projective fourfold over an algebraically closed
field, and let $H$ be an ample rational Cartier divisor.  Fix a smooth
complete intersection flag, cut successively by divisors from
$|l_1H|$, $|l_2H|$, and $|l_3H|$, ending in a smooth integral curve $C$;
here each $l_iH$ is an integral ample divisor.  Write
$d_i=l_iH\cdot C$.  If $E$ is a vector bundle of rank $e$ and
$\mu_{\max}(E|_C)\le M$ for a real number $M\ge0$, then
\begin{equation}\label{eq:frob-flag-sections}
 h^0(W,E)\le e(M+1)\prod_{i=1}^3\left(1+\frac{M}{d_i}\right).
\end{equation}
\end{lemma}

\begin{proof}
Successive applications of the divisor exact sequences along the flag
bound $h^0(W,E)$ by
\begin{equation}\label{eq:frob-flag-sum}
 \sum_{k_1,k_2,k_3\ge0}
 h^0\left(C,E|_C\otimes
 \OO_C\left(-\sum_{i=1}^3 k_i l_iH\right)\right).
\end{equation}
To justify the iterations, at each flag stage and with the preceding
twists fixed, sufficiently negative ample twists have no sections.  A
vector bundle on that stage embeds into a finite direct sum of copies
of a sufficiently positive ample line bundle, by dualizing a global
generation map for its dual.  This gives the required vanishing and
allows each iteration of the exact sequence to terminate.  Enlarging
the resulting finite sums to \eqref{eq:frob-flag-sum} can only increase
the bound.

A summand vanishes when $\sum_i k_i d_i>M$, because its maximum slope is
then negative.  Each summand is at most $e(M+1)$: evaluation at
$\lfloor M\rfloor+1$ distinct points is injective, since subtracting
these points makes the maximum slope negative even before the additional
nonpositive twists.  There are at most
$\prod_i(1+M/d_i)$ possible triples with $\sum_i k_i d_i\le M$.
This proves \eqref{eq:frob-flag-sections}.
\end{proof}

\begin{lemma}\label{lem:frob-external-orders}
For $1\le i\le t$, let $V_i$ be an integral projective variety over an
algebraically closed field, let $x_i\in V_i$ be a smooth point, and let
$M_i$ be a line bundle with $H^0(V_i,M_i)\ne0$.  Put
\[
 a_i=\max_{0\ne s\in H^0(V_i,M_i)}\ord_{x_i}(s).
\]
Every nonzero section $s$ of the external product
$\bigotimes_i\operatorname{pr}_i^*M_i$ satisfies
\[
 \ord_{(x_1,\ldots,x_t)}(s)\le\sum_{i=1}^t a_i.
\]
\end{lemma}

\begin{proof}
Each space of sections is finite dimensional, and its order filtration
is separated; the numbers $a_i$ are therefore finite.  Choose a basis
adapted to this filtration in each slot.  The initial forms of basis
elements having a fixed order are linearly independent in that degree
of the associated graded local ring.  By the K\"unneth formula, every
section of the external product is a linear combination of the tensor
products of these bases.  In the least total degree occurring in a
nonzero combination, group its initial forms by their slot multidegrees.
Different multidegrees cannot cancel.  Within each multidegree the
products are independent, since they are tensor products of independent
initial forms in disjoint sets of smooth parameters.  The section's
order is consequently the least total degree of one of its nonzero
basis coefficients, and this is at most $\sum_i a_i$.
\end{proof}

\begin{proof}[Proof of Proposition~\ref{prop:quadratic-cost}]
The map $\theta_t$ is finite, so $P_t$ is ample.  Fix $X,L,r$, and $t$;
all auxiliary data below may depend on these choices.  Set
\begin{equation}\label{eq:frob-bq}
 b=\frac{b_0}{t},\qquad q=2tr.
\end{equation}
Suppose, for a contradiction, that
$\varepsilon(P_t;z)>5t/b$ at a smooth very general point $z$, whose image
in every $X$-slot lies in $X_{\mathrm{reg}}$.

\emph{Jets and the canonical torsor.}
Choose a rational number $c$ with
\[
 5t<c<b\varepsilon(P_t;z).
\]
On the blow-up $\sigma:\widetilde Z_t\to Z_t$ at $z$, with exceptional
divisor $E_z$, the rational divisor $\sigma^*(bP_t)-cE_z$ is ample.  To
see this, choose a larger coefficient below
$b\varepsilon(P_t;z)$ for which the corresponding divisor is nef, and
interpolate it with the ample divisor obtained at a sufficiently small
positive coefficient.  Serre vanishing, applied to sufficiently large
divisible multiples of $\sigma^*(bP_t)-cE_z$, makes the restriction map
to the corresponding thick exceptional divisor surjective.  Pullback
identifies sections of $k_0bP_t$ with sections of its pullback, by
normality of $Z_t$.  At the smooth point $z$, ordinary jets of order
strictly less than $k_0c$ inject into the sections on $k_0cE_z$.
It follows, after choosing $k_0$ large and divisible, that
\begin{equation}\label{eq:frob-initial-jets}
 \begin{gathered}
 Q=k_0bL\text{ is integral Cartier},\\
 H^0(Z_t,k_0bP_t)\longrightarrow
 \OO_{Z_t,z}/\mathfrak m_z^{5tk_0+1}
 \text{ is surjective},
 \end{gathered}
\end{equation}
where a local frame is understood in the displayed evaluation map.

Fix a section $\eta$ trivializing $rK_X$.  Over $X_{\mathrm{reg}}$,
the cover $Y\to X$ is the torsor of frames $\tau$ of the canonical line
with $\tau^r=\eta$.  Over tuples in this open locus the diagonal action
is free, and $Y^t\to Z_t$ is etale.  Choose a lift $\widetilde z$ of $z$.
Pulling back \eqref{eq:frob-initial-jets} preserves jet surjectivity at
$\widetilde z$.  The pulled-back sections are diagonal invariants.
By the K\"unneth formula and decomposition into characters, they are
linear combinations of pure tensors whose slot factors have the form
\begin{equation}\label{eq:frob-eigen-tensors}
 \begin{gathered}
 \tau^{-m_i}s_i,
 \qquad
 s_i\in H^0\bigl(X_{\mathrm{reg}},\OO(Q+m_iK_X)\bigr),
 \\
 0\le m_i<r,
 \qquad
 \sum_i m_i\equiv0\pmod r.
 \end{gathered}
\end{equation}
Indeed multiplication of an eigenvector by the corresponding power of
the tautological frame makes it invariant and hence descending to the
indicated line bundle.

Take a smooth projective resolution $u:W\to X$, isomorphic over
$X_{\mathrm{reg}}$, and write
\[
 K_W=u^*K_X+A,\qquad A\ge0\text{ exceptional}.
\]
We use the same letters $L,Q$ for their pullbacks to $W$.
Every $s_i$ in \eqref{eq:frob-eigen-tensors} extends to a section of
$Q+m_iK_W$.  In fact, its effective zero divisor extends by closure on
the normal variety $X$; this divisor is rational Cartier, since its
class is $Q+m_iK_X$.  Pulling it back and adding $m_iA$ gives the
effective divisor of the corresponding rational section on $W$.
Similarly $\eta$ extends to
\[
 \eta_W\in H^0(W,rK_W),\qquad \Div(\eta_W)=rA.
\]
Let $U=u^{-1}(X_{\mathrm{reg}})$.  On the torsor of roots of $\eta_W$
over $U$, a fixed finite list of the tensors
\eqref{eq:frob-eigen-tensors}, now using these sections on $W$, spans
the jets in \eqref{eq:frob-initial-jets} at the chosen lift.

\emph{A polarization, a flag, and a movable test.}
Fix an ample divisor $H_{\mathrm{amp}}$ on $W$.  Since $L$ is nef and
$K_WL^3=0$, a sufficiently small positive rational $\delta$ makes
$H=L+\delta H_{\mathrm{amp}}$ ample and ensures
\begin{equation}\label{eq:frob-polarization}
 H^4\le3,
 \qquad LH^3\le H^4,
 \qquad (q+t-1)K_WH^3\le b_0H^4.
\end{equation}
The equality $K_WL^3=0$ follows from $K_X\sim_{\Q}0$ and the
exceptionality of $A$.  The canonical class of $W$ is pseudoeffective.
By \cite[Theorem~2.1]{CampanaPaun2015}, $\Omega_W^1$ is generically
semipositive for this polarization.  Apply
\cite[Theorem~6.1]{MehtaRamanathan1982} to its finitely many
Harder--Narasimhan factors.  We can choose a smooth complete
intersection flag, cut successively by sufficiently positive integral
multiples $l_iH$, $1\le i\le3$, ending in a smooth integral curve $C$
such that
\begin{equation}\label{eq:frob-flag-nonnegative}
 \mu_{\min}(\Omega_W^1|_C)\ge0.
\end{equation}
The general curve avoids the codimension-two loci where the factors
or their filtration quotients fail to be locally free.  Their
restricted slopes are the original slopes multiplied by the same
positive factor $l_1l_2l_3$, so the restriction theorem gives
\eqref{eq:frob-flag-nonnegative}.

Choose also a very general point $x\in U$ at which the all-degree bound
defining $\gamma(L;X)\le C_1$ holds.  Let
$\pi_x:\widehat W\to W$ be its blow-up, with exceptional divisor $J$,
and put
\begin{equation}\label{eq:frob-test-divisor}
 D^+=bL+(q+1)K_W,\qquad d_0=b(C_1+1).
\end{equation}
The divisor $D^+$ is big.  In sufficiently divisible degrees its
section spaces agree with those of $bL$, after multiplication by a
fixed section whose zero divisor is exceptional over $X$.  This is
because $rK_X\sim0$, $A$ is effective and exceptional, and
$u_*\OO_W(E)=\OO_X$ for any effective integral exceptional divisor
$E$ on this resolution.  At $x$ these fixed sections have order zero.
Thus every nonzero section of a sufficiently divisible multiple
$mD^+$ has order at most $mbC_1$ at $x$.

It follows that $\pi_x^*D^+-d_0J$ is not pseudoeffective.  Otherwise,
interpolating it with the big divisor $\pi_x^*D^+$ would make
$\pi_x^*D^+-dJ$ big for a rational $d$ with $bC_1<d<d_0$.  An effective
sufficiently divisible multiple would give a section of a multiple of
$D^+$ with an excluded order at $x$.
The projective algebraic duality theorem of
\cite[Theorem~2.2, arXiv version]{BDPP2013} therefore supplies a strongly movable
class $\beta$ on $\widehat W$ such that
\begin{equation}\label{eq:frob-movable-test}
 D^+\cdot\beta<d_0J\cdot\beta,
 \qquad J\cdot\beta>0.
\end{equation}
Here and subsequently divisors on $W$ are pulled back when paired with
$\beta$.  More precisely, we may take
$\beta=v_*(A_1A_2A_3)$ for a smooth projective birational model
$v:T\to\widehat W$ and very ample divisors $A_i$ on $T$.  The dual cone
is the closure of the cone generated by such classes; strict negativity
on that cone yields negativity on one generator.  Also
$D^+\cdot\beta\ge0$, since $L$ is nef and $K_W$ is pseudoeffective.
The second inequality in \eqref{eq:frob-movable-test} follows from the
first and $d_0>0$.  In particular,
\begin{equation}\label{eq:frob-test-nonnegative}
 L\cdot\beta\ge0,\qquad K_W\cdot\beta\ge0.
\end{equation}

\emph{Reduction of the fixed data.}
Spread the preceding finite algebraic data over an integral finitely
generated $\Z$-subalgebra of $\C$.  This includes the smooth varieties,
the flag, the points and the blow-up, the morphism $v$ and its very ample
divisors, suitable integral multiples of the rational line bundles,
the sections used in \eqref{eq:frob-eigen-tensors}, the section
$\eta_W$, and the torsor and its marked lift over $U$.  After shrinking
the base, the relevant fibres of $W,C,T$, and $\widehat W$ are smooth,
projective, and geometrically integral; the testing morphism remains
dominant; and all required ampleness and intersection data persist.
The fixed finite jet evaluation remains surjective: it is a map of
finite locally free modules obtained by restriction to a fixed-order
infinitesimal neighbourhood of a section of a smooth morphism, and
surjectivity is open.  Invert $r$ and the finitely many denominators
used in the data.

We can also retain \eqref{eq:frob-flag-nonnegative} on every geometric
fibre of this smaller base.  To see this without imposing any
Frobenius semistability condition, choose a fixed relative ample twist
that makes the family $\Omega_W^1|_C$ relatively globally generated.
For every positive-rank quotient, its degree is then bounded below
uniformly, in terms of its rank and this twist.  Hence there are only
finitely many possible negative degrees and ranks, and therefore
finitely many Hilbert polynomials, for quotients of negative degree.
The corresponding relative Quot schemes are projective.  None has a
point on the geometric generic fibre: such a point would persist after
extension to the original complex field and contradict
\eqref{eq:frob-flag-nonnegative}.  This also excludes quotients with
torsion and negative degree, since their torsion-free quotients have
still smaller degree.  Remove the images of these finitely many Quot
schemes.  The remaining geometric fibres have
$\mu_{\min}(\Omega_W^1|_C)\ge0$.

There are geometric fibres in arbitrarily large characteristics $p$:
the base is nonempty, of finite type, and dominant over $\Spec\Z$.
We take such fibres over algebraic closures of residue fields and keep
the preceding notation.  The genus of $C$ and all the intersection
numbers used above are constant.  The class $\beta$ continues to pair
nonnegatively with every effective divisor: pull the divisor back
under the dominant morphism $v$ and intersect with $A_1A_2A_3$.
This argument applies also after a base-field twist.  Thus only the
fixed test, and not pseudoeffective-cone duality, is being specialized.

\emph{Full rank away from the diagonal.}
In one of these reductions, write $F:W\to W'$ for relative Frobenius,
where a prime denotes base-field twist.  It is finite flat of degree
$p^4$.  Put
\begin{equation}\label{eq:frob-ranks}
 \begin{gathered}
 N=\lfloor p/k_0\rfloor,
 \qquad B=NQ,
 \qquad \mathcal E=F_*\OO_W(B),\\
 \qquad s=p^4=\rk(\mathcal E),
 \qquad R=s^t.
 \end{gathered}
\end{equation}
Products of $N$ of the fixed jet tensors span jets through degree
$5tk_0N$ at the chosen lift.  Indeed every monomial of degree at most
$5tk_0N$ is a product of $N$ monomials of degree at most $5tk_0$.
Choose representatives with these leading terms from the original
jet-spanning space.  Their products give a triangular spanning set for
the jet filtration, in every characteristic.  For all sufficiently
large $p$,
\[
 5tk_0N\ge4t(p-1).
\]
The local monomial inclusion used here is the one in
\cite[proof of Proposition~2.12, equation~(2.8)]{MustataSchwede2014},
with smooth dimension $4t$. The products therefore span on the thick point defined by the $p$-th
powers of $4t$ smooth parameters.  Etaleness identifies the completed
local ring at the torsor lift with the completed local ring at its
tuple in $U^t$.  Under this identification the thick point is the
fibre of $F^t:W^t\to(W')^t$ over the Frobenius image of that tuple.

In a product of $N$ pure tensors, let $m_i'$ denote the new exponent
in slot $i$.  Then
\[
 0\le m_i'\le N(r-1)<pr,
 \qquad \sum_i m_i'\equiv0\pmod r.
\]
Since $p$ is prime to $r$, choose integers $a_i$ with
\[
 r\le a_i\le2r-1,\qquad pa_i\equiv m_i'\pmod r.
\]
Their sum is a multiple of $r$ and is at most $2tr-t<q$.  Increase
$a_1$ by the nonnegative multiple $q-\sum_i a_i$ of $r$.  We now have
\begin{equation}\label{eq:frob-padded-weights}
 a_i\ge0,\qquad \sum_i a_i=q,\qquad pa_i\ge m_i'.
\end{equation}
Multiply the corresponding slot section of $B+m_i'K_W$ by
$\eta_W^{(pa_i-m_i')/r}$.  This gives a section of $B+pa_iK_W$, and
division by $\tau^{pa_i}$ gives exactly its previous torsor expression.
In a local canonical frame, the coefficient of $\tau$ is a unit.  Its
$p$-th power is a nonzero constant on the thick Frobenius fibre, since
the $p$-th powers of all local parameters vanish there.

It follows that the projection-formula map
\begin{equation}\label{eq:frob-column-map}
 \bigoplus_{\substack{a_i\ge0\\\sum_i a_i=q}}
 \left(\bigotimes_{i=1}^tH^0(W,\OO_W(B+pa_iK_W))\right)
 \otimes
 \OO_{(W')^t}\left(-\sum_{i=1}^t a_i\operatorname{pr}_i^*K_{W'}\right)
 \longrightarrow\mathcal E^{\boxtimes t}
\end{equation}
has rank $R$ at one point, hence at the generic point.  The products
just constructed give the spanning values, up to the nonzero constant
frame factors.  By finite pushforward and base change these values are
the fibre of the target.  The identity used in the projection formula
is $F^*\OO_{W'}(K_{W'})\simeq\OO_W(pK_W)$; it is an identity of line
bundles, not the pullback map on differential forms.  In compatible
local frames its transitions are $p$-th powers.

\emph{The loss of rank on the diagonal.}
On the diagonal $W'\subset(W')^t$, all the source line bundles in
\eqref{eq:frob-column-map} equal $\OO_{W'}(-qK_{W'})$.  Let
\[
 \Delta_F=\underbrace{W\times_{W'}\cdots\times_{W'}W}_{t\ \mathrm{factors}},
 \qquad h:\Delta_F\to W',\qquad g:\Delta_F\to W
\]
be the common structure map and the first projection, respectively.
Both are finite.  The line bundle
\begin{equation}\label{eq:frob-diagonal-line}
 \mathcal L=\OO_{\Delta_F}\left(
       \sum_{i=1}^t\operatorname{pr}_i^*B
       +pq\operatorname{pr}_1^*K_W\right)
\end{equation}
satisfies
\[
 h_*\mathcal L\simeq
 (\mathcal E^{\boxtimes t})|_{W'}\otimes\OO_{W'}(qK_{W'}).
\]
This is finite base change from $F^t$, followed by the projection
formula; all the line bundles $\operatorname{pr}_i^*\OO_W(pK_W)$
are the pullback of $\OO_{W'}(K_{W'})$.  The columns of
\eqref{eq:frob-column-map}, after twisting, thus restrict to global
sections of $h_*\mathcal L$.  At every point of the diagonal their rank
is at most $h^0(\Delta_F,\mathcal L)$.

The two-factor diagonal filtration and its truncated symmetric pieces
are the canonical filtration of
\cite[Theorem~3.7]{Sun2008} and
\cite[Definition~3.1, Remark~3.2, and Corollary~3.5]{KitadaiSumihiro2008}.
We use the same local construction in $t-1$ difference slots, and give
the resulting ranks and determinant comparison explicitly.
Filter $\mathcal L$ by powers of the ideal of the reduced diagonal
$W\subset\Delta_F$ and push forward by $g$.  The associated pieces are
\begin{equation}\label{eq:frob-diagonal-graded}
 \begin{gathered}
 \mathcal G_j=\OO_W(tB+pqK_W)\otimes A_j(V),
 \qquad V=(\Omega_W^1)^{\oplus(t-1)},\\
 0\le j\le u_0=4(t-1)(p-1).
 \end{gathered}
\end{equation}
For completeness, in etale smooth coordinates the other slots add
$4(t-1)$ difference variables with their $p$-th powers zero and no
other relations.  This follows from the Cartesian relative Frobenius
square for an etale chart.  Modulo squares the differences transform
by ordinary differentials, giving $V$; the associated graded algebra
is therefore intrinsically $A_\bullet(V)$.  The line bundle contributes
its restriction $\OO_W(tB+pqK_W)$ to each piece.  In particular, with
$e_j=\rk A_j(V)$,
\begin{equation}\label{eq:frob-diagonal-section-sum}
 h^0(\Delta_F,\mathcal L)\le\sum_{j=0}^{u_0}h^0(W,\mathcal G_j),
 \qquad \sum_{j=0}^{u_0}e_j=p^{4(t-1)}=s^{t-1}.
\end{equation}

Set $n=4(t-1)$, $G=\max\{0,2g(C)-2\}$, and $c_2=n(n-1)G$.
By the retained inequality \eqref{eq:frob-flag-nonnegative},
$\mu_{\min}(V|_C)\ge0$.  Lemma~\ref{lem:frob-truncated-slopes} gives,
uniformly in $j$ and $p$,
\begin{align}
 \mu_{\max}(\mathcal G_j|_C)
 &\le
 \bigl(tB+pqK_W+(p-1)(t-1)K_W\bigr)\cdot C+c_2\notag\\
 &\le M_p:=2b_0p\,l_1l_2l_3H^4+c_2.
 \label{eq:frob-Mp}
\end{align}
Indeed $BH^3\le pbLH^3$, $K_WH^3\ge0$, and
\eqref{eq:frob-polarization} bounds the two contributions by
$pb_0H^4$ each, before multiplication by $l_1l_2l_3$.

Put $a=l_1l_2l_3H^4$ and $d_i=l_iH\cdot C=l_i a$.  Applying
Lemma~\ref{lem:frob-flag-sections} to each $\mathcal G_j$ yields
\begin{align}
 h^0(W,\mathcal G_j)
 &\le e_j(M_p+1)\prod_{i=1}^3(1+M_p/d_i)\notag\\
 &=e_jp^4\bigl((2b_0)^4H^4+O(p^{-1})\bigr).
 \label{eq:frob-uniform-sections}
\end{align}
The error is uniform in $j$, since $c_2,a,d_i$ were fixed before $p$
varied.  Explicitly, the leading coefficient is
$(2b_0)^4a^4/\prod_i d_i=(2b_0)^4H^4$.
Summing \eqref{eq:frob-uniform-sections}, using
\eqref{eq:frob-diagonal-section-sum}, $H^4\le3$, and the strict first
inequality in \eqref{eq:frob-choice-b0}, gives, for all sufficiently
large $p$,
\begin{equation}\label{eq:frob-diagonal-rank}
 h^0(\Delta_F,\mathcal L)\le R/4.
\end{equation}
Thus the rank of \eqref{eq:frob-column-map} is at most $R/4$ at every
point of the diagonal.

\emph{The determinant contradiction.}
Choose $R$ columns of \eqref{eq:frob-column-map} with nonzero
determinant, and write $a_i^{(\nu)}$ for the weight of slot $i$ in
column $\nu$.  Their determinant is a nonzero section of the external
product of the line bundles
\[
 \mathcal Q_i=(\det\mathcal E)^{\otimes R/s}
     \otimes\OO_{W'}\left(\sum_{\nu=1}^R a_i^{(\nu)}K_{W'}\right).
\]
Let $\lambda_i=R^{-1}\sum_{\nu=1}^R a_i^{(\nu)}$; then
$0\le\lambda_i\le q$.
Identifying numerical classes through base-field twist, we have
\begin{equation}\label{eq:frob-determinant-class}
 \frac{c_1(\mathcal Q_i)}{R}
 =\frac{B}{p}
   +\left(\frac{p-1}{2p}+\lambda_i\right)K_W.
\end{equation}
The first-Chern coefficient is the usual Frobenius direct-image
coefficient; compare \cite[Lemma~4.2]{Sun2008}. Here only its
numerical divisor class is needed. To check it directly, finite base change identifies
$F^*\mathcal E$ with the pushforward from $W\times_{W'}W$ of the line
bundle pulled back from its second slot.  The same truncated-diagonal
filtration used in \eqref{eq:frob-diagonal-graded}, now with two slots,
has pieces $\OO_W(B)\otimes A_j(\Omega_W^1)$.
Across all $p^4=s$ truncated monomials, the total exponent of each of
the four variables is $s(p-1)/2$.  Taking determinants gives
\[
 c_1(F^*\mathcal E)=sB+\frac{s(p-1)}2K_W.
\]
Relative Frobenius pullback multiplies divisor classes by $p$ after
the base-twist identification.  Since the determinant of
$\mathcal E^{\boxtimes t}$ contributes
$(\det\mathcal E)^{R/s}$ in each slot, this proves
\eqref{eq:frob-determinant-class}.

Use now the base twist of the blow-up of $x$ and of the test class
$\beta$.  By $B/p\le bL$ on this test,
\eqref{eq:frob-test-nonnegative}, and $\lambda_i\le q$, the normalized
class in \eqref{eq:frob-determinant-class} pairs with the test to at
most $D^+\cdot\beta$.  Let $s_i$ be any nonzero section of
$\mathcal Q_i$, and let $m_i=\ord_{x'}(s_i)$.  Its divisor has effective
strict transform of class
$\pi_{x'}^*c_1(\mathcal Q_i)-m_iJ'$ on the twisted blow-up.  The
nonnegativity of the test on effective divisors and
\eqref{eq:frob-movable-test} imply
\[
 m_i(J\cdot\beta)
 \le c_1(\mathcal Q_i)\cdot\beta
 \le R D^+\cdot\beta
 <R d_0(J\cdot\beta).
\]
Here the pairings have been identified with their untwisted numerical
values.  As $J\cdot\beta>0$, every such section has order strictly less
than $Rd_0$ at $x'$.

The determinant section is nonzero, so the K\"unneth formula ensures
that each slot section space is nonzero.  By
Lemma~\ref{lem:frob-external-orders}, its order at
$(x',\ldots,x')$ is therefore strictly less than
\begin{equation}\label{eq:frob-order-upper}
 Rt d_0=R b_0(C_1+1)<3R/4.
\end{equation}
On the other hand, \eqref{eq:frob-diagonal-rank} bounds the rank of its
column matrix at this point by $R/4$.  Trivialize the source and target
bundles locally, and perform constant invertible row operations so
that at least $R-\lfloor R/4\rfloor$ rows vanish at the point.  Every
entry of these rows lies in the maximal ideal.  Each term of the
determinant is consequently in its
$(R-\lfloor R/4\rfloor)$-th power.  The determinant has order at least
$3R/4$, contradicting \eqref{eq:frob-order-upper}.

This excludes the assumed Seshadri inequality and proves
\eqref{eq:frob-cost}.  All varieties, line bundles, flags, points, and
testing data were fixed before the characteristic grew.  Their
dependence on $X,r,t$ therefore affects only how large the
characteristic must be, whereas the constant $5/b_0$ in the conclusion
depends only on $C_1$.
\end{proof}

\section{Many factors and birational transports}
\label{sec:transports}

We next turn the quadratic cost bound into a birational obstruction.
The argument uses many factors, but their number is fixed before the order
of the covering group tends to infinity.  This order of choices will be
important in the application. Rational evaluation on diagonal products
also appears in algebraic formulations of Lie's superposition theorem;
compare \cite[\S\S3,6--7 and Theorem~7.1]{BlazquezSanzMorales2010}.
We prove the rational evaluation needed here directly from a generically finite
forgetting map of degree one.

\begin{proposition}\label{prop:many-factor-contradiction}
There is no sequence of finite cyclic covers
\[
 \pi_\nu:Y_\nu\longrightarrow X_\nu
\]
of normal projective fourfolds, generically torsors for cyclic groups
\(G_\nu\) of orders \(r_\nu\to\infty\), with the following properties.
The group \(\Bir(X_\nu)\) is countable, and there are ample rational
Cartier divisors \(L_\nu\) on \(X_\nu\) and constants \(c,C>0\),
independent of \(\nu\), such that at very general points
\[
 \varepsilon(L_\nu)\ge c,
 \qquad
 \varepsilon(\pi_\nu^*L_\nu)\ge c r_\nu^{1/4}.
\]
Moreover, for each fixed integer \(t\ge2\), put
\[
 Z_{\nu,t}=Y_\nu^t/G_{\nu,\mathrm{diag}},\qquad
 \theta_{\nu,t}:Z_{\nu,t}\longrightarrow X_\nu^t,
 \qquad
 P_{\nu,t}=\theta_{\nu,t}^*L_\nu^{\boxplus t}.
\]
Here \(L^{\boxplus t}=\sum_{i=1}^t\operatorname{pr}_i^*L\).
The further required bound is
\[
 \varepsilon(P_{\nu,t})\le Ct^2
\]
at a very general point, for all sufficiently large \(\nu\), with this
threshold allowed to depend on \(t\).
\end{proposition}

\begin{proof}
Choose an integer \(T\), to be made sufficiently large in terms of
\(c,C\).  Passing to a tail of the sequence, we may use all the upper
bounds with \(2\le t\le T\) simultaneously.  We suppress \(\nu\)
from the notation until it is necessary to vary the cover.  Write
\(L_Y=\pi^*L\).

Since \(G\) is abelian, the finite map \(\theta_t\) is the quotient
by \(G^t/G_{\mathrm{diag}}\).  Indeed the quotient by the full product
group has function field \(\C(X^t)\), by the generic torsor property,
and the resulting finite birational map to the normal variety \(X^t\)
is an isomorphism.  In particular, \(P_t\) is ample.

\medskip\noindent
\emph{Curve families and compatible leaves.}
Set
\[
 B=CT^2+1.
\]
On each \(Z_t\), choose a marked covering family of integral curves
whose degree divided by multiplicity at the mark is at most \(B\).
Such a family is obtained by choosing a curve at a geometric generic
point witnessing a strict bound below \(Ct^2+1\), and spreading the
curve and its mark.  We use the geometric generic conventions of
Lemma~\ref{lem:generic-chain}: finite data can be defined over a
countable field, and the relevant generic marked families can be
realized over \(\C\).

Enlarge these lists by all deck translates, all permutations of the
slots, and all nonconstant image families under ordered projections
\(Z_t\to Z_s\), for \(2\le s<t\le T\).  Close under these
operations.  The lists remain finite, although their sizes need not be
bounded uniformly in the covering order.

Projection does not increase the marked curve cost.  To see this, let
\(D\to D'\) be the nonconstant map of an integral curve to its reduced
image, let \(e\) be its degree, and let \(z\) and \(z'\) be the marks.
A general local hyperplane through \(z'\), pulled to the normalizations
of the curves, has total order
\(e\mult_{z'}D'\) over \(z'\).  The contribution of the branches
through \(z\) is at least \(\mult_zD\).  Hence
\[
 e\mult_{z'}D'\ge\mult_zD.
\]
Since \(P_t\) minus the pullback of \(P_s\) is nef, we obtain
\[
 \frac{P_s\cdot D'}{\mult_{z'}D'}
 \le
 \frac{P_t\cdot D}{\mult_zD}.
\]
After shrinking parameter spaces we may use reduced integral geometric
generic image curves.  Their projected mark maps remain dominant.

Let \(H_t\) be a geometric generic leaf furnished by
Lemma~\ref{lem:generic-chain}, and put \(h_t=\dim H_t\).
Lemma~\ref{lem:chain-order} gives
\begin{equation}\label{eq:transport-chain-degree}
  h_t>0,\qquad P_t^{h_t}\cdot H_t\le(h_tB)^{h_t}.
\end{equation}
Deck invariance and Corollary~\ref{cor:chain-functoriality} imply that
\(W_t=\theta_t(H_t)\) is the geometric generic fibre closure of a
rational quotient of \(X^t\).  We write
\[
 \xi_t:X^t\dashrightarrow S_t
\]
for such a quotient, with geometrically integral generic fibre.
The leaves and their quotients are equivariant under permutations.

The leaves are also compatible with coordinate projections: an upper
leaf projects into the lower leaf through the projected tuple.  Here is
the genericity justification.  Retain the parameter space of a curve
family when forming its projected image family.  Every parameter and
endpoint occurring in a generic upper chain is generic in its original
parameter space and incidence over the initial field of definition.
Its projected endpoint is generic in the image curve.  For this lower
family, the equality of leaf names at a mark and a generic endpoint is
an identity on the generic incidence.  Consequently it applies to the
projected upper step, regardless of the additional upper history;
a stationary projection gives the same equality immediately.  A path
reaching a generic point of the upper leaf therefore has constant lower
name.  This proves the asserted inclusion.  Whenever different slot
sizes are compared below, we take leaves through a generic tuple and
its projections and extend their name fields to a common algebraically
closed field.

\medskip\noindent
\emph{Finite projections to each slot.}
For \(r\) sufficiently large in terms of \(T,c,C\), the map from
\(H_t\) to each single \(X\)-slot is generically finite onto its
image.  Fix one slot, let its image be \(I\), and write
\[
 \ell=\dim I,\qquad d=h_t-\ell.
\]
The image contains an ambient generic point \(x\) of \(X\), which
is also generic in \(I\) over the leaf name.  The Seshadri lower bound
implies
\begin{equation}\label{eq:transport-image-degree}
 L^\ell\cdot I\ge c^\ell,
\end{equation}
where the assertion is immediate for \(\ell=0\).

Suppose \(d>0\).  Over this generic point \(x\), the geometric fibre
of \(Z_t\to X\) is \(Y^{t-1}\): choose a lift of \(x\) in its
torsor to identify the first coordinate.  The restriction of \(P_t\)
is the external sum of the remaining copies of \(L_Y\).  Every one
of the marked coordinates is generic on \(Y\).  Projecting any curve
through that mark to a nonconstant slot, with the preceding
degree--multiplicity comparison, shows that the Seshadri constant of
this restricted polarization is at least \(cr^{1/4}\).

Let \(J\) be a \(d\)-dimensional component of the geometric generic
fibre of \(H_t\to I\) containing the chosen generic point.  Then
\begin{equation}\label{eq:transport-fibre-degree}
 P_t^d\cdot J\ge (cr^{1/4})^d.
\end{equation}
We recall the elementary degree bound being used twice here.  If an
ample rational divisor has Seshadri constant at least \(a\) at a
smooth point of an \(e\)-dimensional subvariety, blow up that point
and restrict the nef class consisting of the pulled-back divisor minus
\(a\) times the exceptional divisor.  Its top intersection on the
strict transform gives degree at least \(a^e\), since the point has
multiplicity one.  A limiting argument permits a Seshadri bound given
as a supremum.  Our image and fibre points are smooth generic points
of their subvarieties.  The ambient very general bounds continue to
hold on the geometric generic data after the field extensions in use.

The pushforward of \(P_t^d\cdot H_t\) to \(I\) has coefficient
equal to the total degree on the geometric generic fibre, at least
the contribution of \(J\).  Since
\(P_t-\operatorname{pr}^*L\) is nef, equations
\eqref{eq:transport-image-degree} and
\eqref{eq:transport-fibre-degree} give
\[
 P_t^{h_t}\cdot H_t
 \ge (\operatorname{pr}^*L)^\ell P_t^d\cdot H_t
 \ge c^\ell(cr^{1/4})^d.
\]
For fixed \(T\), this contradicts
\eqref{eq:transport-chain-degree} as \(r\to\infty\), because
\(h_t\le4T\).  There are only finitely many slot sizes and slots to
consider.  This proves the assertion simultaneously for all of them
on a tail of the sequence.  In particular,
\begin{equation}\label{eq:transport-ranks}
 1\le h_t\le4.
\end{equation}

\medskip\noindent
\emph{A forgetting map of degree one.}
The same single-slot finiteness holds for \(W_t\).  Let \(D_t\)
denote the degree of its projection to its first-slot image.  The
degree of \(H_t\) over that image contains \(D_t\) as a factor, so
\eqref{eq:transport-chain-degree} and
\eqref{eq:transport-image-degree} imply
\begin{equation}\label{eq:transport-degree-bound}
 D_t\le \frac{P_t^{h_t}\cdot H_t}{c^{h_t}}
       \le K T^8
\end{equation}
where one may take
\[
 K=\max_{1\le h\le4}\left(\frac{h(C+1)}c\right)^h.
\]
Indeed \(h_t\le4\), \(B\le(C+1)T^2\), and \(T^{2h_t}\le T^8\).
Thus \(K\) depends only on \(c,C\), before \(T\) is selected.

Projection compatibility and finite first-slot projection give
\(h_t\le h_{t-1}\).  If these ranks are equal, the image of \(W_t\)
under forgetting its last slot equals the lower leaf: it is contained
in that integral leaf and has the same dimension.  The first-slot
images therefore agree as well.  The degree formula for the two
generically finite maps is
\begin{equation}\label{eq:transport-integer-ratio}
 \frac{D_t}{D_{t-1}}
 =\deg(W_t\longrightarrow W_{t-1})\in\N.
\end{equation}
All degrees are taken after the common geometric field extension
described above, which preserves them.

By \eqref{eq:transport-ranks}, among the \(T-1\) ranks there is a
constant stretch of length at least \((T-1)/4\).  Choose \(T\),
depending only on \(c,C\), so that
\[
 \frac{T-1}{4}-1>\log_2(KT^8).
\]
If every adjacent ratio on that stretch were
larger than one, \eqref{eq:transport-integer-ratio} would force the
degrees to grow by a factor at least two at each step, contradicting
\eqref{eq:transport-degree-bound}.  Hence for some \(k\ge3\),
\begin{equation}\label{eq:transport-birational-omission}
 h_k=h_{k-1}=h>0,
 \qquad W_k\longrightarrow W_{k-1}
 \quad\hbox{is birational}.
\end{equation}
Permutation symmetry gives the same conclusion for every omission of
one slot, with the corresponding lower leaf.  In particular, no
finite correspondence of degree greater than one remains in this
step.

\medskip\noindent
\emph{Rational transport of tangent directions.}
At a generic tuple of \(X^k\), the tangent space along the leaf
\(W_k\) is a rational rank-\(h\) distribution
\(\mathcal D_k=\ker(d\xi_k)\).  It projects injectively to each
single-slot tangent space and isomorphically to the lower distribution
for every omission.  Indeed the corresponding maps of geometric
generic leaves are generically finite, and in characteristic zero
their differentials have rank \(h\) at the generic point.  By the
equal-rank inclusion already proved, their images are exactly the
lower tangent spaces.

Write \(A_i\subset T_{X,x_i}\) for the image in slot \(i\).
This plane depends rationally only on \(x_i\).  To justify the
descent, trivialize \(T_X\) rationally and use a Grassmann chart.
The coordinates of \(A_i\) belong to the function field of each
\((k-1)\)-tuple containing slot \(i\), because they are obtained
from that lower distribution.  The intersection of these coordinate
subfields is the function field of slot \(i\).  Similarly, the graph
of the isomorphism from \(A_i\) to \(A_j\), obtained by lifting to
\(\mathcal D_k\) and projecting to slot \(j\), depends only on
\((x_i,x_j)\).

For completeness, the coordinate-field intersection can be checked
by replacing an omitted coordinate by a fresh independent generic
coordinate.  A function pulled from the other coordinates is unchanged
under that replacement.  Repeating for every omitted coordinate, and
then specializing the fresh coordinates to general constants, shows
that the function belongs to the field of the common coordinates.
This applies to the rational chart coordinates of the planes and
graphs.  It applies also when \(k=3\).

Permutation symmetry thus gives a common rational plane distribution
\(A(x)\) and isomorphisms
\[
 R(x,y):A(x)\xrightarrow{\sim}A(y).
\]
Lifting a vector through the same \(h\)-plane and using three slots
gives the rational cocycle identity
\begin{equation}\label{eq:transport-cocycle}
 R(y,z)R(x,y)=R(x,z).
\end{equation}
Choose a general complex reference point \(a\) so that
\(R(a,-)\) is generically defined and invertible and the specialized
cocycle identity holds.  Transporting a basis of \(A(a)\) gives
rational vector fields \(v_1,\ldots,v_h\) on \(X\), independent at
a generic point.  Equation~\eqref{eq:transport-cocycle} says exactly
that the simultaneous vector fields
\[
 (v_j,\ldots,v_j),\qquad 1\le j\le h,
\]
span \(\mathcal D_k\).  No commutativity assertion about the
\(v_j\) is needed.

\medskip\noindent
\emph{Rational evaluation and local flows.}
Equal ranks were enough to construct the tangent transports. We now
use the stronger degree-one assertion: it turns evaluation of the
missing coordinate into a rational function rather than a finite
correspondence.
Consider the rational map to its image closure
\[
 \Lambda=(\operatorname{pr}_{<k},\xi_k):X^k\dashrightarrow D.
\]
It is birational.  Its restriction to a geometric generic leaf of
\(\xi_k\) is the forgetting map in
\eqref{eq:transport-birational-omission}, and hence has degree one.
More explicitly, let
\[
 F_0\subset E\subset F
\]
be the function fields of \(S_k,D,X^k\), included by pullback.
The field \(E\) is generated over \(F_0\) by the first
\((k-1)\)-slot functions. Since the geometric generic forgetting map
is generically finite, \(F\) and \(E\) have the same transcendence
degree over \(F_0\); hence \(F/E\) is finite. Geometric integrality
of the generic fibre makes \(F/F_0\) regular.  Tensoring with an algebraic closure
\(\overline F_0\) gives fields: each tensor is the union of the
fields obtained by finite algebraic base extension.  Their fraction
fields are the function fields of the geometric generic leaf and its
projected image.  The degree-one geometric projection therefore gives
\([F:E]=1\).  This also rules out an extra generic component or an
unseen algebraic degree in \(\Lambda\).

Let \(\mathrm{ev}\) be the last coordinate of its rational inverse.
Choose a general complex tuple \((u,y_0)\), with
\(u\in X^{k-1}\), in smooth open sets on which the vector fields,
the quotient, and the inverse evaluation are regular.  For small
\(z=(z_1,\ldots,z_h)\in\C^h\), let \(E_z\) be the ordered
composition of the local time-\(z_j\) flows of \(v_j\).  Let
\(E_z^-\) be the reversed composition of the negative-time flows.
These flows exist holomorphically in both the initial point and the
times on sufficiently small neighbourhoods of all chosen coordinates.
One may obtain this directly by Picard iteration in smooth coordinate
polydiscs: bounded coefficients and derivatives make the integral
operator a uniform contraction for small time.  Uniqueness gives the
opposite-time inverse, and reversing the order gives the inverse of
the composition.

The simultaneous flows preserve \(\xi_k\).  Put
\(\lambda(z)=E_z(u)\), acting in all the reference slots.  For
\(y\) near \(y_0\) and sufficiently small \(z\), birational
evaluation therefore gives
\begin{equation}\label{eq:transport-flow-evaluation}
 E_z(y)=\mathrm{ev}\bigl(\lambda(z),\xi_k(u,y)\bigr).
\end{equation}
For each fixed such \(z\), the right side is rational in \(y\).
The equality on an analytic open shows that its input lies in \(D\)
and meets the domain of evaluation; an analytic open is Zariski dense.
Thus it defines a rational self-map of \(X\), regular at \(y_0\)
after shrinking the neighbourhoods and times.

The same argument gives
\[
 E_z^-(y)=\mathrm{ev}\bigl(E_z^-(u),\xi_k(u,y)\bigr)
\]
on a neighbourhood of \(y_0\).  The resulting rational maps are mutually inverse
because they are inverse on a smaller analytic neighbourhood.  They
are therefore birational self-maps of \(X\).  Finally,
\(z\mapsto E_z(y_0)\) has nonzero derivative, since the fields are
independent and \(h>0\).  Its image is uncountable, so the birational
maps, all regular at \(y_0\), are uncountably many distinct maps.
This contradicts the assumed countability of \(\Bir(X)\).
\end{proof}

\begin{lemma}\label{lem:birational-countability}
Let \(X\) be a normal projective terminal variety with
\(K_X\sim_\Q0\).  If a smooth projective resolution of \(X\) has
irregularity zero, then \(\Bir(X)\) is countable.
\end{lemma}

\begin{proof}
Let \(p:W\to X\) be a smooth projective resolution. Since
\(H^1(W,\OO_W)=0\), the exponential sequence embeds the analytic
Picard group in \(H^2(W,\Z)\). By GAGA the analytic and algebraic
Picard groups agree, so \(\Pic(W)\) is countable. Divisor pushforward
induces a surjection \(\Pic(W)\to\Cl(X)\): the strict transform of
every Weil divisor on \(X\) is Cartier on \(W\). Thus \(\Cl(X)\)
is countable.

Every birational self-map of \(X\) is an isomorphism in codimension
one.  To see this, take a smooth common resolution of its graph, with
projections \(p_1,p_2\) to \(X\).  Choose a nonzero trivialization
\(\eta\) of a sufficiently divisible pluricanonical bundle of
\(X\).  Its two pullbacks are nonzero regular pluricanonical forms
on the common resolution.  That section space has dimension one:
canonicality gives pullback regularity, and restriction to a big open
of \(X\), followed by divisorial extension, identifies the sections
with those of the trivial pluricanonical bundle on \(X\).
Consequently the two pullbacks are proportional.  Terminality says
that their zero divisors have support exactly the exceptional prime
divisors of \(p_1\) and \(p_2\), respectively.  These supports must
coincide.  A prime divisor contracted by the self-map or its inverse
would belong to only one of the two supports, which is impossible.

Fix a very ample divisor \(A\) on \(X\).  Group the birational
self-maps according to the Weil linear equivalence class of the strict
transform of \(A\) back to the source.  Within one class, the
complete spaces of sections are identified by the codimension-one
isomorphisms.  The two maps to the embedding of \(X\) by \(|A|\)
therefore differ by a projective linear automorphism preserving that
embedded variety.  It is enough to prove that this stabilizer has
finite effective image on \(X\).

Let $S$ be the identity component of the stabilizer in the projective
linear group of the embedding. Its action on $X$ induces a rational
action on the smooth projective resolution $W$. The effective image is
a connected linear algebraic group: the kernel of the action is a closed
normal subgroup, and a quotient of a linear algebraic group is linear.
As proved above, $W$ has a nonzero positive pluricanonical section.
Matsumura's obstruction therefore says that this effective image is
trivial \cite[Corollary~1]{Matsumura1963}. More precisely, that corollary
excludes positive-dimensional linear algebraic subgroups of $\Bir(W)$
for a smooth complete variety with a nonzero pluricanonical system.
It applies here to the rational action, without requiring that every
birational self-map of $W$ be regular.

The identity component of the stabilizer consequently acts trivially.
An algebraic group has only finitely many components, so the effective
stabilizer is finite.  There are countably many possible divisor
classes and finitely many birational maps in each such class.  Hence
\(\Bir(X)\) is countable.
\end{proof}

\section{Completion of the canonical-index proof}\label{sec:index-zero}

\begin{proof}[Proof of Theorem~\ref{thm:fourfold-index-zero}]
Apply the canonical cover and the decomposition
\eqref{eq:index-product-cover} to \(X\).  Proposition~\ref{prop:index-product-case}
settles all decompositions with at least two positive dimensional
blocks and the purely abelian case.  The remaining case is a single
four-dimensional nonflat factor.  Corollary~\ref{cor:index-noncanonical}
settles its noncanonical quotients.  It remains to bound the orders
when \(X\) is canonical and has such a single-factor cover.

Suppose that these orders are unbounded, and choose a sequence whose
orders \(r\) tend to infinity.  Replace each member by its crepant
\(\Q\)-factorial terminalization \(V\), using
Lemma~\ref{lem:index-terminalization}.  The order \(r\) is unchanged.
The conclusions of Lemma~\ref{lem:index-single-factor-images} are
birational properties, so every proper positive dimensional rational
image of \(V\) has rationally connected resolution, and a smooth
projective model of \(V\) has irregularity zero.

By Lemma~\ref{lem:index-moving-divisor}, any member carrying an
effective nonbig divisor of positive Iitaka dimension has canonical
order dividing the fixed integer \(N_{\rm mov}\).  After discarding
a finite initial part of the sequence, every effective Weil divisor
of positive Iitaka dimension is therefore big.  For the cyclic
index cover \(\pi:Y\to V\), Lemma~\ref{lem:index-general-type-fibres}
now verifies the precise intermediate-fibration hypothesis of
Theorem~\ref{thm:scalar-bounds}, both for \(V\) with the trivial group
and for \(Y\) with its cyclic deck group.

Choose an ample rational Cartier divisor \(L\) on each \(V\), scaling
it by a positive rational number so that \(1\le L^4\le2\).
The scaling is possible without any uniform Cartier denominator.
The scalar estimates give constants \(C,c>0\), independent of the
member of the sequence, such that at very general points
\[
 \gamma(L;V)\le C2^{1/4},\qquad
 \varepsilon(L;x)\ge c,
 \qquad
 \varepsilon(\pi^*L;y)\ge c r^{1/4}.
\]
For the last inequality we used
\((\pi^*L)^4=rL^4\); the polarization upstairs is ample and invariant
under the deck group. Choose once and for all a rational number
$C_1\ge C2^{1/4}$ and put $b_0=1/(4(C_1+1))$.
Then $b_0(C_1+1)=1/4<3/4$ and
$3(2b_0)^4<3/16<1/4$, so both strict conditions in
\eqref{eq:frob-choice-b0} hold. Proposition~\ref{prop:quadratic-cost}
therefore gives, for every fixed integer $t\ge2$,
\[
 \varepsilon(P_t)\le 20(C_1+1)t^2.
\]
The constants are independent of the member of the sequence.
For each fixed product, its auxiliary choices are made before the
characteristic tends to infinity in that proposition. We then use
only finitely many such product sizes before letting the canonical
orders $r$ tend to infinity in the transport argument.

Lemma~\ref{lem:birational-countability} applies to the terminal
\(V\), since \(K_V\sim_{\Q}0\) and its smooth projective model has
irregularity zero.  Thus \(\Bir(V)\) is countable.  All hypotheses of
Proposition~\ref{prop:many-factor-contradiction} have now been checked:
the cover is a finite cyclic torsor on a big smooth open, the
polarizations have the two displayed lower bounds, the diagonal
quotients have the uniform quadratic upper bounds, and the
downstairs birational group is countable.  That proposition rules
out \(r\to\infty\), a contradiction.

The orders are consequently bounded in this final case as well.
There are only the cases just treated.  Taking a common multiple of
their bounded orders gives a single positive integer \(N_4\) with
\(N_4K_X\sim0\) for every input \(X\), as required.
\end{proof}

\section{Proof of the uniform Iitaka theorem}
\label{sec:main-proof}

\begin{proof}[Proof of Theorem~\ref{thm:main}]
First we deduce nonvanishing from pseudo-effectivity. The companion
good-model theorem, with the \(\Q\)-factorial dlt choice in its
construction \cite[Theorem~11.1 and Proposition~2.5]{OpenAIAbundance},
gives a log minimal model \((Y,\Delta_Y)\) of \((X,\Delta)\).
Its boundary is the strict transform of \(\Delta\) together with any
extracted prime divisors of coefficient one. Hence it is rational,
and \(D_Y=K_Y+\Delta_Y\) is nef and \(\Q\)-Cartier. The companion's
rational abundance theorem \cite[Theorem~1.1]{OpenAIAbundance} makes
\(D_Y\) semiample. On a common smooth resolution
\[
 X\xleftarrow{p}W\xrightarrow{q}Y
\]
with compatible canonical divisors, its comparison lemma
\cite[Lemma~2.2]{OpenAIAbundance} gives
\[
 p^*D=q^*D_Y+E,\qquad E\geq0,\qquad E\text{ is }q\text{-exceptional}.
\]
Choose a positive integer \(r\) such that \(rD_Y\) is Cartier and
globally generated, and a nonzero section
\(s\in H^0(Y,\OO_Y(rD_Y))\), viewed in the common rational function field.
Pulling the effective Cartier divisor
\(\Div_Y(s)+rD_Y\) back by \(q\), adding \(rE\), and pushing forward
by \(p\) gives
\[
 \Div_X(s)+rD\geq0.
\]
The pole test \eqref{eq:floor-pole-test} therefore gives
\(0\ne s\in V_r(D)\). Choose a positive integer \(t\) such that
\(trD\) is integral Cartier. Then \(s^t\) is a nonzero section of
\(\OO_X(trD)\), so \(\kappa(X,D)\geq0\).
These auxiliary degrees need not be uniform: they only establish the
hypothesis for the uniform argument, whose all-degree comparison
will retain the full Iitaka field.

Enlarge the fixed coefficient set to $\Phi'=\Phi\cup\{1\}$.
Proposition~\ref{prop:semiample-model} gives a semiample dlt model
$(X',\Delta')$ with coefficients in $\Phi'$ and identical rounded
section spaces in every integer degree. Write its contraction as
$f:X'\to Z$ and its adjoint as
\[
 K_{X'}+\Delta'\sim_{\Q}f^*A,
\]
where $A$ is ample rational Cartier and
$\dim Z=\kappa(X,K_X+\Delta)$.

If $\dim Z>0$, apply Proposition~\ref{prop:effective-base} to
$f:(X',\Delta')\to Z$, with total-space dimension four and fixed
coefficient set $\Phi'$. Since $A$ is ample, it is big. The proposition
gives one degree depending only on $\Phi'$ whose every positive
multiple is nonempty and generates exactly $\C(Z)=K(D)$ inside the
common rational function field. This same degree covers all positive
base dimensions. When $\dim Z=4$, the contraction is birational and
its field is $\C(X')$.

Suppose instead that $Z$ is a point. The adjoint of $(X',\Delta')$
is then \(\Q\)-linearly trivial. There are three exhaustive cases.
If the pair is not klt, the four-dimensional non-klt index theorem
\cite[Corollary~1.7]{JiangLiu2021} gives a common linearly trivial
multiple depending only on $\Phi'$. If it is klt and
$\Delta'\ne0$, use \cite[Theorem~1.14]{Xu2019}. Its hyperstandard
coefficient convention includes any prescribed finite rational set:
one may take the defining set $\{1-b:b\in\Phi'\}$ and denominator
one. If it is klt and $\Delta'=0$, use
Theorem~\ref{thm:fourfold-index-zero}. A bound for the positive
torsion order gives a common trivial multiple by taking the least
common multiple of the finitely many possible orders.

Take a common multiple of these zero-dimensional degrees and of the
degree obtained for positive-dimensional bases. The choice depends
only on $\Phi$ and dimension four. In dimension zero it yields a
nonempty system with image a point. In positive dimension,
the every-multiple clause of Proposition~\ref{prop:effective-base}
retains the whole Iitaka field.
Finally \eqref{eq:all-rounded-sections} transfers the complete
systems back to the original pair. This proves every assertion,
including the rounded reflexive convention and recovery of the
full Iitaka base field. If a canonical representative is replaced by
$K_X+\Div(h)$, then the floor divisor changes by $m\Div(h)$ in
every integer degree. Multiplication by $h^{-m}$ identifies the two
section spaces and leaves all ratios unchanged. The same integer
therefore works for all canonical representatives.
\end{proof}

\begin{thebibliography}{99}

\bibitem{BierstoneMilman2008}
E.~Bierstone and P.~D.~Milman,
\emph{Functoriality in resolution of singularities},
Publ. Res. Inst. Math. Sci. \textbf{44} (2008), no.~2, 609--639.
\href{https://doi.org/10.2977/PRIMS/1210167338}{doi:10.2977/PRIMS/1210167338}.

\bibitem{BCHM2010}
C.~Birkar, P.~Cascini, C.~D.~Hacon, and J.~McKernan,
\emph{Existence of minimal models for varieties of log general type},
J. Amer. Math. Soc. \textbf{23} (2010), no.~2, 405--468.
\href{https://doi.org/10.1090/S0894-0347-09-00649-3}{doi:10.1090/S0894-0347-09-00649-3}.
\href{https://arxiv.org/abs/math/0610203v2}{arXiv:math/0610203v2}.

\bibitem{BirkarZhang2016}
C.~Birkar and D.-Q.~Zhang,
\emph{Effectivity of Iitaka fibrations and pluricanonical systems of polarized pairs},
Publ. Math. Inst. Hautes \'Etudes Sci. \textbf{123} (2016), 283--331.
\href{https://doi.org/10.1007/s10240-016-0080-x}{doi:10.1007/s10240-016-0080-x}.

\bibitem{Birkar2025}
C.~Birkar,
\emph{Singularities on Fano fibrations and beyond},
preprint, \href{https://arxiv.org/abs/2305.18770v2}{arXiv:2305.18770v2}
(2025), Corollary~1.8.

\bibitem{BlazquezSanzMorales2010}
D.~Bl\'azquez-Sanz and J.~J.~Morales-Ruiz,
\emph{Local and global aspects of Lie superposition theorem},
J. Lie Theory \textbf{20} (2010), no.~3, 483--517.
\href{https://jolt.centre-mersenne.org/item/10.5802/jolt.606.pdf}{Published article}.

\bibitem{BDPP2013}
S.~Boucksom, J.-P.~Demailly, M.~P\u{a}un, and Th.~Peternell,
\emph{The pseudo-effective cone of a compact K\"ahler manifold and
varieties of negative Kodaira dimension},
J. Algebraic Geom. \textbf{22} (2013), no.~2, 201--248.
\href{https://doi.org/10.1090/S1056-3911-2012-00574-8}{doi:10.1090/S1056-3911-2012-00574-8}.
Cited preprint version:
\href{https://arxiv.org/abs/math/0405285v1}{arXiv:math/0405285v1}.

\bibitem{CampanaPaun2015}
F.~Campana and M.~P\u{a}un,
\emph{Orbifold generic semi-positivity: an application to families of canonically polarized manifolds},
Ann. Inst. Fourier (Grenoble) \textbf{65} (2015), no.~2, 835--861.
\href{https://doi.org/10.5802/aif.2945}{doi:10.5802/aif.2945}.

\bibitem{Campana2004}
F.~Campana,
\emph{Orbifolds, special varieties and classification theory: an appendix},
Ann. Inst. Fourier (Grenoble) \textbf{54} (2004), no.~3, 631--665.
\href{https://doi.org/10.5802/aif.2028}{doi:10.5802/aif.2028}.

\bibitem{ChenHanLiu2023}
G.~Chen, J.~Han, and J.~Liu,
\emph{On effective log Iitaka fibrations and existence of complements},
Int. Math. Res. Not. IMRN \textbf{2024} (2024), no.~10, 8329--8349.
\href{https://doi.org/10.1093/imrn/rnad253}{doi:10.1093/imrn/rnad253}.
Also available as \href{https://arxiv.org/abs/2301.04813v1}{arXiv:2301.04813v1}.

\bibitem{ChenTsakanikas2023}
G.~Chen and N.~Tsakanikas,
\emph{On the termination of flips for log canonical generalized pairs},
Acta Math. Sin. (Engl. Ser.) \textbf{39} (2023), 967--994.
\href{https://doi.org/10.1007/s10114-023-0116-3}{doi:10.1007/s10114-023-0116-3}.
Preprint: \href{https://arxiv.org/abs/2011.02236v2}{arXiv:2011.02236v2}.

\bibitem{Druel2018}
S.~Druel, \emph{A decomposition theorem for singular spaces with trivial
canonical class of dimension at most five}, Invent. Math. \textbf{211}
(2018), no.~1, 245--296.
\href{https://doi.org/10.1007/s00222-017-0748-y}{doi:10.1007/s00222-017-0748-y}.
\href{https://druel.perso.math.cnrs.fr/textes/bobo.pdf}{Author's copy of the published text},
Theorem 1.2, with the
\href{https://druel.perso.math.cnrs.fr/textes/bobo_erratum.pdf}{author's erratum}.

\bibitem{EinKuchleLazarsfeld1995}
L.~Ein, O.~K\"uchle, and R.~Lazarsfeld,
\emph{Local positivity of ample line bundles},
J. Differential Geom. \textbf{42} (1995), no.~2, 193--219.
\href{https://www.math.stonybrook.edu/~roblaz/Reprints/Ein.Kuchle.Laz.LPALB.pdf}{Author's copy of the published article}.

\bibitem{FujinoGongyo2014}
O.~Fujino and Y.~Gongyo,
\emph{Log pluricanonical representations and the abundance conjecture},
Compos. Math. \textbf{150} (2014), no.~4, 593--620.
\href{https://doi.org/10.1112/S0010437X13007495}{doi:10.1112/S0010437X13007495}.
\href{https://www.math.kyoto-u.ac.jp/~fujino/fg-comp-final.pdf}{Final author manuscript}.

\bibitem{FujinoGongyoModuli}
O.~Fujino and Y.~Gongyo,
\emph{On the moduli b-divisors of lc-trivial fibrations},
Ann. Inst. Fourier (Grenoble) \textbf{64} (2014), no.~4, 1721--1735.
\href{https://doi.org/10.5802/aif.2894}{doi:10.5802/aif.2894}.

\bibitem{FujinoHashizume}
O.~Fujino and K.~Hashizume,
\emph{Existence of log canonical modifications and its applications},
Eur. J. Math. \textbf{9} (2023), no.~1, Paper No.~13.
\href{https://doi.org/10.1007/s40879-023-00598-0}{doi:10.1007/s40879-023-00598-0}.
Accepted author version of 28 April 2022, version 0.21:
\href{https://www.math.kyoto-u.ac.jp/~fujino/lc-model-fujino-hashizume-draft-7-EJM-rev-normalstylefile.pdf}{author manuscript}.

\bibitem{FujinoWeakPositivity}
O.~Fujino,
\emph{Notes on the weak positivity theorems},
in \emph{Algebraic Varieties and Automorphism Groups},
Adv. Stud. Pure Math., vol.~75,
Math. Soc. Japan, Tokyo, 2017, pp.~73--118.
\href{https://doi.org/10.2969/aspm/07510073}{doi:10.2969/aspm/07510073}.
Author version 0.54, dated 30 June 2015:
\href{https://www.math.kyoto-u.ac.jp/~fujino/weak-posi11.pdf}{author manuscript}.

\bibitem{GHS2003}
T.~Graber, J.~Harris, and J.~Starr,
\emph{Families of rationally connected varieties},
J. Amer. Math. Soc. \textbf{16} (2003), no.~1, 57--67.
\href{https://doi.org/10.1090/S0894-0347-02-00402-2}{doi:10.1090/S0894-0347-02-00402-2}.

\bibitem{GGK2019}
D.~Greb, H.~Guenancia and S.~Kebekus, \emph{Klt varieties with trivial
canonical class: Holonomy, differential forms, and fundamental groups},
Geom. Topol. \textbf{23} (2019), 2051--2124.
\href{https://doi.org/10.2140/gt.2019.23.2051}{doi:10.2140/gt.2019.23.2051}.
Corrected preprint version:
\href{https://arxiv.org/abs/1704.01408v5}{arXiv:1704.01408v5}, June 15, 2020.
Theorem B and Proposition D.

\bibitem{EGAIV3}
A.~Grothendieck, with the collaboration of J.~Dieudonn\'e,
\emph{\'El\'ements de g\'eom\'etrie alg\'ebrique. IV.
\'Etude locale des sch\'emas et des morphismes de sch\'emas.
Troisi\`eme partie},
Publ. Math. Inst. Hautes \'Etudes Sci. \textbf{28} (1966), 5--255.
\href{https://www.numdam.org/item/PMIHES_1966__28__5_0/}{Published article}.

\bibitem{HMX2014}
C.~D.~Hacon, J.~McKernan, and C.~Xu,
\emph{ACC for log canonical thresholds},
Ann. of Math. (2) \textbf{180} (2014), no.~2, 523--571.
\href{https://doi.org/10.4007/annals.2014.180.2.3}{doi:10.4007/annals.2014.180.2.3}.

\bibitem{HoringPeternell2019}
A.~H\"oring and T.~Peternell,
\emph{Algebraic integrability of foliations with numerically trivial canonical bundle},
Invent. Math. \textbf{216} (2019), no.~2, 395--419.
\href{https://doi.org/10.1007/s00222-018-00853-2}{doi:10.1007/s00222-018-00853-2}.
\href{https://math.univ-cotedazur.fr/~hoering/articles/a33-algebraic-integrability.pdf}{Author version}.

\bibitem{JiangLiu2021}
C.~Jiang and H.~Liu,
\emph{Boundedness of log pluricanonical representations of log Calabi--Yau
pairs in dimension 2},
Algebra Number Theory \textbf{15} (2021), no.~2, 545--567.
\href{https://doi.org/10.2140/ant.2021.15.547}{doi:10.2140/ant.2021.15.547}.
Preprint: \href{https://arxiv.org/abs/2002.11928v1}{arXiv:2002.11928v1}.

\bibitem{Jiang2021}
C.~Jiang,
\emph{A gap theorem for minimal log discrepancies of noncanonical
singularities in dimension three},
J. Algebraic Geom. \textbf{30} (2021), 759--800.
\href{https://doi.org/10.1090/jag/759}{doi:10.1090/jag/759}.
Cited final author version:
\href{https://arxiv.org/abs/1904.09642v3}{arXiv:1904.09642v3}.

\bibitem{Kawamata1992}
Y.~Kawamata,
\emph{Abundance theorem for minimal threefolds},
Invent. Math. \textbf{108} (1992), 229--246.
\href{https://doi.org/10.1007/BF02100604}{doi:10.1007/BF02100604}.

\bibitem{KitadaiSumihiro2008}
Y.~Kitadai and H.~Sumihiro,
\emph{Canonical filtrations and stability of direct images by Frobenius morphisms},
Tohoku Math. J. (2) \textbf{60} (2008), no.~2, 287--301.
\href{https://doi.org/10.2748/tmj/1215442876}{doi:10.2748/tmj/1215442876}.

\bibitem{FlipsAbundance1992}
J.~Koll\'ar (ed.),
\emph{Flips and abundance for algebraic threefolds:
A summer seminar at the University of Utah (Salt Lake City, 1991)},
Ast\'erisque \textbf{211},
Soci\'et\'e Math\'ematique de France, 1992.
\href{https://doi.org/10.24033/ast.198}{doi:10.24033/ast.198}.

\bibitem{Kollar2013}
J.~Koll\'ar, with the collaboration of S.~Kov\'acs,
\emph{Singularities of the minimal model program},
Cambridge Tracts in Mathematics, vol.~200,
Cambridge University Press, Cambridge, 2013.
\href{https://doi.org/10.1017/CBO9781139547895}{doi:10.1017/CBO9781139547895}.

\bibitem{KollarMori1998}
J.~Koll\'ar and S.~Mori,
\emph{Birational Geometry of Algebraic Varieties},
Cambridge Tracts in Mathematics \textbf{134},
Cambridge University Press, Cambridge, 1998.
\href{https://doi.org/10.1017/CBO9780511662560}{doi:10.1017/CBO9780511662560}.

\bibitem{Langer2004}
A.~Langer, \emph{Semistable sheaves in positive characteristic},
Ann. of Math. (2) \textbf{159} (2004), no.~1, 251--276.
\href{https://doi.org/10.4007/annals.2004.159.251}{doi:10.4007/annals.2004.159.251}.

\bibitem{Matsumura1963}
H.~Matsumura,
\emph{On algebraic groups of birational transformations},
Atti Accad. Naz. Lincei Rend. Cl. Sci. Fis. Mat. Natur. (8)
\textbf{34} (1963), no.~2, 151--155.
\href{https://www.bdim.eu/item?id=RLINA_1963_8_34_2_151_0}{Digital archive}.

\bibitem{MehtaRamanathan1982}
V.~B.~Mehta and A.~Ramanathan,
\emph{Semistable sheaves on projective varieties and their restriction to curves},
Math. Ann. \textbf{258} (1982), no.~3, 213--224.
\href{https://doi.org/10.1007/BF01450677}{doi:10.1007/BF01450677}.

\bibitem{MustataSchwede2014}
M.~Musta\c{t}\u{a} and K.~Schwede,
\emph{A Frobenius variant of Seshadri constants},
Math. Ann. \textbf{358} (2014), no.~3--4, 861--878.
\href{https://doi.org/10.1007/s00208-013-0976-4}{doi:10.1007/s00208-013-0976-4}.
Cited preprint version:
\href{https://arxiv.org/abs/1203.1081v1}{arXiv:1203.1081v1}.

\bibitem{OpenAILifting}
OpenAI,
\emph{Lifting sections from the reduced support of an adjoint},
OpenAI Math Release preprint
\href{https://github.com/openai/math/blob/main/preprints/Lifting-sections-from-the-reduced-support-of-an-adjoint-September-27-2026/paper.pdf}{OAI:Lifting-sections-from-the-reduced-support-of-an-adjoint-September-27-2026},
2026.

\bibitem{OpenAIAbundance}
OpenAI,
\emph{Log abundance in characteristic zero},
OpenAI Math Release preprint
\href{https://github.com/openai/math/blob/main/preprints/Log-abundance-in-characteristic-zero-September-24-2026/paper.pdf}{OAI:Log-abundance-in-characteristic-zero-September-24-2026},
2026.

\bibitem{OpenAIRelative}
OpenAI,
\emph{Relative denominators and effective systems for log Calabi-Yau fibrations},
OpenAI Math Release preprint
\href{https://github.com/openai/math/blob/main/preprints/Relative-denominators-and-effective-systems-for-log-Calabi-Yau-fibrations-September-27-2026/paper.pdf}{OAI:Relative-denominators-and-effective-systems-for-log-Calabi-Yau-fibrations-September-27-2026},
2026.

\bibitem{RamananRamanathan1984}
S.~Ramanan and A.~Ramanathan,
\emph{Some remarks on the instability flag},
Tohoku Math. J. (2) \textbf{36} (1984), no.~2, 269--291.
\href{https://doi.org/10.2748/tmj/1178228852}{doi:10.2748/tmj/1178228852}.

\bibitem{Sun2008}
X.~Sun,
\emph{Direct images of bundles under Frobenius morphism},
Invent. Math. \textbf{173} (2008), no.~2, 427--447.
\href{https://doi.org/10.1007/s00222-008-0125-y}{doi:10.1007/s00222-008-0125-y}.
Cited preprint version:
\href{https://arxiv.org/abs/math/0611360v2}{arXiv:math/0611360v2}.

\bibitem{StacksProject}
The Stacks Project Authors,
\emph{The Stacks Project},
Lemma~58.21.4, Tag
\href{https://stacks.math.columbia.edu/tag/0BMB}{0BMB},
accessed 25 September 2026.

\bibitem{TodorovXu2009}
G.~Todorov and C.~Xu, \emph{Effectiveness of the log Iitaka fibration
for 3-folds and 4-folds}, Algebra Number Theory \textbf{3} (2009),
697--710. \href{https://doi.org/10.2140/ant.2009.3.697}{doi:10.2140/ant.2009.3.697}.

\bibitem{ViehwegZhang2009}
E.~Viehweg and D.-Q.~Zhang, \emph{Effective Iitaka fibrations},
J. Algebraic Geom. \textbf{18} (2009), 711--730.
\href{https://doi.org/10.1090/S1056-3911-09-00515-3}{doi:10.1090/S1056-3911-09-00515-3}.
Cited preprint version:
\href{https://arxiv.org/abs/0707.4287v3}{arXiv:0707.4287v3}.

\bibitem{Xu2019}
Y.~Xu,
\emph{Some results about the index conjecture for log Calabi--Yau pairs},
preprint, \href{https://arxiv.org/abs/1905.00297v2}{arXiv:1905.00297v2}
(2019).

\end{thebibliography}

\end{document}
